\documentclass[11pt]{article}
\usepackage[T1]{fontenc}
\usepackage[utf8]{inputenc}
\usepackage{lmodern}
\usepackage{geometry}
\usepackage{hyperref}
\usepackage{booktabs}
\usepackage{enumitem}
\usepackage{array}
\geometry{margin=1in}

\title{Goertzel 4CP Route:\\Historical Audit and Current Lean Status}
\author{Oruzi / Codex}
\date{April 23, 2026}

\begin{document}
\maketitle
\sloppy

\begin{abstract}
This note records the current formal status of Ben Goertzel's four-color proof
route as represented locally by the older branch history, the present
\path{4cp_proof_v23.pdf} manuscript, and the corresponding Lean development in
\path{Mettapedia/GraphTheory/FourColor}.  The main correction made here is both
historical and logical: the currently formalized Lemma~4.3 obstruction matches
an older per-run purification claim, not the exact v23 Lemma~4.3 / Step~2
surface.  The next important update is positive: the exact v23 Step~2
single-component purification statement is now formalized and proved in Lean.
Separately, the current Lean development still exposes genuine manuscript-level
bugs at the spanning-forest leaf reduction and the naive well-founded peeling
step, but it also formalizes a stronger repaired replacement based on
boundary-root adjacency-distance data, canonical shortest-path parents, and
annulus-construction lowerings.  The remaining live blocker is the
boundary-order / annulus-geometry interface needed to instantiate that stronger
route.  At the weak collar-geometry layer, Lean now also shows that canonical
lowering to the annulus construction does not by itself provide the strict
parent-witness orientation needed by that repaired route.  A further
closed-walk obstruction now shows that the current all-roots-on-the-outer-
boundary adjacency-distance package is incompatible with the honest
two-component boundary split extracted from the closed-walk source; the first
ground-up repair is now also formalized as a two-sided annulus-root package
that allows roots to meet either boundary component while still lowering to
the existing interior-dual Theorem~4.9 pipeline.  A newer correction
also isolates the interior-vertex side condition in
Theorem~4.9: boundary erasure preserves Kirchhoff equations only after the
target vertex set is purified away from the selected boundary support.  So the
route remains blocked later on, but the evidence has to be stated precisely.
The repaired previous-boundary lane also has a new calibration: honest
closed-walk source data plus an arbitrary same-embedding weak collar geometry
does not force the missing current-boundary witness placement.
The diamond-with-triangle benchmark now sharpens the root-distance side of this
audit: it has an honest source and repaired previous-boundary / witness-cover
geometry, but its forcing interior edge blocks the current interior-dual
root-distance packages on the same embedding.
The current crux evidence is strong enough to reject several weak interfaces,
but not strong enough to certify a \(>90\%\) block of every Goertzel-route
continuation; the next responsible move is ground-up formalization of the
stronger annulus and boundary-order hypotheses.
\end{abstract}

\section{Current Lean status}

As of this note, the FourColor directory builds cleanly and contains no direct
\path{sorry} or \path{axiom} placeholders.  This matters because the negative
results discussed below are not informal suspicions hidden behind unfinished
proofs; they are exposed as checked lemmas, examples, and regression files.

The formalization currently contains:
\begin{itemize}[leftmargin=2em]
\item the algebraic generator layer for the Goertzel route;
\item an exact Lean proof of the v23 Step~2 single-component purification
identity;
\item theorem-surface formalization work around the v23 Theorem~4.9 annulus and
dual-peeling route;
\item explicit regression files showing where claimed reductions are too weak or
false;
\item a boundary-order note documenting the present embedding-side gap.
\end{itemize}

\section{Historical lineage}

It helps to separate two very different histories.

First, there is already a canonical accepted formal proof of the Four Color
Theorem: Gonthier's Coq development, now available in the
\path{rocq-community/fourcolor} repository.  That line is not based on the
Goertzel/Kauffman boundary-algebra route.

Second, there is the specific experimental lineage of the Goertzel route
itself.  In the local workspace history, that lineage runs through earlier
machine-checked and semi-formal efforts in Megalodon, Lean, and Isabelle.
Those efforts were not merely parallel reimplementations; they were part of the
process by which the route was stress-tested, revised, and restated.

The key historical point for the present note is this:
\begin{itemize}[leftmargin=2em]
\item an earlier boundary-algebra effort already ran into a contradiction around
an older per-run Lemma~4.3 purification identity;
\item the older Lean branch explicitly records this as a fatal flaw and points
to a Megalodon-checked refutation of that claim;
\item an archival Isabelle/HOL branch then explored a revised
disk/Kempe-closure-spanning presentation of the route, emphasizing strong-dual
and run-purification structure;
\item Goertzel v23 should be understood against that background: it is already a
repair attempt after an earlier machine-exposed failure mode;
\item in particular, v23 changes the local theorem surface: its Lemma~4.3 is a
run-invariance statement, and its later purification step is phrased using the
full boundary intersection \(D \cap \partial f\), not a single distinguished
run \(R\);
\item the current mettapedia Lean work now finds that v23 itself still has
formal obstructions.
\end{itemize}

So the present bug report is not an isolated first stumble.  It is the next
stage of an ongoing prover-guided audit of the route.

In the wider prover landscape, systems such as Isabelle, Lean, Coq/Rocq, and
Mizar all matter as comparison points for formal mathematical practice.  But
for the specific Goertzel-route lineage discussed here, the concrete artifacts
currently in hand are the earlier Megalodon-checked contradiction, an earlier
Lean branch, an archival Isabelle branch, and the present mettapedia Lean
development.

\section{Verified obstructions and audit status}

\subsection{Older per-run purification obstruction}

The file
\path{Mettapedia/GraphTheory/FourColor/GoertzelLemma43Obstruction.lean}
formalizes the older per-run purification surface using one distinguished run
\(R\) together with two complementary arcs \(A\) and \(A'\).

The key checked theorem is:
\begin{quote}
\path{goertzel_lemma43_boundary_claim_false}
\end{quote}

Its content is that the two indicator-chain contributions cancel on the run
itself, but their sum survives on the complementary arcs.  So the claim
``the sum equals \(\gamma \cdot 1_R\)'' fails in that abstraction once
complementary-arc mass is genuinely present.

This is a real formal obstruction, but it should be attributed carefully.  It
matches the older per-run purification pattern already recorded in the archival
branch history.  It does \emph{not} by itself refute the exact v23 Step~2 claim
whose target support is \(D \cap \partial f\).

\subsection{Exact v23 Lemma 4.3 / Step 2 status}

The exact v23 local story is subtler:
\begin{itemize}[leftmargin=2em]
\item v23 Lemma~4.3 is a face-boundary run-invariance statement;
\item the later purification step defines
\(Y_{f,D} := X_f^{ab}(C) \oplus X_f^{ab}(C^D)\) and claims
\(Y_{f,D} = c \cdot 1_{D \cap \partial f}\);
\item the current Lean tree now contains an exact formalization of this local
step in \path{GoertzelLemma44.lean}.
\end{itemize}

More precisely, the new definition
\path{v23PolarizedFaceGenerator} formalizes the exact Step~1 polarized
generator that marks only the \(a\)-colored half of the boundary, and the
theorem
\begin{quote}
\path{v23_single_component_purification}
\end{quote}
proves the claimed identity
\[
X_f^{ab}(C) \oplus X_f^{ab}(C^D) = c \cdot 1_{D \cap \partial f}
\]
for one `(a,b)` Kempe component \(D\).  A second theorem,
\path{sum_v23_single_component_purifications_eq_boundaryOnlyGenerator}, shows
that summing these exact v23 Step~2 identities over all boundary-hitting
components recovers the existing simplified boundary-only generator already used
downstream.

So the present formal evidence is now enough to say something sharper: v23 does
\emph{not} simply repeat the older Lemma~4.3 / purification error.  The old
per-run claim is formally refuted, while the revised v23 local purification
step is formally vindicated in the present abstraction.

\subsection{Step 2 spanning-forest leaf bug and its replacement}

The file
\path{Mettapedia/GraphTheory/FourColor/Theorem49ForestRegression.lean}
contains a concrete regression showing that a leaf of a chosen spanning forest
need not be a leaf in the ambient dual graph.

The example uses the path graph on four vertices as a spanning-tree candidate
inside the four-cycle.  Vertex \(0\) is a leaf of the path, but the ambient
cycle still has an extra adjacency from \(0\) to \(3\).  Therefore a reduction
step of the form
\begin{quote}
``leaf in the chosen spanning forest'' \(\Rightarrow\) ``all other incident dual
edges are already accounted for''
\end{quote}
is unsound without additional argument.

However, this is no longer the live Lean endpoint.  The current formal route
replaces the false leaf implication with the packaged theorem surface
\path{InteriorDualBoundaryRootAdjDistancePeelData} in
\path{Theorem49InteriorDualForest.lean}.  Instead of asserting that a forest
leaf controls all remaining incident ambient edges, that interface only asks
for a weaker and honest local condition: every non-witness edge on a peeled
face lies either on the selected boundary support or on an adjacent peeled
child face at strictly larger root distance in the chosen spanning forest.

From there the code canonically recovers the stronger parent-based packages:
boundary-root parent data, explicit root-distance parent data, height-ordered
peeling, and well-founded parent peeling.  On the annulus side, these repaired
packages are already lowered all the way to the BFS annulus construction layer
in \path{PlanarBoundaryAnnulusConstruction.lean}, including the outer-boundary-
root specialization.  So the regression is best read as a refutation of a
manuscript heuristic, not as a proof that no spanning-forest-based repair
exists.

\subsection{Well-founded peeling bug and its replacement}

The same regression file also proves that
\begin{quote}
acyclic forest \(+\) parent-edge adjacency
\end{quote}
does not by itself imply a well-founded parent peeling relation.

Even on the two-vertex tree, one may define a parent map sending each endpoint
to the other.  Every parent edge lies in the forest, the forest is acyclic, and
yet the induced parent relation has a 2-cycle and is not well founded.

So any correct peeling argument needs an extra descending invariant such as
root-distance, explicit height, or another strictly decreasing measure.

Again, the current Lean route has already internalized this repair.  The
packages just mentioned lower from adjacency-distance data to explicit
root-distance parent data, then to height-ordered parent data, and finally to
well-founded parent data.  So the surviving burden is not to rediscover that an
invariant is needed; it is to extract the stronger invariant-bearing package
from the annulus geometry used later in the manuscript.

\subsection{Annulus witness orientation gap}

There is now a second checked crux at the annulus interface, separate from the
endpoint-carrier issues discussed later.  The file
\path{Theorem49PlanarBoundaryAnnulusWitnessRegression.lean} contains an
explicit \path{counterCollarGeometry} satisfying the current
\path{PlanarBoundaryAnnulusCollarGeometry} fields.  This object canonically
lowers through \path{planarBoundaryAnnulusConstruction_of_annulusCollarGeometry}
to \path{counterAnnulusConstruction}.

The theorem
\path{not_parentWitnessOrientation_counterAnnulusConstruction} proves that the
lowered construction fails
\path{PlanarBoundaryAnnulusConstruction.ParentWitnessOrientation}.  The bad
middle-collar face has witness edge \path{eqEdge}; the only lower-distance face
does not contain that edge, while every face that does contain it has distance
at least as large as the bad face.  The packaged existential theorem
\path{exists_planarBoundaryAnnulusConstruction_without_parentWitnessOrientation}
records this as a construction-layer counterexample.

Thus the current weak collar-geometry axioms cannot be used as a black box to
obtain the parent-oriented, well-founded annulus route.  A continuation must
prove the missing orientation or positive-distance condition from stronger
geometry, or replace this part of the strategy.

That repair is now checked more sharply too.  The new file
\path{Theorem49PlanarBoundaryAnnulusGeometryRegression.lean} builds a tiny
two-face model \path{counterPreviousBoundaryWitnessGeometry} satisfying the
repaired surface
\path{PlanarBoundaryAnnulusPreviousBoundaryWitnessGeometry}, equivalently the
current-boundary form \path{WitnessOnCurrentBoundary}.  But its canonical
lowering
\path{planarBoundaryAnnulusConstruction_of_annulusCollarGeometry} still fails
\path{ParentWitnessOrientation}: the outer collar face itself becomes a peeled
construction face of distance \(0\), so it cannot have any strictly earlier
parent in the construction metric.  The same repaired geometry now also
explicitly yields an attached boundary-rooted face-peel certificate on that
same embedding, so this is not just a failure of one packaging choice: the
canonical construction-parent shell can fail even after the honest
certificate-side endpoint is already present.  So even the repaired annulus
geometry should not be funneled back through the canonical
construction-parent shell.  The honest ground-up annulus continuation is the
direct route from weak collar geometry to attached certificates, and from
repaired geometry to well-founded peeling, already formalized in
\path{Theorem49PlanarBoundaryAnnulusGeometry.lean}.

That conclusion is now calibrated on an honest source model too.  The new file
\path{Theorem49PlanarBoundaryAnnulusHonestGeometryRegression.lean} puts a weak
one-collar annulus geometry \path{diamondWithTriangleOneCollarGeometry} on the
genuine closed-walk source model \path{diamondWithTriangleEmbedding}.  On that
same embedding the one-collar geometry already yields an attached
boundary-rooted face-peel certificate via
\path{diamondWithTriangleOneCollarGeometry_hasAttachedRootedFacePeelCertificate},
but the canonical lowering
\path{diamondWithTriangleOneCollarConstruction} still fails
\path{ParentWitnessOrientation} by
\path{not_diamondWithTriangleOneCollarConstructionParentWitnessOrientation}.
More sharply, the source-side summary theorem
\path{closedWalkAnnulusBoundarySource_and_annulusCollarGeometry_and_attachedRootedFacePeelCertificate_without_interiorDualBoundaryRootAdjDistancePeelData_diamondWithTriangle}
shows that the honest facial closed-walk source, weak annulus geometry, and
attached certificate can coexist on one embedding even while the current
interior-dual boundary-root adjacency-distance package is absent there.  So
the ground-up annulus certificate lane is not only earlier than the canonical
construction-parent shell; it is also formally earlier than the present
interior-dual route on a genuine cyclic source model.  The same file now also
proves
\path{closedWalkAnnulusBoundarySource_and_annulusCollarGeometry_and_attachedRootedFacePeelCertificate_without_planarBoundaryAnnulusRootParentPeelData_diamondWithTriangle},
so even that honest source plus weak one-collar annulus geometry and attached
certificate still does not force the current parent-oriented annulus-root
target on the same cyclic source model.  Thus the direct certificate lane is
already strictly earlier than the live parent-peel burden too.  The same file
now also proves
\path{diamondWithTriangle_explicitTait_currentSufficientSameEmbeddingGeometry_consolidatedRouteDiagnosis},
which packages the whole same-embedding calibration on that genuine cyclic
source model: the full current positive geometry stack and explicit-Tait
Theorem~4.9 synthesis coexist with a nonempty raw interior-edge endpoint
carrier, an empty purified carrier, failure of the projected nonempty
endpoint, failure of generic
\path{InteriorDualBoundaryRootAdjDistancePeelData}, and failure of the live
parent-peel package.  The more specialized theorems
\path{diamondWithTriangle_explicitTait_currentSufficientSameEmbeddingGeometry_without_interiorDualBoundaryRootAdjDistancePeelData}
and
\path{diamondWithTriangle_explicitTait_currentSufficientSameEmbeddingGeometry_without_planarBoundaryAnnulusRootParentPeelData}
remain available as focused projections of that same diagnosis.  The same file
now also proves
\path{diamondWithTriangle_explicitTait_synthesis_without_planarBoundaryAnnulusRootParentPeelData},
so even the full Theorem~4.9 synthesis endpoint under the explicit Tait coloring
coexists with failure of the current parent-peel package on that same
embedding.  Thus this live parent-peel burden is not necessary even for a
genuine cyclic positive synthesis model.  This is now lifted to a reusable
converse-failure theorem too:
\path{Theorem49ForcingInteriorEdgeObstruction.lean} proves
\path{not_forall_nonempty_planarBoundaryAnnulusRootAdjDistancePeelData_of_exists_embedding_taitEdgeColoring_and_theorem49BoundaryRootSynthesis_and_forcingInteriorEdgeWitness},
\path{not_forall_nonempty_planarBoundaryAnnulusRootParentPeelData_of_exists_embedding_taitEdgeColoring_and_theorem49BoundaryRootSynthesis_and_forcingInteriorEdgeWitness},
and \path{Theorem49ForcingInteriorEdgeObstructionRegression.lean}
instantiates them on the diamond model via
\path{exists_embedding_taitEdgeColoring_and_theorem49BoundaryRootSynthesis_and_forcingInteriorEdgeWitness_diamondWithTriangleGraph}
and
\path{not_forall_nonempty_planarBoundaryAnnulusRootAdjDistancePeelData_of_taitEdgeColoring_and_theorem49BoundaryRootSynthesis_diamondWithTriangle_via_forcingInteriorEdgeWitness}
and
\path{not_forall_nonempty_planarBoundaryAnnulusRootParentPeelData_of_taitEdgeColoring_and_theorem49BoundaryRootSynthesis_diamondWithTriangle_via_forcingInteriorEdgeWitness}.
So any synthesis model carrying a forcing interior edge now formally refutes
the universal converse from full Theorem~4.9 synthesis back to either
same-embedding annulus-root adjacency-distance data or same-embedding
parent-peel data.  The same file now also proves
\path{diamondWithTriangle_explicitTait_currentSufficientSameEmbeddingGeometry_without_nonemptyProjectedSynthesis},
which seals that comparison on the same genuine cyclic source model: the full
current same-embedding positive geometry stack, together with explicit-Tait
Theorem~4.9 synthesis, already exists on
\path{diamondWithTriangleEmbedding}, yet the newer projected nonempty
endpoint still fails there because the selected-boundary-purified carrier is
empty.  The new wheel theorems
\path{exists_theorem49BoundaryRootSynthesis_wheelWithInnerTriangleGraph_via_degenerateAttachedCertificate}
and
\path{wheelWithInnerTriangleGraph_explicitTait_graphLevelSynthesis_with_fixedEmbedding_IDBRAD_obstruction}
show the complementary graph-vs-embedding separation: the wheel graph still
admits the full Theorem~4.9 synthesis endpoint for the explicit Tait coloring
on some embedding via the degenerate all-boundary certificate route, while the
honest source-preserving wheel embedding with the same coloring and an
explicit unblocked interior endpoint still fails generic same-embedding
\path{InteriorDualBoundaryRootAdjDistancePeelData}.  So the current wheel
obstructions are source-preserving route failures, not graph-level
impossibility of Theorem~4.9 synthesis for that colored graph.  The new theorem
\path{wheelWithInnerTriangleGraph_explicitTait_degenerateSynthesis_blocks_nonemptyProjectedSynthesis}
sharpens the positive side of that comparison: the current graph-level
witness is exactly the degenerate all-boundary embedding, so it reaches full
synthesis only on an embedding whose selected-boundary-purified interior-edge
endpoint carrier is empty and therefore cannot witness
\path{Theorem49BoundaryRootNonemptyProjectedSynthesis}.  Thus the currently
known wheel graph-level positive theorem still sits entirely off the
projected nonempty lane.  The new theorem
\path{wheelWithInnerTriangleGraph_explicitTait_graphLevelSynthesis_requires_embedding_change}
sharpens that diagnosis one step further: the currently known graph-level
witness is provably not the honest source-preserving wheel embedding at all,
because the degenerate witness has empty selected-boundary-purified carrier
while the honest wheel embedding has a surviving purified carrier.  So the
present wheel positivity result necessarily changes embeddings before it can
fire.  The stronger
theorem
\path{wheelWithInnerTriangleGraph_explicitTait_graphLevelSynthesis_with_fixedEmbedding_without_selector_or_acyclicEndpoint_or_IDBRAD}
shows that this same separation already covers every currently sufficient
same-embedding selector/acyclic endpoint on the honest wheel embedding too:
with the same explicit Tait coloring, honest source, nonempty purified
carrier, and explicit unblocked interior endpoint, the fixed wheel embedding
still has no boundary-free selector, no
well-founded/height-ordered/collar-layer witness surface, no weak annulus-
collar geometry, and no generic same-embedding
\path{InteriorDualBoundaryRootAdjDistancePeelData}.  The new theorem
\path{wheelWithInnerTriangleGraph_explicitTait_graphLevelSynthesis_with_fixedEmbedding_without_sameEmbedding_projectedSourceRoutes}
pushes that same wheel diagnosis all the way to the currently formalized
projected endpoint routes: the graph still admits full Theorem~4.9 synthesis
on some embedding, but on the honest source-preserving wheel embedding the
same explicit Tait coloring still cannot instantiate either the IDBRAD-based
projected endpoint route or the source-fixed canonical-parent projected
endpoint route.  So the benchmark now separates graph-level Theorem~4.9
positivity not only from the live same-embedding sufficient shells, but from
both concrete same-embedding projected source routes already formalized in
Lean.  The new theorem
\path{wheelWithInnerTriangleGraph_explicitTait_graphLevelSynthesis_with_fixedEmbedding_without_replacementPositiveProjectedGeometry_or_previousBoundaryWitness}
extends that same graph-vs-fixed-embedding separation to the newer repaired
positive route as well: the graph still admits full Theorem~4.9 synthesis on
some embedding, but on the honest source-preserving wheel embedding the same
explicit Tait coloring, honest source, and surviving purified carrier still
fail both direct replacement-positive packages and the repaired
previous-boundary witness geometry.  So current graph-level positivity on
this benchmark remains strictly upstream of every same-embedding positive
route surface now formalized around the live annulus manuscript lane.
Lean.

That route has now been pushed one step further on the honest source
interface itself.  In
\path{PlanarBoundaryClosedWalkSourceRoute.lean} the theorems
\path{zero_on_allEdges_of_exists_closedWalkAnnulusBoundarySource_and_annulusCollarGeometry_direct_of_annihilatesKempeClosureGeneratorFamily}
and
\path{zero_on_allEdges_of_exists_boundaryReachabilityData_and_dartSuccessorCycleEmbeddingData_and_selectedBoundaryArc_and_annulusCollarGeometry_direct_of_annihilatesKempeClosureGeneratorFamily}
show that once the same embedding carries weak annulus collar geometry,
boundary vanishing on the selected support, and the Kempe-generator
annihilation hypothesis, Step~2 edge-vanishing follows directly.  So the
remaining burden on the honest annulus lane is no longer to recover the old
interior-dual package: it is to derive same-embedding weak collar geometry,
or an equally strong direct substitute, from genuine cyclic boundary-order
hypotheses.  This source-facing collar lane now also reaches the full
Theorem~4.9 endpoint: \path{PlanarBoundaryClosedWalkSourceRoute.lean} proves
\path{exists_theorem49BoundaryRootSynthesis_of_exists_closedWalkAnnulusBoundarySource_and_annulusCollarGeometry_direct}
and
\path{exists_theorem49BoundaryRootSynthesis_of_exists_boundaryReachabilityData_and_dartSuccessorCycleEmbeddingData_and_selectedBoundaryArc_and_annulusCollarGeometry_direct}.
So once honest closed-walk or successor-cycle source data is paired with
same-embedding weak collar geometry, no additional theorem-4.9 algebra or
certificate packaging remains.

The route picture is now even sharper at the pure geometry level.  In
\path{Theorem49PlanarBoundaryAnnulusDecomposition.lean} the definition
\path{planarBoundaryAnnulusDecomposition_of_boundaryData} shows that under the
current weak boundary API, an annulus boundary split already yields a trivial
one-collar annulus decomposition with all ambient faces in the unique collar.
Then \path{PlanarBoundaryClosedWalkSource.lean} proves
\path{admitsPlanarBoundaryAnnulusDecomposition_of_admitsPlanarBoundaryClosedWalkAnnulusBoundarySource}.
So the honest closed-walk source already reaches the pure annulus
decomposition surface with no interior-dual or witness-edge hypotheses.  The
live burden is therefore not pure decomposition anymore; it is witness-edge
ownership and the stronger same-embedding collar geometry needed to continue
the annihilator branch.

That burden is now formalized more sharply as well.  In
\path{Theorem49PlanarBoundaryAnnulusWitness.lean} now proves the generic
criterion
\path{not_nonempty_planarBoundaryAnnulusWitnessAssignment_of_exists_two_distinct_interior_edges_on_faceBoundary_of_canonicalDecomposition}:
for the canonical one-collar decomposition built from annulus boundary data,
any ambient face carrying two distinct interior edges already blocks witness
assignment.  The concrete three-face model in
\path{Theorem49PlanarBoundaryAnnulusWitnessRegression.lean} is now just an
instance of that general obstruction.  Thus the new
source-to-decomposition lowering does not secretly solve Step~2: even after
decomposition becomes trivial, witness-edge ownership remains a genuine
additional geometric payload, and the obstruction is not tied to one
hand-crafted counterexample.

That obstruction now reaches the honest source lane itself.  In
\path{Theorem49PlanarBoundaryAnnulusWitness.lean} the new theorem
\path{not_nonempty_planarBoundaryAnnulusWitnessAssignment_of_exists_two_distinct_interior_edges_on_faceBoundary_of_closedWalkAnnulusBoundarySource}
pushes the same argument through the canonical one-collar decomposition
extracted from an honest
\path{PlanarBoundaryClosedWalkAnnulusBoundarySource}.  Then the new regression
file \path{Theorem49PlanarBoundaryAnnulusHonestWitnessRegression.lean} builds a
finite source model \path{sharedInteriorPairAnnulusEmbedding} with two outer
faces sharing the consecutive interior edges \path{sip01} and \path{sip12},
plus a disjoint inner triangular boundary component.  This model carries the
full honest source package, but
\path{not_nonempty_planarBoundaryAnnulusWitnessAssignment_sharedInteriorPairCanonicalDecomposition}
shows that the canonical decomposition extracted from that very source still
admits no witness assignment.  So the tempting repair ``derive boundary data
from an honest source, then use the canonical one-collar decomposition'' is
now blocked on an honest same-embedding source model, not just on an abstract
annulus-boundary witness.

The same obstruction now reaches beyond the canonical decomposition and into
the weak geometry lane itself, as long as that lane still tries to stay
one-collar.  \path{Theorem49PlanarBoundaryAnnulusWitness.lean} now proves
\path{not_nonempty_planarBoundaryAnnulusWitnessAssignment_of_numCollars_eq_one_and_exists_two_distinct_interior_edges_on_faceBoundary},
and \path{Theorem49PlanarBoundaryAnnulusGeometry.lean} lifts this to
\path{not_exists_planarBoundaryAnnulusCollarGeometry_of_numCollars_eq_one_and_exists_two_distinct_interior_edges_on_faceBoundary}.
Instantiating those results on the same honest source model gives
\path{not_exists_oneCollar_planarBoundaryAnnulusCollarGeometry_sharedInteriorPair}
in \path{Theorem49PlanarBoundaryAnnulusHonestWitnessRegression.lean}.  So the
remaining repair ``derive weak annulus collar geometry directly from the honest
source'' is now formally blocked on the entire one-collar branch for that
model, not only on the canonical decomposition recovered from it.  Any
workaround now has to produce a genuinely multi-collar geometry or leave the
current weak-collar theorem surface altogether.

That honest-source obstruction first reached the same-split multi-collar
repairs via the current-boundary bridge
\path{not_exists_planarBoundaryAnnulusCollarGeometry_of_not_exists_planarBoundaryAnnulusCurrentBoundaryData_of_numCollars_eq_and_boundaryData}.
The same witness now has a stronger obstruction that no longer fixes the
ambient boundary split.  In
\path{Theorem49PlanarBoundaryAnnulusWitness.lean},
\path{false_of_mem_terminal_collarFaces_and_exists_two_distinct_interior_edges_on_faceBoundary}
shows that a terminal collar face cannot carry two distinct interior edges:
once one of them is the local witness, the other has no deeper frontier on
which to sit and would have to be an ambient boundary edge.  The honest
regression file applies this to the two outer faces of
\path{sharedInteriorPairAnnulusEmbedding}; hence any terminal collar is forced
to be exactly the inner triangular face, formalized as
\path{terminal_collarFaces_eq_singleton_sipFace2}.

That terminal-singleton fact blocks the remaining two- and three-collar weak
geometries even when their boundary split is chosen freely.  If a predecessor
stage has only \path{sipFace2} in its remainder, the theorem
\path{false_of_intermediateBoundary_before_singleton_sipFace2} forces an
intermediate boundary edge to lie both in \path{interiorEdgeSupport} and in the
selected boundary support of the inner triangle, contradicting the global
ambient-boundary/interior disjointness.  This yields
\path{not_exists_twoCollar_planarBoundaryAnnulusCollarGeometry_sharedInteriorPair}
and
\path{not_exists_threeCollar_planarBoundaryAnnulusCollarGeometry_sharedInteriorPair}
without any equality to \path{sharedInteriorPairBoundaryData}.  Together with
the one-collar blocker and the finite-face bound
\path{PlanarBoundaryAnnulusCollarGeometry.numCollars_le_card_faces}, specialized
as
\path{numCollars_le_three_of_planarBoundaryAnnulusCollarGeometry_sharedInteriorPair},
the regression now proves
\path{not_nonempty_planarBoundaryAnnulusCollarGeometry_sharedInteriorPair}.
Thus the packaged honest-source model
\path{closedWalkAnnulusBoundarySource_without_annulusCollarGeometry_sharedInteriorPair}
and its graph-level form
\path{exists_embedding_closedWalkAnnulusBoundarySource_without_annulusCollarGeometry_sharedInteriorPair}
refute any universal same-embedding theorem from honest closed-walk source data
to weak annulus collar geometry on this theorem surface, regardless of how the
annulus boundary split is later chosen.  The corresponding failed converse is
\path{not_forall_exists_planarBoundaryAnnulusCollarGeometry_of_closedWalkAnnulusBoundarySource_sharedInteriorPair}.
Consequently the still stronger source-alone route to the repaired
previous-boundary witness surface is false as well: the same regression now
proves
\path{not_nonempty_planarBoundaryAnnulusPreviousBoundaryWitnessGeometry_sharedInteriorPair},
\path{closedWalkAnnulusBoundarySource_without_previousBoundaryWitness_sharedInteriorPair},
and
\path{not_forall_exists_planarBoundaryAnnulusPreviousBoundaryWitnessGeometry_of_closedWalkAnnulusBoundarySource_sharedInteriorPair}
by projecting any repaired previous-boundary witness geometry back to weak
collar geometry.  This refutes only the unqualified honest-source implication;
the stricter same-embedding question with weak collar geometry already assumed
remains the separate repaired-witness burden.  The repaired-witness
investigation module now removes the earlier placeholder by proving the exact
fixed-collar formulation
\path{honest_source_plus_collar_geometry_extending_previous_boundary_witness_iff_witnessOnCurrentBoundary}:
for a given weak collar geometry, an extension of that very collar to the
previous-boundary surface is equivalent to proving its
\path{WitnessOnCurrentBoundary} property.  The same investigation now proves
that this property is not forced by merely adding an honest source.  On the
connected chained-diamonds source, Lean defines the valid weak two-collar
geometry \path{chainedDiamondsBadTwoCollarGeometry}; its terminal triangular
face chooses \path{cdt97}, while the previous intermediate boundary is
the singleton \path{cdt45}.  The theorems
\path{not_witnessOnCurrentBoundary_chainedDiamondsBadTwoCollarGeometry},
\path{not_forall_honest_source_plus_collar_geometry_forces_previous_boundary_witness},
and
\path{not_nonempty_previousBoundaryWitnessGeometry_extending_chainedDiamondsBadTwoCollarGeometry}
therefore refute the route ``honest source + arbitrary weak collar
\(\Rightarrow\) repaired previous-boundary witness'' on a connected honest
source embedding.  The same file now sharpens this calibration:
\path{exists_atMostOne_closedWalkAnnulusBoundarySource_and_annulusCollarGeometry_without_witnessOnCurrentBoundary_chainedDiamonds},
\path{not_forall_atMostOne_source_plus_collar_geometry_forces_previous_boundary_witness},
and
\path{not_forall_canonicalWitnessChoice_source_plus_collar_geometry_forces_previous_boundary_witness}
show that even the local at-most-one source package, and even an already
inhabited canonical witness choice on the source boundary split, do not make a
separately supplied weak collar satisfy the repaired placement condition.  The
positive route must use the collar constructed from that choice or carry the
placement condition explicitly.

The same model now blocks the weaker witness-cover repair surfaces that the
positive synthesis layer currently accepts.  The theorem
\path{not_nonempty_planarBoundaryHeightOrderedFacePeelWitnessData_sharedInteriorPair}
does not mention annulus boundary splits or collar remainders.  It uses only the
height-ordered witness-cover fields: covering \path{sip01} forces an outer face
to choose \path{sip01}; the other shared interior edge \path{sip12} must then be
chosen by a strictly later peeled face; applying the same argument to that face
forces a still-later \path{sip01} witness.  Since only \path{sipFace0} and
\path{sipFace1} can choose either shared edge, this gives a strict height cycle.
The finite collar-layer form is blocked by
\path{not_nonempty_planarBoundaryCollarLayerFacePeelWitnessData_sharedInteriorPair},
via its canonical lowering to the height-ordered surface.  The obstruction now
also removes the height hypothesis entirely.  The generic list lemma
\path{not_exists_mem_of_twoWitnessPrefixCycle} proves that a finite peel order
cannot contain a two-witness prefix cycle, even if the list repeats faces.  Its
support-level strengthening
\path{not_mem_facePeelWitnessSupport_pair_of_twoWitnessPrefixCycle} now records
that both mutually dependent witness edges are excluded from the realized peel
support.  The regression instantiates this as
\path{not_nonempty_attachedBoundaryRootedFacePeelCertificate_sharedInteriorPair}
for the raw attached-face \path{BoundaryRootedFacePeelCertificate} consumed by
the boundary-root synthesis theorem.  The well-founded-parent witness-cover
surface is then blocked by
\path{not_nonempty_planarBoundaryWellFoundedFacePeelWitnessData_sharedInteriorPair},
since it canonically lowers to that attached certificate.  The source-packaged
counterexamples
\path{closedWalkAnnulusBoundarySource_without_heightOrderedFacePeelWitnessData_sharedInteriorPair}
and
\path{closedWalkAnnulusBoundarySource_without_collarLayerFacePeelWitnessData_sharedInteriorPair}
now have the attached-certificate and well-founded variants
\path{closedWalkAnnulusBoundarySource_without_attachedBoundaryRootedFacePeelCertificate_sharedInteriorPair}
and
\path{closedWalkAnnulusBoundarySource_without_wellFoundedFacePeelWitnessData_sharedInteriorPair}.
This obstruction is now calibrated at the coloring layer as well:
\path{sharedInteriorPairTaitEdgeColoring} gives a concrete Tait edge coloring
of the same graph, proved by
\path{sharedInteriorPairTaitEdgeColoring_isTait}.  The package
\path{closedWalkAnnulusBoundarySource_and_taitEdgeColoring_without_attachedBoundaryRootedFacePeelCertificate_sharedInteriorPair}
therefore combines an honest closed-walk annulus source, an actual Tait
coloring of its graph, and failure of the same-embedding raw attached
certificate.  The obstruction is not an artifact of a non-colorable source
graph; the missing datum is genuinely the same-embedding geometric
certificate.
They show
that the repaired positive endpoint ``honest source plus certificate or
witness-cover data'' has a real extra geometric burden.  Honest closed-walk
annulus source data alone does not force even the raw same-embedding attached
certificate, let alone the weaker finite witness-cover interfaces now sufficient
for Step~2 and Theorem~4.9 synthesis.  The same regression now realizes the
stronger successor-dart boundary source on this embedding as
\path{sharedInteriorPairDartSuccessorCycleGeometry}; the selected-boundary arc
proof is
\path{sharedInteriorPairDartSuccessorCycleGeometry_selectedBoundaryArcOnFace}.
Consequently
\path{exists_embedding_boundaryReachabilityData_and_dartSuccessorCycleEmbeddingData_and_selectedBoundaryArc_without_attachedBoundaryRootedFacePeelCertificate_sharedInteriorPair}
and
\path{not_forall_exists_attachedBoundaryRootedFacePeelCertificate_of_boundaryReachabilityData_and_dartSuccessorCycleEmbeddingData_and_selectedBoundaryArc_sharedInteriorPair}
show that even replacing the closed-walk source by the route-facing
successor-cycle source does not force the raw attached certificate on the same
embedding.  This has now been strengthened by
\path{exists_embedding_boundaryReachabilityData_and_dartSuccessorCycleEmbeddingData_and_selectedBoundaryArc_and_taitEdgeColoring_without_attachedBoundaryRootedFacePeelCertificate_sharedInteriorPair}
and
\path{not_forall_exists_attachedBoundaryRootedFacePeelCertificate_of_boundaryReachabilityData_and_dartSuccessorCycleEmbeddingData_and_selectedBoundaryArc_and_taitEdgeColoring_sharedInteriorPair}:
even the successor-dart source shell plus selected-boundary arcs plus an
actual Tait coloring of the graph does not force the same-embedding attached
certificate.  The stronger bundled witness
\path{exists_embedding_boundaryReachabilityData_and_dartSuccessorCycleEmbeddingData_and_selectedBoundaryArc_and_taitEdgeColoring_without_currentSufficientSameEmbeddingGeometry_sharedInteriorPair}
packages the same source and coloring together with simultaneous failure of
height-ordered witness data, collar-layer witness data, the raw attached
certificate, well-founded witness data, and annulus collar geometry.  The
corresponding disjunctive failed converse
\path{not_forall_exists_some_currentSufficientSameEmbeddingGeometry_of_boundaryReachabilityData_and_dartSuccessorCycleEmbeddingData_and_selectedBoundaryArc_and_taitEdgeColoring_sharedInteriorPair}
shows that even one of those currently sufficient same-embedding geometric
endpoints is not forced by the route-facing source plus Tait coloring.  The
same model now also proves that this is not merely another empty-carrier
phenomenon:
\path{vertex_one_mem_selectedBoundaryInteriorEdgeEndpointVertices_sharedInteriorPair}
exhibits vertex \(1\) as an endpoint of an interior shared edge while avoiding
all selected-boundary edges, and hence
\path{selectedBoundaryInteriorEdgeEndpointVertices_nonempty_sharedInteriorPair}
certifies a nonempty purified carrier.  The strengthened package
\path{exists_embedding_boundaryReachabilityData_and_dartSuccessorCycleEmbeddingData_and_selectedBoundaryArc_and_taitEdgeColoring_and_nonempty_selectedBoundaryInteriorCarrier_without_currentSufficientSameEmbeddingGeometry_sharedInteriorPair}
and failed converse
\path{not_forall_exists_some_currentSufficientSameEmbeddingGeometry_of_boundaryReachabilityData_and_dartSuccessorCycleEmbeddingData_and_selectedBoundaryArc_and_taitEdgeColoring_and_nonempty_selectedBoundaryInteriorCarrier_sharedInteriorPair}
therefore show that honest successor-cycle source data, Tait colorability, and
nonempty selected-boundary interior carrier still do not force any of the
currently sufficient same-embedding geometric endpoints.  The
same successor-source package now has parallel failed-converse endpoints for
the individual upstream surfaces too:
\path{not_forall_exists_planarBoundaryAnnulusCollarGeometry_of_boundaryReachabilityData_and_dartSuccessorCycleEmbeddingData_and_selectedBoundaryArc_sharedInteriorPair},
\path{not_forall_exists_planarBoundaryHeightOrderedFacePeelWitnessData_of_boundaryReachabilityData_and_dartSuccessorCycleEmbeddingData_and_selectedBoundaryArc_sharedInteriorPair},
\path{not_forall_exists_planarBoundaryCollarLayerFacePeelWitnessData_of_boundaryReachabilityData_and_dartSuccessorCycleEmbeddingData_and_selectedBoundaryArc_sharedInteriorPair},
and
\path{not_forall_exists_planarBoundaryWellFoundedFacePeelWitnessData_of_boundaryReachabilityData_and_dartSuccessorCycleEmbeddingData_and_selectedBoundaryArc_sharedInteriorPair}.
Thus the stronger route-facing source shell fails all currently sufficient
same-embedding geometric endpoints in this shared-interior-pair model, not only
the raw attached certificate.  In Lean both the honest closed-walk source
failures and the successor-source failures are now factored through reusable
same-embedding failed-converse lemmas, with generic source principles in
\path{PlanarBoundaryClosedWalkSource.lean} and target-specific wrappers in
\path{Theorem49PlanarBoundaryAnnulusGeometry.lean}; the shared-interior-pair
theorems are instances obtained by supplying the explicit missing-target
witnesses.

\subsection{Closed-walk source versus all-outer roots}

There is now a sharper non-endpoint obstruction at the intersection of the
boundary-order source and the surviving root-distance annulus route.  The new
file
\path{Theorem49ClosedWalkOuterRootObstruction.lean}
proves
\path{false_of_closedWalkAnnulusBoundarySource_and_interiorDualBoundaryRootAdjDistancePeelData_of_rootsOuterBoundary}.
If the annulus boundary split is the honest split obtained from a
\path{PlanarBoundaryClosedWalkAnnulusBoundarySource}, then
\path{InteriorDualBoundaryRootAdjDistancePeelData} cannot have all of its
chosen roots meeting the extracted outer boundary component.

The formal reason is short.  Inner-boundary faces are disjoint from the peeled
interior, hence the root-distance package makes them root faces.  If every
root face touches the extracted outer boundary, an inner-boundary face is also
an outer-boundary face.  But the closed-walk source supplies the facewise
boundary-component separation that makes the extracted outer and inner
boundary-face layers disjoint.  The theorem
\path{false_of_closedWalkAnnulusBoundarySource_and_outerBoundaryRootAdjDistancePeelData_sameBoundaryData}
states the same contradiction for the current
\path{PlanarBoundaryOuterBoundaryRootAdjDistancePeelData} package whenever it
uses the same boundary split as the closed-walk source, and
\path{not_exists_embedding_closedWalkAnnulusBoundarySource_and_outerBoundaryRootAdjDistancePeelData_sameBoundaryData}
packages the graph-level nonexistence claim.  The same file now also proves
\path{false_of_closedWalkAnnulusBoundarySource_and_outerBoundaryRootParentPeelData_sameBoundaryData}
and
\path{not_exists_embedding_closedWalkAnnulusBoundarySource_and_outerBoundaryRootParentPeelData_sameBoundaryData},
so the same-split obstruction already kills the stronger
\path{PlanarBoundaryOuterBoundaryRootParentPeelData} package too, not only the
weaker adjacency-distance interface.

This substantially narrows the surviving route.  The abstract
\path{PlanarBoundaryOuterBoundaryRootAdjDistancePeelData} synthesis package is
still a formally valid theorem surface, but it cannot be instantiated from the
honest closed-walk two-component boundary source while keeping the current
``all roots touch the outer boundary'' hypothesis.  A continuation has to
replace that root condition with a two-sided/root-per-boundary formulation, or
derive a different boundary split not governed by the closed-walk source's
component separation.  The new parent-level obstruction shows that simply
strengthening the same outer-root package downstream does not repair this lane.
The same file now also proves
\path{false_of_boundaryReachabilityData_and_interiorDualBoundaryRootParentPeelData_of_rootsOuterAmbientBoundary_and_boundarySelectedBoundaryArcOnFace},
so once a boundary-order shell already carries facewise selected-boundary arcs,
the stronger outer-root parent package is impossible there even before one asks
for an honest closed-walk annulus source.  The same file now also proves the
same-boundary local nonexistence theorems
\path{not_exists_planarBoundaryOuterBoundaryRootAdjDistancePeelData_sameBoundaryData_of_boundaryReachabilityData_and_boundarySelectedBoundaryArcOnFace}
and
\path{not_exists_planarBoundaryOuterBoundaryRootParentPeelData_sameBoundaryData_of_boundaryReachabilityData_and_boundarySelectedBoundaryArcOnFace},
so the live selected-arc boundary-order shell already rules out both outer-root
packages as soon as they try to reuse that shell's own annulus boundary split.

The first replacement surface is now in Lean.  The file
\path{Theorem49PlanarBoundaryAnnulusRootInteriorDual.lean} defines
\path{PlanarBoundaryAnnulusRootAdjDistancePeelData} and
\path{PlanarBoundaryAnnulusRootParentPeelData}, where each root face may meet
either the extracted outer or inner ambient boundary component.  The
constructors
\path{planarBoundaryAnnulusRootAdjDistancePeelData_of_boundaryData_and_interiorDualBoundaryRootAdjDistancePeelData}
and
\path{planarBoundaryAnnulusRootAdjDistancePeelDataOfCoverSeparatedAnnulusBoundaryRootsCanonicalRootedSharedEdgeFaceMembershipCover}
show that this is not a new obstruction layer: an annulus boundary split plus
generic boundary-root interior-dual data already gives the two-sided package,
and raw cover/separated-root data can use the full outer/inner boundary union
as its root support.  The closed-walk constructors
\path{planarBoundaryAnnulusRootAdjDistancePeelData_of_closedWalkAnnulusBoundarySource_and_interiorDualBoundaryRootAdjDistancePeelData}
and
\path{admitsPlanarBoundaryAnnulusRootAdjDistancePeelData_of_exists_closedWalkAnnulusBoundarySource_and_interiorDualBoundaryRootAdjDistancePeelData}
make the positive companion explicit: the honest closed-walk boundary source
can feed the two-sided root surface, provided the generic interior-dual data is
available.  The lowering theorems
\path{admitsPlanarBoundaryInteriorDualBoundaryRootAdjDistancePeelData_of_admitsPlanarBoundaryAnnulusRootAdjDistancePeelData}
and
\path{admitsPlanarBoundaryCollarLayerFacePeelWitnessData_of_admitsPlanarBoundaryAnnulusRootAdjDistancePeelData}
connect the new surface back to the existing collar-layer and Theorem~4.9
pipeline.  The route file
\path{PlanarBoundaryClosedWalkSourceRoute.lean} now also exposes this positive
path all the way to the purified theorem-4.9 endpoint:
\path{exists_theorem49AnnulusRootSynthesis_of_exists_closedWalkAnnulusBoundarySource_and_interiorDualBoundaryRootAdjDistancePeelData_direct}
and its successor-cycle analogue
\path{exists_theorem49AnnulusRootSynthesis_of_exists_boundaryReachabilityData_and_dartSuccessorCycleEmbeddingData_and_selectedBoundaryArc_and_interiorDualBoundaryRootAdjDistancePeelData_direct}
show that, on one embedding, honest boundary-order source data plus generic
interior-dual data already reaches the two-sided synthesis surface without
restating the older shifted-construction side conditions.
The same file now also exposes the more direct certificate-side corollaries
\path{admitsPlanarBoundaryAttachedRootedFacePeelCertificate_of_exists_closedWalkAnnulusBoundarySource_and_interiorDualBoundaryRootAdjDistancePeelData_direct},
\path{exists_theorem49BoundaryRootSynthesis_of_exists_closedWalkAnnulusBoundarySource_and_interiorDualBoundaryRootAdjDistancePeelData_direct},
and their successor-cycle analogues, showing that once the generic
boundary-root interior-dual package is present the honest source shells can
land directly on the attached-certificate and boundary-root synthesis surfaces
without routing back through annulus-construction parent orientation.  The same
successor-cycle shell now also reaches the non-vacuous two-sided annulus-root
surface:
\path{theorem49AnnulusRootNonemptySynthesis_of_boundaryReachabilityData_and_dartSuccessorCycleEmbeddingData_and_selectedBoundaryArc_and_interiorDualBoundaryRootAdjDistancePeelData_of_peelFaces_endpoint_disjoint_selectedBoundarySupport}
and
\path{exists_theorem49AnnulusRootNonemptySynthesis_of_exists_boundaryReachabilityData_and_dartSuccessorCycleEmbeddingData_and_selectedBoundaryArc_and_interiorDualBoundaryRootAdjDistancePeelData_with_peelFaces_endpoint_disjoint_selectedBoundarySupport_and_nonempty_interiorEdgeEndpointSupport}
show that if the canonical annulus construction on that same embedding has
peeled-face endpoint no-touch and the raw interior-edge endpoint carrier is
nonempty, then the live boundary-order shell already reaches
\path{Theorem49AnnulusRootNonemptySynthesis} before any source-fixed parent
package is introduced.
The same live successor-cycle shell now also reaches the generic source-fixed
IDBRAD surface itself:
\path{interiorDualBoundaryRootAdjDistancePeelData_of_boundaryReachabilityData_and_dartSuccessorCycleEmbeddingData_and_selectedBoundaryArc_and_boundaryFaceRootsCanonicalRootedSharedEdgeFaceMembershipCover}
is the same-embedding constructor, and
\path{admitsPlanarBoundaryInteriorDualBoundaryRootAdjDistancePeelData_of_exists_boundaryReachabilityData_and_dartSuccessorCycleEmbeddingData_and_selectedBoundaryArc_and_boundaryFaceRootsCanonicalRootedSharedEdgeFaceMembershipCover}
is the graph-level route theorem.  So once the successor-cycle boundary-order
shell supplies the source boundary-face root cover/separation, boundary-free
peeled faces, rooted shared-edge coverage, and strict root-distance
child-membership fields, no additional source-to-IDBRAD repackaging step
remains.  The same rooted-cover package now also closes directly at the
theorem-4.9 endpoint through
\path{exists_theorem49BoundaryRootSynthesis_of_exists_boundaryReachabilityData_and_dartSuccessorCycleEmbeddingData_and_selectedBoundaryArc_and_boundaryFaceRootsCanonicalRootedSharedEdgeFaceMembershipCover_direct},
so no further theorem-4.9-side repackaging remains below that concrete
source-facing surface.
The annulus-root file now also records the basic two-sided geometric forcing
fact that was missing from the earlier route story: every ambient face meeting
either distinguished annulus boundary component is automatically a non-peeled
face and therefore a root face for the generic interior-dual adjacency-distance
package.  In particular,
\path{exists_root_mem_outerBoundaryFaces_of_closedWalkAnnulusBoundarySource_and_interiorDualBoundaryRootAdjDistancePeelData}
and
\path{exists_root_mem_innerBoundaryFaces_of_closedWalkAnnulusBoundarySource_and_interiorDualBoundaryRootAdjDistancePeelData}
show that an honest closed-walk annulus source plus the generic
boundary-root interior-dual data forces at least one chosen root face on each
boundary component of the extracted annulus split.
The same file now also exposes the upstream parent-orientation surface itself:
\path{parentWitnessOrientation_of_closedWalkAnnulusBoundarySource_and_interiorDualBoundaryRootAdjDistancePeelData_rootDistance}
is the same-embedding theorem, and the graph-level wrappers
\path{admitsPlanarBoundaryAnnulusConstructionParentWitnessOrientation_of_exists_closedWalkAnnulusBoundarySource_and_interiorDualBoundaryRootAdjDistancePeelData_direct}
and
\path{admitsPlanarBoundaryAnnulusConstructionParentWitnessOrientation_of_exists_boundaryReachabilityData_and_dartSuccessorCycleEmbeddingData_and_selectedBoundaryArc_and_interiorDualBoundaryRootAdjDistancePeelData_direct}
show that honest closed-walk and successor-cycle boundary-order sources already
reach \path{PlanarBoundaryAnnulusConstruction.ParentWitnessOrientation}
through the root-distance construction, before invoking annihilator or
synthesis theorems.

The converse route is now explicitly blocked on the diamond-with-triangle
benchmark.  The new regression theorem
\path{diamondWithTriangle_source_previousBoundaryWitness_and_witnessCover_without_interiorDual_rootDistance}
packages, on one embedding, an honest
\path{PlanarBoundaryClosedWalkAnnulusBoundarySource}, repaired
\path{PlanarBoundaryAnnulusPreviousBoundaryWitnessGeometry}, well-founded and
height-ordered witness-cover data, and a
\path{ForcingInteriorEdgeWitness}; the same theorem also proves absence of both
\path{InteriorDualBoundaryRootAdjDistancePeelData} and the two-sided
\path{PlanarBoundaryAnnulusRootAdjDistancePeelData}.  The failed converses
\path{not_forall_interiorDualBoundaryRootAdjDistancePeelData_of_previousBoundaryWitnessGeometry_diamondWithTriangle}
and
\path{not_forall_planarBoundaryAnnulusRootAdjDistancePeelData_of_closedWalkAnnulusBoundarySource_and_previousBoundaryWitnessGeometry_diamondWithTriangle}
record the route implication directly: repaired previous-boundary witness
geometry, even with the honest closed-walk source present, is not enough to
derive the interior-dual root-distance route data.  Thus root-distance data
should be treated as a separate sufficient construction input, not as a hidden
consequence of the current repaired witness geometry.

On the annulus-geometry side, the development now also exposes the hidden
parent-carrying step inside the previous-boundary witness proof itself.
\path{Theorem49PlanarBoundaryAnnulusDecomposition.lean} defines the canonical
collar index
\path{PlanarBoundaryAnnulusDecomposition.layerOf},
while
\path{Theorem49PlanarBoundaryAnnulusWitness.lean}
proves
\path{PlanarBoundaryAnnulusWitnessAssignment.exists_parentFace_of_previousBoundaryWitness}
and packages the chosen parent map as
\path{PlanarBoundaryAnnulusWitnessAssignment.parentFaceOfPreviousBoundaryWitness}.
So the current live geometric gap is no longer hidden inside the theorem-4.9
endpoint proof: the remaining task is to derive honest previous-boundary
witness placement, or a stronger replacement such as annular no-touch, from
real boundary-order geometry.

\subsection{Boundary-order gap}

The note
\path{fourcolor-boundary-order-notes.tex}
records a separate structural obstruction: the current
\path{PlaneEmbeddingWithBoundary} API only exposes unordered face-boundary
incidence, while the downstream selected-boundary and annulus arguments require
ordered or cyclic face-boundary geometry.

This gap is not merely prose.  The Lean development includes explicit
regressions showing that unordered incidence alone is too weak to recover the
ordered walk data needed downstream.  The current workaround is to introduce
stronger ordered and cyclic source interfaces explicitly, rather than pretending
that the older API already determines them.

\subsection{Interior-vertex discharge}

The later quotient and boundary-erasure bridges in
\path{Theorem49BoundaryProjection.lean} all carry the same side condition: the
vertices where Kirchhoff equations are imposed must not be incident to the
selected boundary support that is erased.  The reviewer escalation around this
point is valid.  Without this condition, erasing selected-boundary coordinates
can change a vertex Kirchhoff sum.

The Lean development now factors this issue into a small checked bridge in
\path{Theorem49InteriorVertices.lean}.  Given any finite candidate vertex set,
\path{selectedBoundaryInteriorVertices} filters out the vertices incident to
selected-boundary edges.  The theorem
\path{selectedBoundaryInteriorVertices_avoids_selectedBoundarySupport}
discharges the side condition for this purified set.  The exact obstruction is
also recorded: \path{verticesAvoidingEdgeSupport_eq_self_iff} says that the
purification is the identity precisely when the original candidate set already
avoids the erased edge support, and
\path{not_boundary_avoids_vertices_of_incident_boundary_edge} shows that a
single candidate vertex incident to a selected-boundary edge falsifies the side
condition.

The route-facing payoff is deliberately narrow rather than another quotient API
packet.  First,
\path{theorem49BoundaryZeroKirchhoffSubspace_selectedBoundaryInteriorVertices_eq_map_boundaryZeroProjection_kirchhoffSubmodule}
discharges the side condition in the boundary-erasure image form of \(W_0(H)\):
for the purified vertex set, the concrete boundary-zero Kirchhoff target is
exactly the selected-boundary-erased Kirchhoff image.  Second, the theorem
\path{theorem49BoundaryZeroKirchhoffSubspace_selectedBoundaryInteriorVertices_chainDot_projectedKempeClosureGeneratorSubspace_eq_zero_iff_of_planarBoundaryInteriorDualBoundaryRootAdjDistancePeelData}
specializes the current \(W_0(H)\) separation endpoint to the purified
interior-vertex set, and the corresponding graph-level existential theorem
exposes the same discharge from
\path{AdmitsPlanarBoundaryInteriorDualBoundaryRootAdjDistancePeelData}.

The annulus route now has a same-statement bridge as well:
\path{theorem49BoundaryZeroKirchhoffSubspace_selectedBoundaryInteriorVertices_image_and_separation_of_planarBoundaryOuterBoundaryRootAdjDistancePeelData}
takes the outer-boundary-rooted Step~2 package and returns both the purified
boundary-erasure image identity and the purified separation endpoint.  The
graph-level theorem
\path{exists_theorem49BoundaryZeroKirchhoffSubspace_selectedBoundaryInteriorVertices_image_and_separation_of_admitsPlanarBoundaryOuterBoundaryRootAdjDistancePeelData}
exposes this directly from the paper-facing annulus target.

The arbitrary candidate has now been replaced at one route endpoint by a
concrete finite graph carrier.  The definitions \path{edgeEndpointSupport} and
\path{graphEdgeEndpointSupport} collect all endpoints of graph edges from a
finite edge set, without assuming a global finite vertex type.  The purified
carrier \path{selectedBoundaryGraphInteriorVertices} then deletes endpoints of
the selected boundary support.  The theorem
\path{theorem49BoundaryZeroKirchhoffSubspace_selectedBoundaryGraphInteriorVertices_image_and_separation_of_planarBoundaryOuterBoundaryRootAdjDistancePeelData}
specializes the annulus bridge to this concrete carrier, and the graph-level
theorem
\path{exists_theorem49BoundaryZeroKirchhoffSubspace_selectedBoundaryGraphInteriorVertices_image_and_separation_of_admitsPlanarBoundaryOuterBoundaryRootAdjDistancePeelData}
exposes it without a user-supplied \path{verticesOf} function.  This is a real
interior-vertex route integration, not a further quotient API extension.

The next refinement uses the annulus route's own face-incidence edge
classification.  The definition \path{interiorEdgeEndpointSupport} collects
endpoints of \path{interiorEdgeSupport} edges, i.e. edges incident to two
ambient faces, and
\path{selectedBoundaryInteriorEdgeEndpointVertices} purifies this more
annular carrier.  The theorem
\path{theorem49BoundaryZeroKirchhoffSubspace_selectedBoundaryInteriorEdgeEndpointVertices_image_and_separation_of_planarBoundaryOuterBoundaryRootAdjDistancePeelData}
specializes the same boundary-erasure and separation endpoint to this carrier,
with graph-level wrapper
\path{exists_theorem49BoundaryZeroKirchhoffSubspace_selectedBoundaryInteriorEdgeEndpointVertices_image_and_separation_of_admitsPlanarBoundaryOuterBoundaryRootAdjDistancePeelData}.
The same surviving carrier is now exposed in map-level and dimension forms.
The theorem
\path{ker_theorem49BoundaryZeroKirchhoffSubspaceProjectedGeneratorPairingMap_selectedBoundaryInteriorEdgeEndpointVertices_eq_bot_of_planarBoundaryOuterBoundaryRootAdjDistancePeelData}
states that the concrete pairing map from this purified \(W_0(H)\) target into
the dual of the projected Definition~4.8 generator span has trivial kernel.
Consequently
\path{finrank_theorem49BoundaryZeroKirchhoffSubspace_selectedBoundaryInteriorEdgeEndpointVertices_le_finrank_projectedKempeClosureGeneratorSubspace_of_planarBoundaryOuterBoundaryRootAdjDistancePeelData}
gives the finite-dimensional bound
\(\dim W_0(H)_{\mathrm{purified}} \le \dim K_{\mathrm{proj}}\), where
\(K_{\mathrm{proj}}\) is the projected Definition~4.8 generator span.
Both statements also have graph-level existential wrappers from
\path{AdmitsPlanarBoundaryOuterBoundaryRootAdjDistancePeelData}.
The same purified carrier now reaches the operational projected-source form of
Theorem~4.9.  The theorem
\path{theorem49BoundaryZeroKirchhoffSubspace_selectedBoundaryInteriorEdgeEndpointVertices_eq_kirchhoffSubmodule_inf_projectedKempeClosureGeneratorSubspace_of_planarBoundaryOuterBoundaryRootAdjDistancePeelData}
identifies this \(W_0(H)\) target with the Kirchhoff equations at the purified
carrier intersected with \(K_{\mathrm{proj}}\), and
\path{theorem49BoundaryZeroKirchhoffSubspace_selectedBoundaryInteriorEdgeEndpointVertices_le_projectedKempeClosureGeneratorSubspace_of_planarBoundaryOuterBoundaryRootAdjDistancePeelData}
gives the direct inclusion into \(K_{\mathrm{proj}}\).  Finally,
\path{exists_projectedKempeClosureGeneratorFamily_finset_sum_of_mem_theorem49BoundaryZeroKirchhoffSubspace_selectedBoundaryInteriorEdgeEndpointVertices_of_planarBoundaryOuterBoundaryRootAdjDistancePeelData}
and
\path{exists_kempeClosureGeneratorFamily_finset_sum_boundaryZeroProjection_and_mem_kirchhoffSubmodule_of_mem_theorem49BoundaryZeroKirchhoffSubspace_selectedBoundaryInteriorEdgeEndpointVertices_of_planarBoundaryOuterBoundaryRootAdjDistancePeelData}
turn membership in the purified \(W_0(H)\) target into finite projected
Definition~4.8 generator sums, respectively finite raw generator sums whose
selected-boundary projection is the target chain and whose raw sum remains
Kirchhoff at the purified carrier.  The graph-level wrappers expose these
finite-combination endpoints from
\path{AdmitsPlanarBoundaryOuterBoundaryRootAdjDistancePeelData}.
The raw-source quotient form is now specialized to the same carrier as well.
\path{theorem49BoundaryZeroKirchhoffSubspace_selectedBoundaryInteriorEdgeEndpointVertices_eq_map_boundaryZeroProjection_kempeClosureGeneratorSubspace_inf_kirchhoffSubmodule_of_planarBoundaryOuterBoundaryRootAdjDistancePeelData}
identifies the purified \(W_0(H)\) target with the selected-boundary erasure of
raw Definition~4.8 generator-span chains that are already Kirchhoff at the
purified carrier.  The preimage theorem
\path{exists_kempeClosureGeneratorSubspace_inf_kirchhoffSubmodule_preimage_boundaryZeroProjection_of_mem_theorem49BoundaryZeroKirchhoffSubspace_selectedBoundaryInteriorEdgeEndpointVertices_of_planarBoundaryOuterBoundaryRootAdjDistancePeelData}
packages this as an explicit source representative in
``raw Kempe span'' intersected with the purified Kirchhoff submodule.  Finally,
\path{exists_kempeClosureGeneratorFamily_finset_sum_boundaryKernel_error_decomposition_of_mem_kirchhoffSubmodule_selectedBoundaryInteriorEdgeEndpointVertices_of_planarBoundaryOuterBoundaryRootAdjDistancePeelData}
records the quotient error term: every raw Kirchhoff chain at the purified
carrier is a finite raw generator sum plus a Kirchhoff chain killed by
selected-boundary erasure.
The new synthesis module \path{Theorem49Synthesis.lean} packages these
positive consequences without adding another carrier layer.  Its core
structure \path{Theorem49BoundaryRootSynthesis} is proved directly from generic
\path{PlanarBoundaryInteriorDualBoundaryRootAdjDistancePeelData}; it collects
the boundary-erased image identity, pointwise separation, kernel-zero and
finite-dimensional forms, projected-generator spanning, and the corrected
raw-source quotient and boundary-kernel decomposition for the selected-
boundary interior-edge endpoint carrier.  The outer-route theorem
\path{exists_theorem49OuterBoundaryRootSynthesis_of_admitsPlanarBoundaryOuterBoundaryRootAdjDistancePeelData}
is now just a wrapper through the outer package's \path{interiorData}, and the
same file also exposes the identical synthesis endpoint on the two-sided
surface via \path{Theorem49AnnulusRootSynthesis} and
\path{exists_theorem49AnnulusRootSynthesis_of_admitsPlanarBoundaryAnnulusRootAdjDistancePeelData}.
So the purified Theorem~4.9 algebra no longer depends on the all-outer-root
localization once generic boundary-root interior-dual data exists.  The same
module now also exposes the primitive two-sided annulus-split interface:
\path{theorem49AnnulusRootSynthesis_of_boundaryData_and_interiorDualBoundaryRootAdjDistancePeelData}
and its graph-level existence wrappers show that once one embedding carries
\path{PlanarBoundaryAnnulusBoundaryData} together with generic
\path{InteriorDualBoundaryRootAdjDistancePeelData}, the theorem-4.9 synthesis
endpoint is already available before any stronger closed-walk source fields are
added.
The next bridge now also exists directly from the concrete annulus
construction surface.  The new file
\path{Theorem49PlanarBoundaryAnnulusConstructionSpanning.lean} proves
\path{boundaryZeroAnnihilatorTrivialForEmbedding_of_planarBoundaryAnnulusConstruction_of_parentWitness}:
if a \path{PlanarBoundaryAnnulusConstruction} satisfies the explicit
\path{ParentWitnessOrientation} condition, then every selected-boundary-zero
chain annihilating the Definition~4.8 family is already zero on all edges.
This is the exact non-geometric hypothesis used by the algebraic spanning
bridge in \path{Theorem49SpanningGap.lean}.  Consequently the same file now
derives
\path{projectedKempeClosureGeneratorSubspace_eq_planarBoundaryZeroSubmodule_of_planarBoundaryAnnulusConstruction_of_parentWitness},
its explicit-family companion
\path{span_projectedKempeClosureGeneratorFamily_eq_planarBoundaryZeroSubmodule_of_planarBoundaryAnnulusConstruction_of_parentWitness},
the corrected target identity
\path{theorem49BoundaryZeroKirchhoffSubspace_eq_kirchhoffSubmodule_inf_projectedKempeClosureGeneratorSubspace_of_planarBoundaryAnnulusConstruction_of_parentWitness},
the corresponding target inclusion, pointwise separation on the purified
endpoint carrier, kernel-triviality of the pairing map, the resulting finrank
bound, and a raw Kempe-span preimage theorem.  So the current annulus route no
longer needs to recreate the older interior-dual boundary-root package merely
to reach the algebraic theorem-4.9 bridge: construction-side geometry plus the
honest parent-witness orientation already suffices.
\path{Theorem49SpanningGap.lean} now exposes the same algebraic bridge directly
from the weaker boundary-aware witness-cover surfaces.  The theorems
\path{boundaryZeroAnnihilatorTrivialForEmbedding_of_planarBoundaryHeightOrderedFacePeelWitnessData}
and
\path{boundaryZeroAnnihilatorTrivialForEmbedding_of_planarBoundaryCollarLayerFacePeelWitnessData}
construct the annihilator endpoint from height-ordered or finite collar-layer
witness-cover data.  The same file now also has graph-level height-ordered and
collar-layer wrappers for both the ordinary and relative submodule-duality
interfaces, with chain-dot variants that explicitly construct the relevant
duality data from the named same-embedding subspaces and identification lemmas.
\path{Theorem49KempeGeneratorSpan.lean} then fixes the
generator side canonically through
\path{theorem49_kempeGeneratorSubspace_spanning_of_planarBoundaryHeightOrderedFacePeelWitnessData}
and its collar-layer analogue.  It now also exposes the graph-level wrappers
\path{exists_theorem49_kempeGeneratorSubspace_spanning_of_admitsPlanarBoundaryHeightOrderedFacePeelWitnessData}
and
\path{exists_theorem49_kempeGeneratorSubspace_spanning_of_admitsPlanarBoundaryCollarLayerFacePeelWitnessData}.
These instantiate the concrete submodule-duality data with the actual
Definition~4.8 Kempe-generator subspace, so the remaining caller-side algebraic
inputs are just target containment and selected-boundary-zero interpretation,
not a separate orthogonality-to-annihilator proof.  Finally
\path{Theorem49BoundaryProjection.lean} pushes both surfaces to the corrected
projected-source conclusion:
\path{projectedKempeClosureGeneratorSubspace_eq_planarBoundaryZeroSubmodule_of_planarBoundaryHeightOrderedFacePeelWitnessData}
and
\path{projectedKempeClosureGeneratorSubspace_eq_planarBoundaryZeroSubmodule_of_planarBoundaryCollarLayerFacePeelWitnessData},
together with the corresponding boundary-zero Kirchhoff target equality and
inclusion theorems.  Thus the spanning gap no longer waits for the stronger
interior-dual rooted surface once either boundary-aware witness-cover surface is
available.  The same file now gives operational finite-witness forms for these
weaker surfaces: every selected-boundary-zero chain, and every concrete
boundary-zero Kirchhoff target chain, has a finite representation by projected
Definition~4.8 generators.  It also gives raw-generator preimage and finite
projection statements before the boundary erasure.  The graph-level existence
wrappers from
\path{AdmitsPlanarBoundaryHeightOrderedFacePeelWitnessData} and
\path{AdmitsPlanarBoundaryCollarLayerFacePeelWitnessData} now include the
explicit-family target identity, projected and explicit-family inclusion
forms, raw Definition~4.8 generator-span preimages, and finite raw-generator
projection witnesses for every boundary-zero Kirchhoff target chain.  The same
projected-source layer now also proves the chain-dot separation form
\[
  W_0(H) \cap (\mathrm{projectedKempeClosureGeneratorSubspace})^\perp
  = 0
\]
from the boundary-zero annihilator endpoint, together with detector and
pointwise orthogonality iff forms.  These theorems have height-ordered,
finite-collar, and graph-level admissibility wrappers, so the projected
Definition~4.8 source supplies actual nondegenerate testing against the
concrete boundary-zero Kirchhoff target without the additional quotient-side
interior-vertex disjointness assumptions required by the raw-source quotient
presentation.  Downstream
theorem 4.9 route files can therefore consume height-ordered or collar-layer
witness data at the corrected projected-source interface without unpacking the
same-embedding witness manually.
\path{Theorem49PlanarBoundaryCollarPeeling.lean} now also exposes the weakest
explicit local-remainder collar surface directly at the generator-family
annihilator endpoint.  The fixed-embedding theorem
\path{zero_on_allEdges_of_planarBoundaryLocalRemainderBoundaryCollarLayerFacePeelWitnessData_of_annihilatesKempeClosureGeneratorFamily}
and the existential companion
\path{zero_on_allEdges_of_exists_planarBoundaryLocalRemainderBoundaryCollarLayerFacePeelWitnessData_of_annihilatesKempeClosureGeneratorFamily}
consume Definition~4.8 generator-family orthogonality without first forgetting
the local-remainder data to ordinary collar-layer data.  Thus the root-free
local-remainder route is now a theorem-facing Step~2 endpoint, not merely an
intermediate representation of finite collars.
\path{Theorem49SpanningGap.lean} and
\path{Theorem49BoundaryProjection.lean} now push this endpoint through the
linear-duality layer as well.  The theorem
\path{boundaryZeroAnnihilatorTrivialForEmbedding_of_planarBoundaryLocalRemainderBoundaryCollarLayerFacePeelWitnessData}
feeds the local-remainder package into the named annihilator-trivial bridge, and
\path{exists_projectedKempeClosureGeneratorFamily_finset_sum_of_mem_theorem49BoundaryZeroKirchhoffSubspace_of_planarBoundaryLocalRemainderBoundaryCollarLayerFacePeelWitnessData}
then gives the concrete projected Definition~4.8 finite-combination form for
every boundary-zero Kirchhoff target chain on the same embedding.  The graph
level companion
\path{exists_projectedKempeClosureGeneratorFamily_finset_sum_of_mem_theorem49BoundaryZeroKirchhoffSubspace_of_admitsPlanarBoundaryLocalRemainderBoundaryCollarLayerFacePeelWitnessData}
records the same endpoint without first converting the local-remainder witness
to an ordinary collar witness.
The same local-remainder branch now also reaches the raw manuscript-facing
source language.  The fixed-embedding theorem
\path{exists_kempeClosureGeneratorFamily_finset_sum_boundaryZeroProjection_of_mem_theorem49BoundaryZeroKirchhoffSubspace_of_planarBoundaryLocalRemainderBoundaryCollarLayerFacePeelWitnessData}
chooses a finite sum from the unprojected Definition~4.8 generator family whose
selected-boundary projection is the target chain, and
\path{exists_kempeClosureGeneratorSubspace_preimage_boundaryZeroProjection_of_mem_theorem49BoundaryZeroKirchhoffSubspace_of_planarBoundaryLocalRemainderBoundaryCollarLayerFacePeelWitnessData}
records the corresponding raw generator-subspace preimage.  Thus the
local-remainder route is not merely a boundary-erased projected-source result:
it now lands back on the raw generator family used by the proof route.
The fixed-embedding strengthening
\path{exists_kempeClosureGeneratorFamily_finset_sum_boundaryZeroProjection_and_mem_kirchhoffSubmodule_of_mem_theorem49BoundaryZeroKirchhoffSubspace_of_planarBoundaryLocalRemainderBoundaryCollarLayerFacePeelWitnessData}
adds the missing Kirchhoff-retention clause: when the chosen target vertices
avoid the selected boundary support, the raw finite Definition~4.8 sum can be
chosen inside the relevant Kirchhoff submodule, and its selected-boundary
projection is still the prescribed boundary-zero target chain.  The graph-level
version records the same endpoint for an
\path{AdmitsPlanarBoundaryLocalRemainderBoundaryCollarLayerFacePeelWitnessData}
witness without passing through the blocked interior-dual root-distance package.
This local-remainder branch now also has the quotient/error endpoint formerly
reserved for the interior-dual package.  The theorem
\path{exists_kempeClosureGeneratorFamily_finset_sum_boundaryZeroProjection_eq_boundaryZeroProjection_of_mem_kirchhoffSubmodule_of_planarBoundaryLocalRemainderBoundaryCollarLayerFacePeelWitnessData}
starts from an arbitrary raw Kirchhoff chain and chooses a finite raw
Definition~4.8 generator-family sum, still Kirchhoff at the chosen target
vertices, with the same selected-boundary-erased image.  The subsequent
boundary-supported and kernel-decomposition theorems record that the difference
is a Kirchhoff correction supported on the selected boundary, equivalently an
element of the boundary-erasure kernel inside the Kirchhoff submodule.  The
graph-level companions expose the same quotient/error package directly from
\path{AdmitsPlanarBoundaryLocalRemainderBoundaryCollarLayerFacePeelWitnessData}.
The route-facing projected synthesis endpoint now exposes the same package
without asking downstream source arguments to reach back into the raw projection
layer:
\path{Theorem49BoundaryRootNonemptyProjectedSynthesis.rawKirchhoffRepresentative_and_boundaryKernelDecomposition}
combines the Kirchhoff-valid raw representative and the boundary-kernel error
decomposition already carried by the nested
\path{Theorem49BoundaryRootSynthesis}.  Thus any closed-walk or collar source
package that has reached the positive same-embedding projected endpoint can
consume the quotient/error interpretation of the raw Definition~4.8
generators directly at that endpoint.
This endpoint has now been named as the abbreviation
\path{Theorem49BoundaryRawQuotientErrorPackage}.  The viable replacement
positive surfaces consume it directly:
\path{theorem49BoundaryRawQuotientErrorPackage_of_heightOrderedPositiveProjectedGeometryOn}
and
\path{theorem49BoundaryRawQuotientErrorPackage_of_collarLayerPositiveProjectedGeometryOn}
give fixed-embedding height-ordered and finite collar-layer versions, while the
namespace methods
\path{Theorem49HeightOrderedPositiveProjectedGeometry.rawQuotientErrorPackage}
and
\path{Theorem49CollarLayerPositiveProjectedGeometry.rawQuotientErrorPackage}
give the graph-level package forms.  Thus the route no longer merely produces a
nonempty projected synthesis package from the replacement geometry; it also
exposes the raw Definition~4.8 quotient/error reading at that same replacement
surface.
The core synthesis surface now also records the corresponding global submodule
decomposition.  The theorem
\path{Theorem49BoundaryRootSynthesis.kirchhoffSubmodule_eq_rawSource_sup_boundaryKernel}
states that, once \path{Theorem49BoundaryRootSynthesis} is available, the
purified Kirchhoff submodule is exactly the supremum of the Kirchhoff-valid raw
Definition~4.8 Kempe-source submodule and the selected-boundary projection
kernel inside that purified Kirchhoff submodule.  This is the synthesis-level
form of the quotient/error conclusion: every purified Kirchhoff chain splits
into a raw source component plus a boundary-kernel error, and the reverse
inclusion is immediate because both summands remain Kirchhoff-valid.
The quotient/surjectivity form is now explicit as
\path{Theorem49BoundaryRootSynthesis.rawSource_mkQ_eq_top}.  When the raw
source is viewed as a submodule of the purified Kirchhoff submodule and the
selected-boundary projection kernel is restricted to that Kirchhoff submodule,
the quotient map sends the raw source onto the top submodule of the quotient.
Equivalently, every purified Kirchhoff quotient class modulo selected-boundary
kernel error has a raw Definition~4.8 Kempe-source representative.  The
statement is deliberately internal to the purified Kirchhoff submodule, avoiding
the stronger and generally unjustified ambient quotient claim.
The same quotient endpoint now has a finrank consequence:
\path{Theorem49BoundaryRootSynthesis.finrank_quotient_le_rawSource} bounds the
dimension of the purified Kirchhoff quotient by the dimension of the
Kirchhoff-valid raw Definition~4.8 source submodule inside it.  This is just the
finite-dimensional shadow of the quotient surjectivity, but it exposes a
rank-level invariant that downstream theorem-4.9 arguments can use without
opening the quotient-map proof.
The target-side version is now explicit as
\path{Theorem49BoundaryRootSynthesis.finrank_target_le_rawSource}: the concrete
boundary-zero Kirchhoff target itself has finrank bounded by the
Kirchhoff-valid raw Definition~4.8 source submodule whose selected-boundary
erasure is the target.  This packages the dimension consequence of
\path{target_eq_rawSourceImage} directly at the synthesis surface.
The quotient-target bridge itself is now named at this surface:
\path{Theorem49BoundaryRootSynthesis.boundaryProjectionQuotientEquivTarget}
identifies the purified Kirchhoff submodule modulo the restricted
selected-boundary projection kernel with the concrete boundary-zero Kirchhoff
target.  Its computation theorem
\path{Theorem49BoundaryRootSynthesis.boundaryProjectionQuotientEquivTarget_apply_mk}
sends the quotient class of a purified Kirchhoff chain to its
selected-boundary-erased chain.  Thus the corrected quotient/error layer is no
longer merely an interpretation of representatives; it is a first-isomorphism
interface for the actual theorem-4.9 target.
The source-side first-isomorphism form is now present as well:
\path{Theorem49BoundaryRootSynthesis.rawSourceQuotientEquivTarget} identifies
the Kirchhoff-valid raw Definition~4.8 source submodule modulo the
selected-boundary projection kernel on that source with the same concrete
boundary-zero Kirchhoff target.  Its computation theorem
\path{Theorem49BoundaryRootSynthesis.rawSourceQuotientEquivTarget_apply_mk}
sends the class of a raw Kirchhoff-valid source chain to its
selected-boundary-erased chain.  This gives the target a direct raw-source
quotient model, separate from the quotient of the entire purified Kirchhoff
submodule.
The two quotient models are now connected directly:
\path{Theorem49BoundaryRootSynthesis.rawSourceQuotientEquivKirchhoffQuotient}
identifies the raw-source quotient with the purified Kirchhoff quotient through
their common concrete theorem-4.9 target.  Its computation theorem
\path{Theorem49BoundaryRootSynthesis.rawSourceQuotientEquivKirchhoffQuotient_apply_mk}
sends the class of a raw Kirchhoff-valid source representative to the class of
the same chain regarded as an element of the purified Kirchhoff submodule.  This
is the synthesis-surface comparison map needed to pass between raw
Definition~4.8 quotient classes and purified Kirchhoff quotient classes without
reopening either first-isomorphism proof.
The representative-level endpoint is now explicit:
\path{Theorem49BoundaryRootSynthesis.exists_rawGeneratorSum_for_kirchhoffQuotientClass}
proves that every class of the purified Kirchhoff quotient has a representative
that is both Kirchhoff-valid and a finite linear combination of actual raw
Definition~4.8 generator-family terms.  The theorem packages the selected-
boundary projection equality as equality of quotient classes modulo the
restricted selected-boundary kernel.  Downstream arguments can therefore choose
a finite raw generator-sum representative for a purified Kirchhoff quotient
class without replaying the raw representative construction or the quotient
kernel conversion.
The separation endpoint has also been transported to the purified quotient:
\path{Theorem49BoundaryRootSynthesis.kirchhoffQuotientClass_eq_zero_iff_projectedGenerator_orthogonal}
states that the class of a purified Kirchhoff chain is zero modulo the
restricted selected-boundary kernel exactly when its selected-boundary-erased
representative pairs to zero with every projected Definition~4.8 generator.
The detector form
\path{Theorem49BoundaryRootSynthesis.exists_projectedGenerator_chainDot_ne_zero_of_ne_zero_kirchhoffQuotientClass}
produces a projected generator with nonzero pairing for every nonzero purified
Kirchhoff quotient class.  Thus nondegeneracy is now available directly on the
quotient/error surface, not only after passing to the concrete target.
The combined downstream form is
\path{Theorem49BoundaryRootSynthesis.exists_rawGeneratorSum_and_projectedGenerator_chainDot_ne_zero_of_ne_zero_kirchhoffQuotientClass}:
each nonzero purified Kirchhoff quotient class admits a finite raw
Definition~4.8 generator-sum representative, and that representative is
detected by a projected Definition~4.8 generator after selected-boundary
erasure.  This packages representative choice, quotient nonzeroness transport,
and projected-generator detection into one theorem at the synthesis surface.
The same nondegeneracy is now packaged as a quotient-level linear map:
\path{Theorem49BoundaryRootSynthesis.kirchhoffQuotientProjectedGeneratorPairingMap}
sends a purified Kirchhoff quotient class to its chain-dot functional on the
projected Definition~4.8 generator subspace after selected-boundary erasure.
The computation theorem
\path{Theorem49BoundaryRootSynthesis.kirchhoffQuotientProjectedGeneratorPairingMap_apply_mk}
evaluates this map on quotient classes represented by purified Kirchhoff chains,
and
\path{Theorem49BoundaryRootSynthesis.ker_kirchhoffQuotientProjectedGeneratorPairingMap_eq_bot}
proves that its kernel is trivial.  This turns the quotient/error separation
endpoint into a standard linear-map kernel statement.  The same map-level
surface now gives
\path{Theorem49BoundaryRootSynthesis.kirchhoffQuotientProjectedGeneratorPairingMap_eq_zero_iff}
and
\path{Theorem49BoundaryRootSynthesis.exists_projectedGenerator_pairingMap_ne_zero_of_ne_zero_kirchhoffQuotientClass}:
an arbitrary purified Kirchhoff quotient class is zero exactly when its
projected-generator functional is zero, and every nonzero quotient class is
detected by evaluating that functional at some projected generator.  The
finite-dimensional consequences are also available at this same surface:
\path{Theorem49BoundaryRootSynthesis.finrank_kirchhoffQuotient_le_projectedGeneratorSubspace}
derives a dimension bound for the purified Kirchhoff quotient directly from
this trivial-kernel pairing map and the dual-rank identity for finite
subspaces, while
\path{Theorem49BoundaryRootSynthesis.finrank_range_kirchhoffQuotientProjectedGeneratorPairingMap}
identifies the dimension of the pairing-map image with the dimension of the
purified Kirchhoff quotient.  The image is now an operational quotient model:
\path{Theorem49BoundaryRootSynthesis.kirchhoffQuotientEquivPairingMapRange}
is the linear equivalence from the purified Kirchhoff quotient onto the
pairing-map image, and
\path{Theorem49BoundaryRootSynthesis.kirchhoffQuotientEquivPairingMapRange_apply_mk}
computes that image functional on representatives by selected-boundary erasure
and the chain-dot pairing.  The raw-source quotient now lands on the same image
model directly:
\path{Theorem49BoundaryRootSynthesis.rawSourceQuotientEquivPairingMapRange}
identifies the Kirchhoff-valid raw Definition~4.8 source quotient with the
pairing-map image, and
\path{Theorem49BoundaryRootSynthesis.rawSourceQuotientEquivPairingMapRange_apply_mk}
computes its representative functional by the same selected-boundary erasure
and chain-dot pairing.  The corresponding separation surface is also explicit:
\path{Theorem49BoundaryRootSynthesis.rawSourceQuotientClass_eq_zero_iff_projectedGenerator_orthogonal}
tests raw-source quotient classes by vanishing against all projected generators,
while
\path{Theorem49BoundaryRootSynthesis.exists_projectedGenerator_chainDot_ne_zero_of_ne_zero_rawSourceQuotientClass}
turns every nonzero raw-source quotient class into a nonzero projected-generator
pairing witness.  The finite representative form
\path{Theorem49BoundaryRootSynthesis.exists_rawGeneratorSum_and_projectedGenerator_chainDot_ne_zero_of_ne_zero_rawSourceQuotientClass}
now gives the same conclusion with an explicit finite raw Definition~4.8
generator-family representative in the original raw-source quotient class.
Thus a downstream nonzero raw-source quotient class can be replaced by a finite
raw generator sum without losing the projected-generator detector.  The
concrete target-level version is now explicit as well:
\path{Theorem49BoundaryRootSynthesis.exists_rawGeneratorSum_and_projectedGenerator_chainDot_ne_zero_of_ne_zero_mem_theorem49BoundaryZeroKirchhoffSubspace}
starts from a nonzero element of the boundary-zero Kirchhoff target, chooses a
finite raw Definition~4.8 generator-family sum whose selected-boundary erasure
is exactly that target, and simultaneously returns a projected generator with
nonzero chain-dot pairing against the target.  The sharpened non-vacuity form
\path{Theorem49BoundaryRootSynthesis.exists_nonempty_rawGeneratorSum_and_nonzero_projectedGenerator_chainDot_ne_zero_of_ne_zero_mem_theorem49BoundaryZeroKirchhoffSubspace}
adds that the finite raw generator-family sum is genuinely nonempty and that
the projected detector is a nonzero chain.  Thus a nonzero theorem-4.9 target
cannot be witnessed by empty raw data or by a zero projected detector.  The
further support-sensitive form
\path{Theorem49BoundaryRootSynthesis.exists_nonzero_rawGeneratorSum_and_nonempty_coeffSupport_and_nonzero_projectedGenerator_chainDot_ne_zero_of_ne_zero_mem_theorem49BoundaryZeroKirchhoffSubspace}
shows that the raw sum itself is nonzero and that the coefficient support is
nonempty, not merely that the chosen ambient finite term set is nonempty.  The
term-level extraction
\path{Theorem49BoundaryRootSynthesis.exists_rawGenerator_with_nonzero_coeff_and_nonzero_projectedGenerator_chainDot_ne_zero_of_ne_zero_mem_theorem49BoundaryZeroKirchhoffSubspace}
then chooses an actual raw Definition~4.8 generator-family chain from that
support with nonzero coefficient, while preserving the same nonzero raw sum,
selected-boundary erasure, and projected detector.  The strengthened
term-level cleanup
\path{Theorem49BoundaryRootSynthesis.exists_nonzero_rawGenerator_with_nonzero_coeff_and_nonzero_projectedGenerator_chainDot_ne_zero_of_ne_zero_mem_theorem49BoundaryZeroKirchhoffSubspace}
further proves that the extracted support chain satisfies \(y \ne 0\).  Thus
the witness is a nonzero raw Definition~4.8 generator-family chain with a
nonzero coefficient, not merely a formal support element attached to the same
raw representative.  The boundary-facing refinement
\path{Theorem49BoundaryRootSynthesis.exists_rawGenerator_with_nonzero_boundaryZeroProjection_and_nonzero_coeff_and_nonzero_projectedGenerator_chainDot_ne_zero_of_ne_zero_mem_theorem49BoundaryZeroKirchhoffSubspace}
then chooses such a support chain with nonzero selected-boundary erasure
\(\pi_{\partial}(y) \ne 0\).  Hence the nonzero theorem-4.9 target is detected
through a raw generator-family term that actually survives the boundary-zero
projection, rather than only through cancellation among projected-zero support
terms.  The
finite-dimensional consequences are now attached to this
raw-source image model too:
\path{Theorem49BoundaryRootSynthesis.finrank_rawSourceQuotient_eq_pairingMapRange}
identifies the raw-source quotient dimension exactly with the pairing-image
dimension, and
\path{Theorem49BoundaryRootSynthesis.finrank_rawSourceQuotient_le_projectedGeneratorSubspace}
then bounds that quotient dimension by the projected Definition~4.8 generator
subspace through the pairing-image/purified-quotient rank calculation.
\path{PlanarBoundaryClosedWalkSourceProjection.lean} now lowers this separation
interface to the route-facing source shells as well.  Canonical witness choice,
the local at-most-one source criterion, and same-boundary one-collar geometry,
for both honest closed-walk sources and successor-cycle boundary-order sources,
now yield
\[
  W_0(H) \cap (\mathrm{projectedKempeClosureGeneratorSubspace})^\perp = 0
\]
through both height-ordered and finite collar-layer witness data.  Thus the
same corrected projected-source interface carries not only projected spanning
but also the nondegenerate chain-dot separation endpoint needed by the
downstream theorem-4.9 route.  The same source-level lowering now also exposes
the pointwise detector and orthogonality-characterization forms: every nonzero
concrete boundary-zero Kirchhoff chain is detected by some projected
Definition~4.8 generator, and vanishing against all projected generators is
equivalent to the chain itself being zero.  The same source packages now also
expose the raw Definition~4.8 preimage and finite raw-generator sum forms:
every concrete boundary-zero Kirchhoff chain is the selected-boundary
projection of a raw generator-span element, and can be presented as the
projection of a finite linear combination of actual raw Definition~4.8
generators.  The constructive non-vacuous source lane is now packaged at this
same interface by
\path{Theorem49BoundaryRootNonemptyProjectedSynthesis}: on one fixed embedding
it records both \path{Theorem49BoundaryRootNonemptySynthesis} and the corrected
projected Definition~4.8 subspace and family-span equalities.  The package is
constructed from height-ordered or finite collar-layer witness data with a
nonempty selected-boundary-purified carrier, and the closed-walk and
successor-cycle source wrappers now expose it for canonical witness choice and
same-boundary one-collar geometry.  No wrapper is added for the obstructed
at-most-one/nonempty-carrier branch.
\path{Theorem49PositiveGeometricSource.lean} records this candidate
construction burden as explicit geometric source packages:
\path{Theorem49ClosedWalkCanonicalPositiveGeometry},
\path{Theorem49ClosedWalkOneCollarPositiveGeometry},
\path{Theorem49SuccessorCycleCanonicalPositiveGeometry}, and
\path{Theorem49SuccessorCycleOneCollarPositiveGeometry}.  Each package carries
the source data, the canonical witness choice or same-boundary one-collar
geometry, and the nonempty purified carrier on one embedding, and each lowers
to \path{Theorem49BoundaryRootNonemptyProjectedSynthesis} for any supplied Tait
coloring.  The fixed-embedding predicates also have direct same-embedding
lowering lemmas, so a candidate construction can target a named embedding
without first passing through a graph-level existential wrapper.
\path{Theorem49PositiveGeometricSourceRegression.lean} now calibrates this
target on the wheel-with-inner-triangle benchmark.  The model still has the
useful pair of facts, \path{wheelWithInnerTriangleTaitEdgeColoring_isTait} and
\path{selectedBoundaryInteriorEdgeEndpointVertices_nonempty_wheelWithInnerTriangle},
but the fixed embedding fails all four package predicates:
\path{not_theorem49ClosedWalkCanonicalPositiveGeometryOn_wheelWithInnerTriangle},
\path{not_theorem49ClosedWalkOneCollarPositiveGeometryOn_wheelWithInnerTriangle},
\path{not_theorem49SuccessorCycleCanonicalPositiveGeometryOn_wheelWithInnerTriangle},
and
\path{not_theorem49SuccessorCycleOneCollarPositiveGeometryOn_wheelWithInnerTriangle}.
The summary theorem
\path{wheelWithInnerTriangle_tait_nonempty_carrier_without_positiveGeometricSourcePackageOn}
therefore isolates the missing positive ingredient as genuine source-tied
canonical-choice or one-collar geometry, not coloring or carrier nonemptiness.
The same fixed benchmark now also calibrates the source-IDBRAD bridge:
\path{wheelWithInnerTriangle_closedWalkSource_tait_hasUnblockedEndpoint_without_interiorDualBoundaryRootAdjDistancePeelData}
bundles an honest closed-walk source, the explicit Tait coloring, and a local
unblocked endpoint, while proving that generic
\path{InteriorDualBoundaryRootAdjDistancePeelData} is impossible on that
embedding.
\path{Theorem49PositiveGeometricSourceImpossibility.lean} then gives the
structural correction: the four current package surfaces are not merely
uninhabited by that benchmark, but inconsistent with the current purified
nonempty carrier.  The theorem
\path{facewiseAtMostOneInteriorEdge_of_planarBoundaryCanonicalWitnessChoice}
shows that canonical witness choice forces facewise at-most-one interior edge;
\path{selectedBoundaryInteriorEdgeEndpointVertices_eq_empty_of_closedWalkAnnulusBoundarySource_and_canonicalWitnessChoice}
combines this with the closed-walk at-most-one obstruction to empty the
carrier.  Since same-boundary one-collar geometry induces the corresponding
canonical witness choice, the summary theorem
\path{not_nonempty_currentTheorem49PositiveGeometricSourcePackages} refutes all
four graph-level package structures.  The next constructive target must
therefore replace this package surface rather than instantiate it.  The same
file now also states this in the live endpoint vocabulary:
\path{not_hasUnblockedInteriorEndpoint_of_closedWalkAnnulusBoundarySource_and_canonicalWitnessChoice}
and
\path{not_hasUnblockedInteriorEndpoint_of_closedWalkAnnulusBoundarySource_and_oneCollarGeometry_with_sourceBoundaryData}
say directly that the old canonical-choice and source-preserving one-collar
branches have no local unblocked endpoint.  Equivalently,
\path{not_nonempty_planarBoundaryCanonicalWitnessChoice_of_closedWalkAnnulusBoundarySource_and_hasUnblockedInteriorEndpoint}
rules out canonical witness choice from an honest source plus
\path{HasUnblockedInteriorEndpoint}, and
\path{not_exists_oneCollarAnnulusCollarGeometry_with_sourceBoundaryData_of_closedWalkAnnulusBoundarySource_and_hasUnblockedInteriorEndpoint}
rules out source-preserving one-collar geometry on the same boundary split.
\path{PlanarBoundaryClosedWalkSourceCanonicalWitness.lean} now records the
exact source-side criterion.  The theorem
\path{nonempty_planarBoundaryCanonicalWitnessChoice_iff_facewiseAtMostOneInteriorEdge_of_closedWalkAnnulusBoundarySource}
says that, for an honest closed-walk annulus source, the source boundary split
carries a canonical witness choice exactly when each face has at most one
interior boundary edge.  Its existential graph-level companion
\path{exists_closedWalkAnnulusBoundarySource_and_planarBoundaryCanonicalWitnessChoice_iff_exists_closedWalkAnnulusBoundarySource_and_facewiseAtMostOneInteriorEdge}
is the corresponding source-search form.  Thus the canonical-witness source
attack has been reduced to the facewise at-most-one obligation itself; it is not
a hidden selector or annulus-construction compatibility condition.
\path{PlanarBoundaryClosedWalkSourceRoute.lean} now also lowers this lane to the
exact split-annulus witness-ownership surface itself:
\path{admitsPlanarBoundaryAnnulusWitnessAssignment_of_exists_closedWalkAnnulusBoundarySource_and_canonicalWitnessChoice_direct},
\path{admitsPlanarBoundaryAnnulusWitnessAssignment_of_exists_closedWalkAnnulusBoundarySource_and_facewiseAtMostOneInteriorEdge_direct},
\path{admitsPlanarBoundaryAnnulusWitnessAssignment_of_exists_boundaryReachabilityData_and_dartSuccessorCycleEmbeddingData_and_selectedBoundaryArc_and_canonicalWitnessChoice_direct},
and
\path{admitsPlanarBoundaryAnnulusWitnessAssignment_of_exists_boundaryReachabilityData_and_dartSuccessorCycleEmbeddingData_and_selectedBoundaryArc_and_facewiseAtMostOneInteriorEdge_direct}.
So on the canonical-witness lane, no annulus-decomposition or witness-assignment
packaging burden remains below the facewise at-most-one condition itself.
\path{Theorem49AtMostOneNonemptyCarrierImpossibility.lean} now instantiates
that exact live successor-cycle facewise-at-most-one shell on the explicit
diamond benchmark via
\path{exists_boundaryReachabilityData_and_dartSuccessorCycleEmbeddingData_and_selectedBoundaryArc_and_facewiseAtMostOneInteriorEdge_diamondWithTriangleGraph},
\path{diamondWithTriangleGraph_admitsPlanarBoundaryAnnulusWitnessAssignment_via_successorCycle_facewiseAtMostOneInteriorEdge},
and
\path{exists_theorem49BoundaryRootSynthesis_diamondWithTriangleGraph_for_explicitTaitColoring_via_successorCycle_facewiseAtMostOneInteriorEdge}.
So the live successor-cycle at-most-one branch is genuinely inhabited and
already reaches witness assignment and theorem-4.9 synthesis on a concrete
graph; the remaining burden on that branch is the nonempty-carrier upgrade, not
basic route reachability.
\path{Theorem49AtMostOneNonemptyCarrierImpossibility.lean} now also packages
the stronger fallback/count/boundary-rest hypothesis surface on that benchmark
via
\path{exists_boundaryReachabilityData_and_dartSuccessorCycleEmbeddingData_and_selectedBoundaryArc_and_atMostOneInteriorEdgePerFace_diamondWithTriangleGraph},
and then instantiates the first corrected projected Definition~4.8 endpoints as
\path{exists_projectedKempeClosureGeneratorSubspace_eq_planarBoundaryZeroSubmodule_diamondWithTriangleGraph_for_explicitTaitColoring_via_successorCycle_atMostOneInteriorEdgePerFace}
and
\path{exists_span_projectedKempeClosureGeneratorFamily_eq_planarBoundaryZeroSubmodule_diamondWithTriangleGraph_for_explicitTaitColoring_via_successorCycle_atMostOneInteriorEdgePerFace}.
So this same live successor-cycle at-most-one shell already reaches the
corrected projected spanning conclusions on an explicit graph as well; the
remaining burden on that branch is the nonempty-carrier positive upgrade, not
theorem-4.9 or projected-spanning reachability.
The same benchmark file now pushes one algebraic layer deeper still:
\path{exists_theorem49BoundaryZeroKirchhoffSubspace_inf_chainDot_orthogonal_projectedKempeClosureGeneratorSubspace_eq_bot_diamondWithTriangleGraph_for_explicitTaitColoring_via_successorCycle_atMostOneInteriorEdgePerFace}
and
\path{exists_theorem49BoundaryZeroKirchhoffSubspace_chainDot_projectedKempeClosureGeneratorSubspace_eq_zero_iff_diamondWithTriangleGraph_for_explicitTaitColoring_via_successorCycle_atMostOneInteriorEdgePerFace}
show that the explicit live shell also reaches the first chain-dot duality
endpoints on the natural purified carrier vertex set.  So the remaining burden
on that constructive branch is not algebraic projection/duality reachability
either; it is the positive nonempty-carrier upgrade itself.
The same file now also reaches the raw finite corrected Definition~4.8
representation there:
\path{exists_kempeClosureGeneratorFamily_finset_sum_boundaryZeroProjection_of_mem_theorem49BoundaryZeroKirchhoffSubspace_diamondWithTriangleGraph_for_explicitTaitColoring_via_successorCycle_atMostOneInteriorEdgePerFace}
packages finite Kempe-closure generator sums realizing every class in the
corresponding boundary-zero Kirchhoff subspace on that same explicit live shell.
So the remaining burden on this constructive branch is not raw
Definition~4.8 realization either; it is still the positive nonempty-carrier
upgrade itself.  The same benchmark file now also reaches the
projected-generator detector there:
\path{exists_projectedKempeClosureGeneratorSubspace_chainDot_ne_zero_of_ne_zero_mem_theorem49BoundaryZeroKirchhoffSubspace_diamondWithTriangleGraph_for_explicitTaitColoring_via_successorCycle_atMostOneInteriorEdgePerFace}
shows that every nonzero class in that same boundary-zero Kirchhoff subspace is
detected by some projected Kempe generator on the explicit live shell.  So the
remaining burden on this constructive branch is not nontrivial
projected-generator detection either; it is still the positive
nonempty-carrier upgrade itself.
\path{Theorem49PositiveGeometricSourceReplacement.lean} records that
replacement explicitly.  The fixed-embedding predicates
\path{Theorem49HeightOrderedPositiveProjectedGeometryOn} and
\path{Theorem49CollarLayerPositiveProjectedGeometryOn} require only the
corresponding boundary-aware witness-cover data plus nonempty purified carrier
on the same embedding, and they lower directly to
\path{Theorem49BoundaryRootNonemptyProjectedSynthesis}.  The graph-level
structures \path{Theorem49HeightOrderedPositiveProjectedGeometry} and
\path{Theorem49CollarLayerPositiveProjectedGeometry} package the same endpoint,
and Lean proves the fixed-embedding equivalence
\path{heightOrderedPositiveProjectedGeometryOn_iff_collarLayerPositiveProjectedGeometryOn}
by the existing height-to-collar and collar-to-height conversions.  The file
also curries the structural obstruction into compatibility lemmas showing that
nonempty purified carrier forbids adding the old closed-walk canonical-choice
or source-preserving one-collar fields on the same embedding.  Thus the next
positive task is construction of these direct witness-data packages from
boundary-order geometry, not another wrapper around the refuted surfaces.  The
same file now factors the construction burden into graph-level bridge theorems:
\path{theorem49CollarLayerPositiveProjectedGeometry_of_annulusCollarGeometry}
and
\path{theorem49HeightOrderedPositiveProjectedGeometry_of_annulusCollarGeometry}
package annulus collar geometry plus a surviving purified carrier directly as
the replacement graph packages, and the corresponding existential theorems do
the same from an arbitrary witness of those hypotheses.  Repaired
previous-boundary witness geometry has matching specializations through
\path{theorem49CollarLayerPositiveProjectedGeometry_of_annulusPreviousBoundaryWitnessGeometry}
and
\path{theorem49HeightOrderedPositiveProjectedGeometry_of_annulusPreviousBoundaryWitnessGeometry}.
The wheel-with-inner-triangle regression now reaches these direct predicates
themselves:
\path{not_theorem49HeightOrderedPositiveProjectedGeometryOn_wheelWithInnerTriangle}
and
\path{not_theorem49CollarLayerPositiveProjectedGeometryOn_wheelWithInnerTriangle}
show that the strict three-spoke height cycle blocks the fixed-embedding
replacement packages despite the honest source, Tait coloring, and purified
carrier.  The source-level failed universal
\path{not_forall_some_replacementPositiveProjectedGeometry_or_previousBoundaryWitness_of_closedWalkAnnulusBoundarySource_and_taitEdgeColoring_and_nonempty_selectedBoundaryInteriorCarrier_wheelWithInnerTriangle}
also rules out deriving either direct replacement package or repaired
previous-boundary witness geometry from those hypotheses alone.
The acyclic parent-witness surface is now connected at the same level:
\path{theorem49HeightOrderedPositiveProjectedGeometryOn_of_wellFoundedFacePeelWitnessData_and_hasUnblockedInteriorEndpoint}
uses the canonical parent-height rank to turn
\path{PlanarBoundaryWellFoundedFacePeelWitnessData} plus
\path{HasUnblockedInteriorEndpoint} into the height-ordered replacement
predicate, and
\path{theorem49CollarLayerPositiveProjectedGeometryOn_of_wellFoundedFacePeelWitnessData_and_hasUnblockedInteriorEndpoint}
gives the finite-collar form by grouping those heights.  The projected endpoint
is exposed directly as
\path{theorem49BoundaryRootNonemptyProjectedSynthesis_of_wellFoundedFacePeelWitnessData_and_hasUnblockedInteriorEndpoint},
with graph-level nonempty and existential versions for the same hypotheses.
The repaired previous-boundary lane now also consumes the named local endpoint
witness directly:
\path{theorem49BoundaryRootNonemptyProjectedSynthesis_of_annulusPreviousBoundaryWitnessGeometry_and_hasUnblockedInteriorEndpoint}
reaches the projected endpoint from previous-boundary witness geometry plus
\path{HasUnblockedInteriorEndpoint}, and the fixed predicates have matching
finite-collar and height-ordered constructors.  The closed-walk route file adds
the source-facing traceability theorem
\path{theorem49BoundaryRootNonemptyProjectedSynthesis_of_closedWalkAnnulusBoundarySource_and_annulusPreviousBoundaryWitnessGeometry_and_hasUnblockedInteriorEndpoint}.
It also adds the acyclic source-facing theorem
\path{theorem49BoundaryRootNonemptyProjectedSynthesis_of_closedWalkAnnulusBoundarySource_and_wellFoundedFacePeelWitnessData_and_hasUnblockedInteriorEndpoint}.
The same acyclic input is now produced directly from the root-distance
interior-dual route: \path{PlanarBoundaryClosedWalkSourceRoute.lean} defines
\path{planarBoundaryWellFoundedFacePeelWitnessData_of_closedWalkAnnulusBoundarySource_and_interiorDualBoundaryRootAdjDistancePeelData},
then
\path{planarBoundaryHeightOrderedFacePeelWitnessData_of_closedWalkAnnulusBoundarySource_and_interiorDualBoundaryRootAdjDistancePeelData}
as the source-facing height-ordered witness-cover object itself, with the
graph-level theorem
\path{admitsPlanarBoundaryHeightOrderedFacePeelWitnessData_of_exists_closedWalkAnnulusBoundarySource_and_interiorDualBoundaryRootAdjDistancePeelData_direct},
and \path{Theorem49PositiveGeometricSourceReplacementRoute.lean} proves
\path{theorem49BoundaryRootNonemptyProjectedSynthesis_of_closedWalkAnnulusBoundarySource_and_interiorDualBoundaryRootAdjDistancePeelData_and_hasUnblockedInteriorEndpoint}.
The successor-cycle source fields have the parallel endpoint
\path{theorem49BoundaryRootNonemptyProjectedSynthesis_of_boundaryReachabilityData_and_dartSuccessorCycleEmbeddingData_and_selectedBoundaryArc_and_interiorDualBoundaryRootAdjDistancePeelData_and_hasUnblockedInteriorEndpoint}.
The same source-facing root-distance bridge now reaches the raw
quotient/error endpoint itself:
\path{theorem49BoundaryRawQuotientErrorPackage_of_closedWalkAnnulusBoundarySource_and_interiorDualBoundaryRootAdjDistancePeelData_and_hasUnblockedInteriorEndpoint}
and
\path{theorem49BoundaryRawQuotientErrorPackage_of_boundaryReachabilityData_and_dartSuccessorCycleEmbeddingData_and_selectedBoundaryArc_and_interiorDualBoundaryRootAdjDistancePeelData_and_hasUnblockedInteriorEndpoint}
return the raw Definition 4.8 representative plus selected-boundary-kernel
decomposition for every Kirchhoff chain once a Tait coloring is supplied.
Thus generic boundary-root interior-dual adjacency-distance data now supplies
the acyclic witness-cover side of the positive projected and raw quotient
routes directly.  If one starts below this layer, from raw closed-walk or successor-cycle
boundary-order source data alone, the minimal missing lemma is a same-embedding
construction of \path{InteriorDualBoundaryRootAdjDistancePeelData}, or an
equivalent cycle-breaking well-founded parent witness.  The file
\path{Theorem49PlanarBoundaryAnnulusRootInteriorDual.lean} now makes this
source-side target explicit through
\path{interiorDualBoundaryRootAdjDistancePeelDataOfClosedWalkAnnulusBoundarySourceCoverSeparatedAnnulusBoundaryRootsCanonicalRootedSharedEdgeFaceMembershipCover}:
an honest closed-walk source plus cover/separated roots, boundary-free peeled
faces, canonical rooted shared-edge coverage of all interior edges, child-face
membership with strict root-distance increase, and the two-incidence bound
constructs the generic interior-dual boundary-root package on the same
embedding.  Its graph-level companion
\path{admitsPlanarBoundaryInteriorDualBoundaryRootAdjDistancePeelData_of_exists_closedWalkAnnulusBoundarySource_and_coverSeparatedAnnulusBoundaryRootsCanonicalRootedSharedEdgeFaceMembershipCover}
states the same obligation as an admissibility theorem.  Thus the source-side
attack has been localized to producing these rooted-peel fields from the
boundary-order source, rather than to a selector-internal compatibility claim
or an unnamed construction shell.  A sharper source-root constructor,
\path{interiorDualBoundaryRootAdjDistancePeelDataOfClosedWalkAnnulusBoundarySourceBoundaryFaceRootsCanonicalRootedSharedEdgeFaceMembershipCover},
now fixes the root set to the outer-or-inner boundary faces extracted from the
closed-walk source, fixes the fallback edge to the source's honest facial
closed-walk fallback, and uses the embedding's built-in two-incidence bound.
Thus the remaining IDBRAD work is exactly cover/separation of the source
boundary-face roots, boundary-free peeled faces, interior-edge coverage by
rooted shared edges, and strict root-distance child membership.  The Lean
development now also exposes the parent-oriented sufficiency route for the
same fixed roots:
\path{PlanarBoundaryClosedWalkAnnulusBoundarySource.wellFounded_parentRel_boundaryFaceRoots_parentTowardsRoot}
proves that cover/separated source boundary-face roots induce a well-founded
canonical shortest-path parent orientation, while
\path{interiorDualBoundaryRootParentPeelDataOfClosedWalkAnnulusBoundarySourceBoundaryFaceRootsCanonicalParentSharedEdgeCover}
packages the same source roots and fallback edge into
\path{InteriorDualBoundaryRootParentPeelData} once rooted shared-edge coverage
is supplied.  The canonical parent child-membership field is now derived inside
Lean from that cover and the ambient two-incidence bound by
\path{canonicalParentChildMembership_of_canonicalParentCover}; the earlier
face-membership-cover constructor remains available for callers that already
have the field.  The graph-level endpoint
\path{admitsPlanarBoundaryInteriorDualBoundaryRootParentPeelData_of_exists_closedWalkAnnulusBoundarySourceBoundaryFaceRootsCanonicalParentSharedEdgeCover}
records this raw-cover fixed-source acyclic-parent sufficiency theorem; the
older graph-level
\path{admitsPlanarBoundaryInteriorDualBoundaryRootParentPeelData_of_exists_closedWalkAnnulusBoundarySourceBoundaryFaceRootsCanonicalParentSharedEdgeFaceMembershipCover}
remains as the compatibility theorem for callers that already supply the
derived child-membership field.  The actual successor-cycle boundary-order
shell now has the matching lowering package
\path{admitsPlanarBoundaryInteriorDualBoundaryRootParentPeelData_of_exists_boundaryReachabilityData_and_dartSuccessorCycleEmbeddingData_and_selectedBoundaryArc_and_boundaryFaceRootsCanonicalParentSharedEdgeFaceMembershipCover},
\path{admitsPlanarBoundaryInteriorDualBoundaryRootParentPeelData_of_exists_boundaryReachabilityData_and_dartSuccessorCycleEmbeddingData_and_selectedBoundaryArc_and_boundaryFaceRootsCanonicalParentSharedEdgeCover},
and the direct graph-level endpoints
\path{exists_theorem49BoundaryRootSynthesis_of_exists_boundaryReachabilityData_and_dartSuccessorCycleEmbeddingData_and_selectedBoundaryArc_and_boundaryFaceRootsCanonicalParentSharedEdgeFaceMembershipCover_direct}
and
\path{exists_theorem49BoundaryRootSynthesis_of_exists_boundaryReachabilityData_and_dartSuccessorCycleEmbeddingData_and_selectedBoundaryArc_and_boundaryFaceRootsCanonicalParentSharedEdgeCover_direct}.
All four lower through
\path{PlanarBoundaryClosedWalkAnnulusBoundarySource.ofDartSuccessorCycleFields},
so once the live successor-cycle shell carries this concrete source-fixed
canonical-parent cover package, no further projection step remains before full
Theorem~4.9 synthesis.
There is now also a proved base case for this sufficiency branch:
\path{interiorDualBoundaryRootParentPeelDataOfClosedWalkAnnulusBoundarySourceBoundaryFaceRootsNoInteriorEdges}
and
\path{admitsPlanarBoundaryInteriorDualBoundaryRootParentPeelData_of_exists_closedWalkAnnulusBoundarySourceBoundaryFaceRootsNoInteriorEdges}
construct the same source-fixed parent package with empty peel set when the
source boundary-face roots cover/separate and the interior-edge support is
empty.  Thus the nondegenerate parent-source task is specifically to produce
or supply the pairwise-unique shared-edge selector input, cover/separated
source roots, boundary-free peeled faces, and a rooted shared-edge cover of
actual interior edges.  It is no longer a separate canonical parent
child-membership problem.  Lean also derives the otherwise-required
pairwise shared-edge uniqueness in this degenerate case via
\path{pairwiseUniqueSharedInteriorEdges_of_interiorEdgeSupport_eq_empty}, so no
extra uniqueness premise remains hidden in the base case.
This raw parent-cover contract is now known to be degenerate at the honest
closed-walk source layer.  The walk source theorem
\path{PlaneEmbeddingWithBoundary.FaceBoundaryClosedWalk.three_le_faceBoundary_card}
shows that every honest nonempty simple facial closed walk has at least three
boundary edges, while the height-terminal lemma for
\path{PlanarBoundaryHeightOrderedFacePeelWitnessData} forces any nonempty
raw-cover parent package to contain a terminal peeled face whose non-witness
boundary edges are already selected-boundary edges.  Since the raw contract
also requires peeled faces to be disjoint from selected-boundary support, that
terminal face would have boundary cardinality at most one.  Lean packages the
contradiction as
\path{false_of_closedWalkAnnulusBoundarySourceBoundaryFaceRootsCanonicalParentSharedEdgeCover_of_interiorEdgeSupport_nonempty}
and the route-facing degeneracy statement as
\path{interiorEdgeSupport_eq_empty_of_closedWalkAnnulusBoundarySourceBoundaryFaceRootsCanonicalParentSharedEdgeCover}.
Thus the current honest closed-walk source-fixed raw-cover contract has only
the no-interior-edge branch; any future nondegenerate sufficiency theorem must
change the source-side contract or replace this terminal parent condition.
This degeneracy is now also enough to close the downstream theorem-4.9 route on
that shell itself.  \path{Theorem49ResidualBoundaryPeeling.lean} proves
\path{theorem49BoundaryRootSynthesis_of_closedWalkAnnulusBoundarySourceBoundaryFaceRootsCanonicalParentSharedEdgeCover_via_residualBoundary}
and
\path{theorem49BoundaryRootSynthesis_of_boundaryReachabilityData_and_dartSuccessorCycleEmbeddingData_and_selectedBoundaryArc_and_boundaryFaceRootsCanonicalParentSharedEdgeCover_via_residualBoundary}.
So the source-fixed raw canonical-parent shared-edge-cover shell, and its
successor-cycle boundary-order form, already reach the full Theorem~4.9
synthesis endpoint on the same embedding before any additional endpoint witness
is added.
The same cardinality obstruction is now stated at the IDBRAD layer itself:
\path{InteriorDualBoundaryRootAdjDistancePeelData.false_of_mem_interiorEdgeSupport_of_faceBoundaryClosedWalkGeometry}
combines honest facial closed walks with the terminal-leaf theorem for
\path{InteriorDualBoundaryRootAdjDistancePeelData}.  A covered interior edge
would force a peeled face of boundary-cardinality at most one, contradicting
the three-edge lower bound for every honest nonempty simple facial closed
walk.  The source-facing theorem
\path{not_nonempty_interiorDualBoundaryRootAdjDistancePeelData_of_interiorEdgeSupport_nonempty_of_faceBoundaryClosedWalkGeometry}
therefore rules out any current boundary-free adjacent-distance/root-parent
package on an honest closed-walk embedding with genuine interior support.
This identifies the interface itself, not only the source-fixed raw-cover
constructor, as the part that must be revised for a nondegenerate route.
The nonempty branch now has a constructor-side sanity theorem.  The generic
lemma
\path{exists_nonroot_peelFace_parent_rootedSharedInteriorEdge_eq_of_canonicalParentCover_of_mem_interiorEdgeSupport}
proves that an actual interior edge covered by the canonical rooted shared-edge
image is covered by a peeled non-root face with a real shortest-path parent,
and that the rooted shared edge on that face is exactly the original interior
edge.  Its nonempty-support corollary
\path{exists_nonroot_peelFace_parent_rootedSharedInteriorEdge_eq_of_canonicalParentCover_of_interiorEdgeSupport_nonempty}
makes the nondegenerate forest branch explicit.  The source-fixed
specializations
\path{exists_nonroot_peelFace_parent_rootedSharedInteriorEdge_eq_of_closedWalkAnnulusBoundarySourceBoundaryFaceRootsCanonicalParentSharedEdgeFaceMembershipCover_of_mem_interiorEdgeSupport}
and
\path{exists_nonroot_peelFace_parent_rootedSharedInteriorEdge_eq_of_closedWalkAnnulusBoundarySourceBoundaryFaceRootsCanonicalParentSharedEdgeFaceMembershipCover_of_interiorEdgeSupport_nonempty}
instantiate this witness using the closed-walk source's boundary-face roots and
fallback edge.  Thus a nonempty canonical-parent cover cannot be witnessed by
the fallback branch.  The same constructor layer now also proves parent-side
membership of canonical shared edges and
\path{canonicalParentChildMembership_of_canonicalParentCover}: every
non-witness edge on a peeled face is either selected boundary support or lies on
a peeled child whose canonical parent is the original face.  Hence the
source-fixed parent constructor no longer takes parent child-membership as a
raw source hypothesis; the remaining source-side burden is the pairwise-unique
shared-edge selector input plus the root cover/separation, boundary-free
peeled-face, and rooted shared-edge cover hypotheses from boundary-order
geometry.
The base case is now exercised on a concrete source rather than only exposed as
an abstract constructor.  The new positive regression
\path{Theorem49PlanarBoundaryAnnulusRootInteriorDualPositiveRegression.lean}
defines the disjoint two-triangle source
\path{twoTriangleClosedWalkAnnulusBoundarySource}, proves
\path{twoTriangleAnnulusBoundaryFaceRoots_eq_univ}, proves root cover and root
separation because the interior dual is \(\bot\), and proves
\path{twoTriangleAnnulus_interiorEdgeSupport_eq_empty}.  It then constructs
\path{twoTriangleClosedWalkSourceInteriorDualBoundaryRootParentPeelData} and
the graph-level
\path{admitsPlanarBoundaryInteriorDualBoundaryRootParentPeelData_twoTriangleAnnulusGraph}.
Thus the empty-support sufficiency branch has both a generic theorem and a
concrete same-source instantiation.  Later the exact-v23 positive regression
adds an explicit Tait coloring on the same graph and proves
\path{theorem49BoundaryRootSynthesis_twoTriangleAnnulusEmbedding_for_explicitTaitColoring_via_noInteriorParentRoute},
so this no-interior parent branch is not merely inhabited: on that benchmark
it already reaches full \path{Theorem49BoundaryRootSynthesis}.  More sharply,
that same benchmark now also proves
\path{exists_embedding_taitEdgeColoring_and_theorem49BoundaryRootSynthesis_without_planarBoundaryAnnulusConstructionPositiveFrontierData_twoTriangleAnnulusGraph}
and the failed universal
\path{not_forall_nonempty_planarBoundaryAnnulusConstructionPositiveFrontierData_of_taitEdgeColoring_and_theorem49BoundaryRootSynthesis_twoTriangleAnnulusGraph},
so even the final theorem-4.9 endpoint under a real Tait coloring does not
force the positive-frontier annulus-construction shell there.  More sharply,
\path{SomePositiveStageConstructionShell},
\path{NoPositiveStageConstructionShells},
\path{not_somePositiveStageConstructionShell_of_noPositiveStageConstructionShells},
and
\path{noPositiveStageConstructionShells_of_interiorEdgeSupport_eq_empty} in
\path{PlanarBoundaryAnnulusConstructionCore.lean},
as well as
\path{twoTriangle_boundaryReachabilityData_and_dartSuccessorCycleEmbeddingData_and_selectedBoundaryArc_and_taitEdgeColoring_and_v23ResidualBoundaryInitialState_and_theorem49BoundaryRootSynthesis_and_noPositiveStageConstructionShells},
and
\path{exists_embedding_taitEdgeColoring_and_theorem49BoundaryRootSynthesis_and_noPositiveStageConstructionShells_twoTriangleAnnulusGraph}
show that on this actual successor-cycle benchmark full theorem-4.9 synthesis
already coexists with simultaneous failure of all four currently formalized
positive-stage construction shells, not only the last positive-frontier one.
The generic
lifting now lives in \path{Theorem49Synthesis.lean} as
\path{not_forall_nonempty_of_exists_embedding_taitEdgeColoring_and_theorem49BoundaryRootSynthesis_without},
\path{not_forall_somePositiveStageConstructionShell_of_exists_embedding_taitEdgeColoring_and_theorem49BoundaryRootSynthesis_and_noPositiveStageConstructionShells},
\path{not_forall_nonempty_planarBoundaryAnnulusConstructionBoundarySupportFaceData_of_exists_embedding_taitEdgeColoring_and_theorem49BoundaryRootSynthesis_without_planarBoundaryAnnulusConstructionBoundarySupportFaceData},
\path{not_forall_nonempty_planarBoundaryAnnulusConstructionFacePartitionData_of_exists_embedding_taitEdgeColoring_and_theorem49BoundaryRootSynthesis_without_planarBoundaryAnnulusConstructionFacePartitionData},
\path{not_forall_nonempty_planarBoundaryAnnulusConstructionFaceLayerData_of_exists_embedding_taitEdgeColoring_and_theorem49BoundaryRootSynthesis_without_planarBoundaryAnnulusConstructionFaceLayerData},
and
\path{not_forall_nonempty_planarBoundaryAnnulusConstructionPositiveFrontierData_of_exists_embedding_taitEdgeColoring_and_theorem49BoundaryRootSynthesis_without_planarBoundaryAnnulusConstructionPositiveFrontierData},
so any same-embedding synthesis witness of that shape already refutes a
universal derivation of the corresponding positive-stage construction package.
The exact two-triangle benchmark now instantiates that pattern not only for
positive-frontier data but also for boundary-support face data,
face-partition data, face-layer data, and the aggregated
\path{SomePositiveStageConstructionShell} surface.  It now also pushes that
same failed-converse diagnosis back up to the exact-v23 source shells
themselves:
\path{twoTriangle_closedWalkSource_tait_and_v23ResidualBoundaryInitialState_and_theorem49BoundaryRootSynthesis_and_noPositiveStageConstructionShells},
\path{exists_embedding_closedWalkAnnulusBoundarySource_and_taitEdgeColoring_and_v23ResidualBoundaryInitialState_and_theorem49BoundaryRootSynthesis_without_somePositiveStageConstructionShell_twoTriangleAnnulusGraph},
and
\path{not_forall_somePositiveStageConstructionShell_of_closedWalkAnnulusBoundarySource_and_v23ResidualBoundaryInitialState_and_taitEdgeColoring_and_theorem49BoundaryRootSynthesis_twoTriangleAnnulusGraph}
do this on the honest closed-walk shell, while
\path{exists_embedding_boundaryReachabilityData_and_dartSuccessorCycleEmbeddingData_and_selectedBoundaryArc_and_taitEdgeColoring_and_v23ResidualBoundaryInitialState_and_theorem49BoundaryRootSynthesis_without_somePositiveStageConstructionShell_twoTriangleAnnulusGraph}
and
\path{not_forall_somePositiveStageConstructionShell_of_boundaryReachabilityData_and_dartSuccessorCycleEmbeddingData_and_selectedBoundaryArc_and_v23ResidualBoundaryInitialState_and_taitEdgeColoring_and_theorem49BoundaryRootSynthesis_twoTriangleAnnulusGraph}
do the same on the live successor-cycle selected-arc shell.  The generic exact-v23
route-side converses now live in \path{Theorem49ResidualBoundaryRoute.lean} as
\path{not_forall_somePositiveStageConstructionShell_of_exists_embedding_closedWalkAnnulusBoundarySource_and_taitEdgeColoring_and_v23ResidualBoundaryInitialState_and_theorem49BoundaryRootSynthesis_and_noPositiveStageConstructionShells}
and
\path{not_forall_somePositiveStageConstructionShell_of_exists_embedding_boundaryReachabilityData_and_dartSuccessorCycleEmbeddingData_and_selectedBoundaryArc_and_taitEdgeColoring_and_v23ResidualBoundaryInitialState_and_theorem49BoundaryRootSynthesis_and_noPositiveStageConstructionShells},
and the two-triangle source-shell failures simply instantiate those
shell-specific helpers from the stronger witness package
\path{NoPositiveStageConstructionShells}.  This positive source now also
reaches the
selector-valued construction surface:
\path{boundaryFreeIncidentFaceSelector_of_interiorDualBoundaryRootParentPeelData}
extracts a \path{BoundaryFreeIncidentFaceSelector} from any parent payload
while preserving the actual parent cover:
\path{boundaryFreeIncidentFaceSelector_of_interiorDualBoundaryRootParentPeelData_faceOf_mem_peelFaces}
says each selected face lies in the parent package's \path{peelFaces}, and
\path{rootedSharedInteriorEdge_boundaryFreeIncidentFaceSelector_of_interiorDualBoundaryRootParentPeelData_faceOf_eq}
says its rooted shared edge is the queried interior edge.  At the finite
construction layer,
\path{boundaryFreeIncidentFaceSelector_of_interiorDualBoundaryRootParentPeelData_peelFaces_subset}
shows that the selector image introduces no new peeled faces beyond the parent
package.  The source-fixed empty-support specialization
\path{boundaryFreeIncidentFaceSelectorOfClosedWalkAnnulusBoundarySourceBoundaryFaceRootsNoInteriorEdges}
packages the same closed-walk hypotheses, and
\path{twoTriangleClosedWalkSourceBoundaryFreeIncidentFaceSelector} instantiates
that selector on the two-triangle source.  Thus the base case produces the
actual boundary-free selector object consumed by the replacement-construction
lane, not only an inhabited parent package.
The parent-to-selector bridge now also supplies the selector strict-descent
condition itself.  The extracted selector is injective on interior edges by
\path{boundaryFreeIncidentFaceSelector_of_interiorDualBoundaryRootParentPeelData_injective};
its inverted selected witness is the canonical rooted shared edge by
\path{boundaryFreeIncidentFaceSelector_of_interiorDualBoundaryRootParentPeelData_selectedWitnessEdge_eq};
and
\path{boundaryFreeIncidentFaceSelector_of_interiorDualBoundaryRootParentPeelData_strictDescent}
turns parent child-membership into the selector \path{hrest} condition using
root distance.  Therefore
\path{theorem49BoundaryRootNonemptyProjectedSynthesis_of_closedWalkAnnulusBoundarySourceBoundaryFaceRootsCanonicalParentSharedEdgeFaceMembershipCover_via_boundaryFreeSelector}
routes the same source-fixed parent hypotheses through the selector-to-annulus
construction surface directly.  This selector route now also has downstream
route-facing forms:
\path{theorem49BoundaryRawQuotientErrorPackage_of_closedWalkAnnulusBoundarySourceBoundaryFaceRootsCanonicalParentSharedEdgeFaceMembershipCover_via_boundaryFreeSelector}
exposes the raw Definition 4.8 quotient/error package, and
\path{exists_theorem49BoundaryRootNonemptyProjectedSynthesis_of_exists_closedWalkAnnulusBoundarySourceBoundaryFaceRootsCanonicalParentSharedEdgeFaceMembershipCover_and_hasUnblockedInteriorEndpoint_via_boundaryFreeSelector}
packages the graph-level source hypotheses with the local
\path{HasUnblockedInteriorEndpoint} witness directly.  But under the current
honest-source semantics that stronger endpoint-consuming shell is also
vacuous:
\path{not_hasUnblockedInteriorEndpoint_of_closedWalkAnnulusBoundarySourceBoundaryFaceRootsCanonicalParentSharedEdgeFaceMembershipCover},
\path{not_hasUnblockedInteriorEndpoint_of_boundaryReachabilityData_and_dartSuccessorCycleEmbeddingData_and_selectedBoundaryArc_and_boundaryFaceRootsCanonicalParentSharedEdgeFaceMembershipCover},
and the matching graph-level no-coexistence theorems
\path{not_exists_embedding_closedWalkAnnulusBoundarySourceBoundaryFaceRootsCanonicalParentSharedEdgeFaceMembershipCover_and_hasUnblockedInteriorEndpoint}
and
\path{not_exists_embedding_boundaryReachabilityData_and_dartSuccessorCycleEmbeddingData_and_selectedBoundaryArc_and_boundaryFaceRootsCanonicalParentSharedEdgeFaceMembershipCover_and_hasUnblockedInteriorEndpoint}
show that the stricter child-membership shell cannot actually coexist with the
named local endpoint witness on the same honest source surfaces.  So these
selector/projected face-membership wrappers are presently route-diagnostic, not
inhabited.  More sharply, the generic converse theorems
\path{not_forall_some_closedWalkAnnulusBoundarySourceBoundaryFaceRootsCanonicalParentSharedEdgeFaceMembershipCover_of_exists_embedding_closedWalkAnnulusBoundarySource_and_hasUnblockedInteriorEndpoint}
and
\path{not_forall_some_successorCycleBoundaryFaceRootsCanonicalParentSharedEdgeFaceMembershipCover_of_exists_embedding_boundaryReachabilityData_and_dartSuccessorCycleEmbeddingData_and_selectedBoundaryArc_and_hasUnblockedInteriorEndpoint},
together with the seed-free shared-interior-pair instances
\path{not_forall_some_closedWalkAnnulusBoundarySourceBoundaryFaceRootsCanonicalParentSharedEdgeFaceMembershipCover_of_closedWalkAnnulusBoundarySource_and_hasUnblockedInteriorEndpoint_sharedInteriorPair}
and
\path{not_forall_some_successorCycleBoundaryFaceRootsCanonicalParentSharedEdgeFaceMembershipCover_of_boundaryReachabilityData_and_dartSuccessorCycleEmbeddingData_and_selectedBoundaryArc_and_hasUnblockedInteriorEndpoint_sharedInteriorPair},
show that no Tait coloring or exact v23 residual seed is needed to refute a
universal derivation of that stronger shell.  The same shared-interior-pair
benchmark now also bundles the whole local selector/canonical-parent burden at
once:
\path{not_forall_some_selectorOrCanonicalParentOrKnownSameEmbeddingGeometry_of_closedWalkAnnulusBoundarySource_and_taitEdgeColoring_and_hasUnblockedInteriorEndpoint_sharedInteriorPair}
and
\path{not_forall_some_selectorOrCanonicalParentOrKnownSameEmbeddingGeometry_of_boundaryReachabilityData_and_dartSuccessorCycleEmbeddingData_and_selectedBoundaryArc_and_taitEdgeColoring_and_hasUnblockedInteriorEndpoint_sharedInteriorPair}
show that honest source data or the live successor-cycle selected-arc shell,
even after adding a Tait coloring, still do not universally force any of the
local \path{BoundaryFreeIncidentFaceSelector}, raw canonical-parent shared-edge
cover, stricter canonical-parent face-membership cover, or the previously known
same-embedding residual, height-ordered, collar-layer, attached-certificate,
well-founded, or annulus-collar geometry surfaces from
\path{HasUnblockedInteriorEndpoint} alone.  So this endpoint-only obstruction
now blocks the whole local selector/parent ladder, not just one stronger shell
at a time.  And this is no longer only a benchmark-facing statement:
\path{not_forall_some_selectorOrCanonicalParentOrKnownSameEmbeddingGeometry_of_exists_embedding_closedWalkAnnulusBoundarySource_and_taitEdgeColoring_and_hasUnblockedInteriorEndpoint}
and
\path{not_forall_some_selectorOrCanonicalParentOrKnownSameEmbeddingGeometry_of_exists_embedding_boundaryReachabilityData_and_dartSuccessorCycleEmbeddingData_and_selectedBoundaryArc_and_taitEdgeColoring_and_hasUnblockedInteriorEndpoint}
lift that whole bundled failure pattern to reusable witness-based converses, so
any same-embedding example with honest-source or live successor-cycle data, a
Tait coloring, and \path{HasUnblockedInteriorEndpoint} but without the local
selector/canonical-parent/known-geometry package already refutes a universal
derivation of that whole burden.  The selector lane now
also closes directly at full theorem 4.9 on its own construction surface:
\path{theorem49BoundaryRootSynthesis_of_planarBoundaryAnnulusConstruction_and_hasUnblockedInteriorEndpoint}
and
\path{BoundaryFreeIncidentFaceSelector.theorem49BoundaryRootSynthesis_of_strictDescent_and_hasUnblockedInteriorEndpoint}
show that once an injective strict-descent selector and a local endpoint
witness exist on an embedding, no further projected-synthesis packaging remains
below the selector-built annulus construction.  The actual successor-cycle
boundary-order shell now has the matching raw-cover selector route too:
\path{theorem49BoundaryRootNonemptyProjectedSynthesis_of_boundaryReachabilityData_and_dartSuccessorCycleEmbeddingData_and_selectedBoundaryArc_and_boundaryFaceRootsCanonicalParentSharedEdgeCover_via_boundaryFreeSelector},
\path{theorem49BoundaryRawQuotientErrorPackage_of_boundaryReachabilityData_and_dartSuccessorCycleEmbeddingData_and_selectedBoundaryArc_and_boundaryFaceRootsCanonicalParentSharedEdgeCover_via_boundaryFreeSelector},
and
\path{exists_theorem49BoundaryRootNonemptyProjectedSynthesis_of_exists_boundaryReachabilityData_and_dartSuccessorCycleEmbeddingData_and_selectedBoundaryArc_and_boundaryFaceRootsCanonicalParentSharedEdgeCover_and_hasUnblockedInteriorEndpoint_via_boundaryFreeSelector}.
So once the live successor-cycle shell carries the raw canonical-parent
shared-edge cover together with \path{HasUnblockedInteriorEndpoint}, it already
reaches the selector-built annulus-construction lane directly, without first
detouring through residual-boundary packaging; the remaining source-facing
selector burden is only to derive that injective strict-descent selector
package itself on the live shell.  The same raw-cover selector lane now also
exposes the route-facing terminal selected-face break itself:
\path{exists_terminal_face_boundary_remainders_and_disjoint_of_interiorDualBoundaryRootParentPeelData_via_boundaryFreeSelector}
lifts the generic selector terminal-face theorem to any canonical-parent
package, and the source/raw/live-shell specializations
\path{exists_terminal_face_boundary_remainders_and_disjoint_of_closedWalkAnnulusBoundarySourceBoundaryFaceRootsCanonicalParentSharedEdgeFaceMembershipCover_via_boundaryFreeSelector},
\path{exists_terminal_face_boundary_remainders_and_disjoint_of_closedWalkAnnulusBoundarySourceBoundaryFaceRootsCanonicalParentSharedEdgeCover_via_boundaryFreeSelector}
and
\path{exists_terminal_face_boundary_remainders_and_disjoint_of_boundaryReachabilityData_and_dartSuccessorCycleEmbeddingData_and_selectedBoundaryArc_and_boundaryFaceRootsCanonicalParentSharedEdgeFaceMembershipCover_via_boundaryFreeSelector},
together with
\path{exists_terminal_face_boundary_remainders_and_disjoint_of_boundaryReachabilityData_and_dartSuccessorCycleEmbeddingData_and_selectedBoundaryArc_and_boundaryFaceRootsCanonicalParentSharedEdgeCover_via_boundaryFreeSelector}
show that both the stricter canonical-parent face-membership shell and the raw
canonical-parent cover shell, plus \path{HasUnblockedInteriorEndpoint},
already yield a peeled face whose non-witness remainders are ambient-boundary
edges while its full face boundary remains disjoint from
\path{selectedBoundarySupport}.  Thus the live selector compatibility lane now
reaches an explicit same-shell terminal boundary-break witness before any
residual-boundary or collar-layer repackaging, even before discarding the
child-membership refinement.
This parent-oriented packet now composes all the way to the nonempty projected
Theorem 4.9 endpoint:
\path{theorem49BoundaryRootNonemptyProjectedSynthesis_of_closedWalkAnnulusBoundarySourceBoundaryFaceRootsCanonicalParentSharedEdgeFaceMembershipCover_direct}
is the fixed-embedding source-to-synthesis theorem, and
\path{exists_theorem49BoundaryRootNonemptyProjectedSynthesis_of_exists_closedWalkAnnulusBoundarySourceBoundaryFaceRootsCanonicalParentSharedEdgeFaceMembershipCover_with_nonempty_selectedBoundaryInteriorEdgeEndpointVertices_direct}
is its graph-level existential form.
The wheel-with-inner-triangle regression now instantiates this source-fixed
parent bridge on a concrete graph:
\path{wheelWithInnerTriangle_theorem49BoundaryRootNonemptyProjectedSynthesis_of_closedWalkSourceBoundaryFaceRootsCanonicalParentSharedEdgeFaceMembershipCover}
specializes the theorem to the wheel embedding and its explicit Tait coloring.
The companion theorem
\path{wheelWithInnerTriangle_source_to_theorem49_and_height_cycle_contradiction_of_closedWalkSourceBoundaryFaceRootsCanonicalParentSharedEdgeFaceMembershipCover}
shows that the same source-fixed parent hypotheses lower to height-ordered
witness-cover data and contradict the known wheel height-cycle obstruction, so
\path{not_exists_wheelWithInnerTriangle_closedWalkSourceBoundaryFaceRootsCanonicalParentSharedEdgeFaceMembershipCover_and_theorem49BoundaryRootNonemptyProjectedSynthesis}
blocks this projected endpoint on the wheel.  This is the concrete soundness
check for the new parent route: the wheel's honest source, Tait coloring, and
unblocked endpoint alone do not make the parent bridge fire.
The same development now separates the first of these obligations: the theorem
\path{PlanarBoundaryClosedWalkAnnulusBoundarySource.rootSetCovers_boundaryFaceRoots_of_interiorDualBoundaryRootAdjDistancePeelData}
shows that any generic boundary-root IDBRAD datum already makes the source
boundary-face root set cover the interior-dual components.  By contrast,
\path{PlanarBoundaryClosedWalkAnnulusBoundarySource.not_rootSetSeparated_boundaryFaceRoots_of_reachable_outerBoundaryFace_innerBoundaryFace}
shows that separation of all source boundary faces is refuted by any interior-dual
reachability from an outer boundary face to an inner boundary face.  This has now
been lifted to the generic IDBRAD layer:
\path{PlanarBoundaryClosedWalkAnnulusBoundarySource.not_reachable_outerBoundaryFace_innerBoundaryFace_of_interiorDualBoundaryRootAdjDistancePeelData}
shows that any IDBRAD datum on the closed-walk source forces the extracted
outer and inner boundary-face layers to be interior-dual reachability-separated,
and
\path{PlanarBoundaryClosedWalkAnnulusBoundarySource.not_nonempty_interiorDualBoundaryRootAdjDistancePeelData_of_reachable_outerBoundaryFace_innerBoundaryFace}
records the corresponding failed-converse form.  The route-facing wrappers
\path{not_exists_closedWalkAnnulusBoundarySource_and_interiorDualBoundaryRootAdjDistancePeelData_and_reachable_outerBoundaryFace_innerBoundaryFace}
and
\path{not_forall_interiorDualBoundaryRootAdjDistancePeelData_of_exists_closedWalkAnnulusBoundarySource_and_reachable_outerBoundaryFace_innerBoundaryFace}
package the same obstruction as a same-embedding no-coexistence theorem and as
a failed universal derivation from closed-walk source data.  This obstruction
now already kills the source boundary-face separation premise itself on the
live successor-cycle shell:
\path{not_rootSetSeparated_boundaryFaceRoots_of_boundaryReachabilityData_and_dartSuccessorCycleEmbeddingData_and_selectedBoundaryArc_and_reachable_outerBoundaryFace_innerBoundaryFace}
and
\path{not_exists_boundaryReachabilityData_and_dartSuccessorCycleEmbeddingData_and_selectedBoundaryArc_and_rootSetSeparated_boundaryFaceRoots_and_reachable_outerBoundaryFace_innerBoundaryFace}
show that outer-to-inner interior-dual reachability blocks the source
\path{hsepRoots} hypothesis before any rooted shared-edge cover or generic
IDBRAD payload is added.  This obstruction
now also appears directly on the live successor-cycle shell:
\path{not_nonempty_interiorDualBoundaryRootAdjDistancePeelData_of_boundaryReachabilityData_and_dartSuccessorCycleEmbeddingData_and_selectedBoundaryArc_and_reachable_outerBoundaryFace_innerBoundaryFace}
and
\path{not_exists_boundaryReachabilityData_and_dartSuccessorCycleEmbeddingData_and_selectedBoundaryArc_and_interiorDualBoundaryRootAdjDistancePeelData_and_reachable_outerBoundaryFace_innerBoundaryFace}
show that if the source extracted from actual boundary reachability and
successor-cycle data still has an interior-dual path from one outer boundary
face to one inner boundary face, then the route already fails before the
source-fixed rooted-cover IDBRAD bridge or any theorem-4.9 packaging is used.
The concrete two-band triangular annulus benchmark now realizes this live
successor-cycle failure on an actual dart-successor shell:
\path{twoBandAnnulusDartSuccessorCycleGeometry} and
\path{twoBandAnnulusDartSuccessorCycleGeometry_selectedBoundaryArcOnFace}
instantiate the live surface,
\path{twoBandDartSuccessorClosedWalkAnnulusBoundarySource_boundaryData_eq} and
\path{twoBandDartSuccessorBoundaryFaceRoots_eq_univ} show that this live shell
extracts the same annulus boundary split and the same full six-face
boundary-root set as the closed-walk source,
\path{rootSetCovers_twoBandDartSuccessorBoundaryFaceRoots} makes the coverage
half explicit there too, and
\path{not_rootSetSeparated_twoBandAnnulusBoundaryFaceRoots_via_dartSuccessor_source_reachability}
and
\path{not_nonempty_interiorDualBoundaryRootAdjDistancePeelData_twoBandAnnulus_via_dartSuccessor_source_reachability}
show that the same outer-to-inner interior-dual adjacency already kills only
the source \path{hsepRoots} premise, and hence generic IDBRAD data, before any
rooted-cover or theorem-4.9 packaging is introduced.
The obstruction is now formalized one premise deeper on the live source-fixed
rooted-cover shell as well:
\path{false_of_boundaryReachabilityData_and_dartSuccessorCycleEmbeddingData_and_selectedBoundaryArc_and_boundaryFaceRootsCanonicalRootedSharedEdgeFaceMembershipCover_and_reachable_outerBoundaryFace_innerBoundaryFace}
and
\path{not_exists_boundaryReachabilityData_and_dartSuccessorCycleEmbeddingData_and_selectedBoundaryArc_and_boundaryFaceRootsCanonicalRootedSharedEdgeFaceMembershipCover_and_reachable_outerBoundaryFace_innerBoundaryFace}
show that if actual successor-cycle data still extract an outer boundary face
reaching an inner boundary face, then no rooted shared-edge face-membership
package can exist there at all.
The same two-band benchmark now instantiates that sharper failure directly:
\path{not_exists_twoBandAnnulus_boundaryFaceRootsCanonicalRootedSharedEdgeFaceMembershipCover_via_dartSuccessor_source_reachability}
rules out every boundary-free peeled-face set, rooted shared-edge cover, and
strict child-membership witness package on the actual successor-cycle shell.
This obstruction is now instantiated on a concrete honest source in
\path{Theorem49PlanarBoundaryAnnulusRootInteriorDualObstructionRegression.lean}.
The two-band triangular annulus
\path{twoBandAnnulusGraph} has six quadrilateral faces, disjoint selected
outer and inner boundary cycles, and the closed-walk source package
\path{twoBandClosedWalkAnnulusBoundarySource}; nevertheless
\path{twoBandAnnulus_interiorDual_adj_outer0_inner0} gives an interior-dual
adjacency between an extracted outer boundary face and an extracted inner
boundary face.  The theorem
\path{not_nonempty_interiorDualBoundaryRootAdjDistancePeelData_twoBandAnnulus_via_source_reachability}
therefore proves the concrete nonexistence of generic IDBRAD data on this
embedding.  The same concrete source now pins down the fixed-root constructor
premise: \path{twoBandAnnulusBoundaryFaceRoots_eq_univ} proves that the
source-extracted \path{boundaryFaceRoots} are all six faces,
\path{rootSetCovers_twoBandAnnulusBoundaryFaceRoots} proves the coverage
premise, and \path{not_rootSetSeparated_twoBandAnnulusBoundaryFaceRoots}
refutes the separation premise using the shared middle edge \path{tbaM01}.
The same concrete source also blocks the primitive boundary-free peel-cover
premise: \path{twoBandAnnulus_peelFaces_eq_empty_of_boundaryFree} proves that
every boundary-free peel-face set is empty, because every face touches selected
boundary support, and
\path{not_interiorEdgeSupport_subset_boundaryFreePeelImage_twoBandAnnulus}
then uses the interior edge \path{tbaM01} to refute any interior-edge-support
cover by an image of such faces.  This local obstruction is now connected to
the shared forcing-edge API by
\path{twoBandAnnulusForcingInteriorEdgeWitness},
\path{not_nonempty_boundaryFreeIncidentFaceSelector_twoBandAnnulus}, and
\path{not_nonempty_interiorDualBoundaryRootAdjDistancePeelData_twoBandAnnulus_via_forcingInteriorEdgeWitness}.
The forcing obstruction now also applies to the source-fixed parent route
itself.  Lean proves
\path{everyInteriorEdgeHasBoundaryFreeIncidentFace_of_interiorDualBoundaryRootParentPeelData},
so any canonical-parent boundary-root package supplies a boundary-free
incident face for every interior edge.  Consequently the same two-band forcing
witness proves
\path{not_nonempty_interiorDualBoundaryRootParentPeelData_twoBandAnnulus_via_forcingInteriorEdgeWitness}
and
\path{not_nonempty_planarBoundaryAnnulusRootParentPeelData_twoBandAnnulus_via_forcingInteriorEdgeWitness}.
The same forcing witness now blocks the live theorem-4.9 raw canonical-parent
cover shell one layer earlier:
\path{false_of_boundaryReachabilityData_and_dartSuccessorCycleEmbeddingData_and_selectedBoundaryArc_and_boundaryFaceRootsCanonicalParentSharedEdgeCover_and_forcingInteriorEdgeWitness}
and
\path{not_exists_embedding_boundaryReachabilityData_and_dartSuccessorCycleEmbeddingData_and_selectedBoundaryArc_and_boundaryFaceRootsCanonicalParentSharedEdgeCover_and_forcingInteriorEdgeWitness}
show that once a successor-cycle shell carries a forcing interior edge, no raw
canonical-parent shared-edge cover package can exist there at all.
The two-band benchmark now instantiates that sharper failure directly as
\path{not_exists_twoBandAnnulus_boundaryFaceRootsCanonicalParentSharedEdgeCover_via_forcingInteriorEdgeWitness}.
The same benchmark also kills the older raw carrier repair on that exact
shell:
\path{not_exists_twoBandAnnulus_boundaryFaceRootsCanonicalParentSharedEdgeCover_and_endpointSupport_disjoint_and_nonempty_interiorEdgeEndpointSupport}
shows that even endpoint-support separation plus a nonempty raw interior-edge
endpoint carrier cannot rescue the live parent-cover lane on the actual
successor-cycle benchmark.
So the remaining source-side IDBRAD burden is a genuine separation principle
plus a nontrivial boundary-free peel source, and the parent-oriented variant
has the same no-forcing obligation; neither follows from the current
closed-walk source package alone.
The generic cycle
obstruction now makes this requirement checkable: a source-level derivation must
show that each apparent remainder-witness cycle is broken by selected boundary
support, by witness identification, or by an actual well-founded orientation.
The list-level theorem
\path{false_of_isChain_height_lt_and_getLast_lt_head} supplies the generic
order obstruction used for arbitrary finite cycles.
Once that datum is
present, the remaining non-vacuity burden is the local unblocked endpoint.
The direct well-founded lane now accepts the older carrier-side hypotheses
without first repackaging them as that endpoint witness:
\path{theorem49BoundaryRootNonemptyProjectedSynthesis_of_wellFoundedFacePeelWitnessData_and_endpointSupport_disjoint_and_nonempty_interiorEdgeEndpointSupport}
uses explicit \path{PlanarBoundaryWellFoundedFacePeelWitnessData},
endpoint-support separation, and nonempty raw
\path{interiorEdgeEndpointSupport} to reach the projected synthesis endpoint,
while the named well-founded route packages and their honest-source wrappers
now also expose full theorem-4.9 through
\path{Theorem49ClosedWalkAnnulusWellFoundedPositiveProjectedGeometry.boundaryRootSynthesis},
\path{Theorem49SuccessorCycleAnnulusWellFoundedPositiveProjectedGeometry.boundaryRootSynthesis},
and
\path{theorem49BoundaryRootSynthesis_of_closedWalkAnnulusBoundarySource_and_wellFoundedFacePeelWitnessData_and_hasUnblockedInteriorEndpoint}.
The route-facing wrappers
\path{theorem49BoundaryRootSynthesis_of_closedWalkAnnulusBoundarySource_and_wellFoundedFacePeelWitnessData_and_endpointSupport_disjoint_and_nonempty_interiorEdgeEndpointSupport}
and
\path{theorem49BoundaryRootSynthesis_of_boundaryReachabilityData_and_dartSuccessorCycleEmbeddingData_and_selectedBoundaryArc_and_wellFoundedFacePeelWitnessData_and_endpointSupport_disjoint_and_nonempty_interiorEdgeEndpointSupport}
expose the same bridge on the honest closed-walk and successor-cycle source
surfaces.  Thus bypassing the blocked interior-dual derivation by supplying a
same-embedding well-founded witness-cover datum no longer leaves a separate
endpoint- or theorem-4.9-packaging obligation; the remaining hard part is still
to obtain that acyclic datum from genuine boundary-order geometry.
\path{Theorem49PositiveGeometricSourceReplacementRoute.lean} now discharges
that local endpoint witness at the same root-distance source layer from two
older geometric surfaces: global endpoint-support separation plus raw
\path{interiorEdgeEndpointSupport} nonemptiness, and peeled-face endpoint
no-touch on the root-distance annulus construction plus the same raw
nonemptiness.  The closed-walk theorem names are
\path{theorem49BoundaryRootNonemptyProjectedSynthesis_of_closedWalkAnnulusBoundarySource_and_interiorDualBoundaryRootAdjDistancePeelData_and_endpointSupport_disjoint_and_nonempty_interiorEdgeEndpointSupport}
and
\path{theorem49BoundaryRootNonemptyProjectedSynthesis_of_closedWalkAnnulusBoundarySource_and_interiorDualBoundaryRootAdjDistancePeelData_of_rootDistance_peelFaces_endpoint_disjoint_selectedBoundarySupport};
the successor-cycle source fields have parallel projected-synthesis wrappers.
The carrier burden has now been lowered one step further:
\path{interiorEdgeEndpointSupport_nonempty_of_interiorEdgeSupport_nonempty}
derives the raw endpoint carrier from a nonempty interior edge, and the
closed-walk root-distance bridge exposes corresponding projected-synthesis
variants with nonempty \path{interiorEdgeSupport} in place of pre-expanded
endpoint support.
\path{Theorem49PositiveGeometricSourceReplacementRoute.lean} then reconnects
this direct replacement surface to the honest closed-walk annulus source shell
and to the route-facing successor-cycle source shell.  The closed-walk package
\path{Theorem49ClosedWalkAnnulusCollarPositiveProjectedGeometry} carries an
honest source, same-embedding annulus collar geometry, and the nonempty
purified carrier, and it lowers to both replacement graph packages while now
also exposing graph-level full theorem-4.9 synthesis directly, so the
nonempty projected synthesis endpoint is no longer the top of that source-facing
surface.  The successor-cycle package
\path{Theorem49SuccessorCycleAnnulusCollarPositiveProjectedGeometry} carries
boundary reachability, successor-cycle facial data, selected-boundary arcs,
same-embedding annulus collar geometry, and the nonempty purified carrier.  It
lowers to the closed-walk route package for traceability and to both
replacement graph packages for the constructive route, and it inherits the same
graph-level full-synthesis closure.  The current positive route is
therefore reduced to producing this successor-cycle
annulus-geometry-plus-carrier input, or its direct annulus-collar analogue,
without passing back through the refuted canonical-choice or same-boundary
one-collar surfaces.  The route-facing file now also supplies direct
closed-walk and successor-cycle full-synthesis constructors from
\path{HasUnblockedInteriorEndpoint}, together with the matching graph-level
existential closures from either a surviving carrier or one local endpoint
witness, namely
\path{exists_theorem49BoundaryRootSynthesis_of_exists_closedWalkAnnulusBoundarySource_and_annulusCollarGeometry_and_hasUnblockedInteriorEndpoint_direct}
and
\path{exists_theorem49BoundaryRootSynthesis_of_exists_boundaryReachabilityData_and_dartSuccessorCycleEmbeddingData_and_selectedBoundaryArc_and_annulusCollarGeometry_and_hasUnblockedInteriorEndpoint_direct}.
Thus source-facing examples can provide the carrier side through a local
endpoint witness while still keeping the annulus-collar geometry explicit.
The older endpoint-separation route now feeds the same local witness without a
separate purification step:
\path{hasUnblockedInteriorEndpoint_of_endpointSupport_disjoint_and_nonempty}
turns raw \path{interiorEdgeEndpointSupport} nonemptiness plus disjointness
from the selected-boundary endpoint support into
\path{HasUnblockedInteriorEndpoint}.  The replacement and route files now
expose the corresponding direct full-synthesis constructors, including
\path{exists_theorem49BoundaryRootSynthesis_of_exists_closedWalkAnnulusBoundarySource_and_annulusCollarGeometry_and_endpointSupport_disjoint_and_nonempty_interiorEdgeEndpointSupport_direct}
and
\path{exists_theorem49BoundaryRootSynthesis_of_exists_boundaryReachabilityData_and_dartSuccessorCycleEmbeddingData_and_selectedBoundaryArc_and_annulusCollarGeometry_and_endpointSupport_disjoint_and_nonempty_interiorEdgeEndpointSupport_direct}.
The same annulus-collar source lane now has the matching full-synthesis
closures from canonical peeled-face endpoint no-touch as well.
The route-facing package layer has matching constructors as well:
\path{theorem49ClosedWalkAnnulusCollarPositiveProjectedGeometryOn_of_closedWalkAnnulusBoundarySource_and_annulusCollarGeometry_and_endpointSupport_disjoint_and_nonempty_interiorEdgeEndpointSupport}
and
\path{theorem49SuccessorCycleAnnulusCollarPositiveProjectedGeometryOn_of_boundaryReachabilityData_and_dartSuccessorCycleEmbeddingData_and_selectedBoundaryArc_and_annulusCollarGeometry_and_endpointSupport_disjoint_and_nonempty_interiorEdgeEndpointSupport}
inhabit the closed-walk and successor-cycle annulus-collar replacement
packages before any Tait coloring is supplied.
\path{Theorem49PositiveGeometricSourceReplacementRegression.lean} calibrates
that sharpened target on the shared-interior-pair annulus.  That model carries
boundary reachability, successor-cycle selected arcs, a Tait coloring, and a
nonempty purified carrier, but Lean proves
\path{not_theorem49HeightOrderedPositiveProjectedGeometryOn_sharedInteriorPair}
and
\path{not_theorem49CollarLayerPositiveProjectedGeometryOn_sharedInteriorPair}.
Consequently the theorem
\path{not_forall_replacementPositiveProjectedGeometryOn_of_boundaryReachabilityData_and_dartSuccessorCycleEmbeddingData_and_selectedBoundaryArc_and_taitEdgeColoring_and_nonempty_selectedBoundaryInteriorCarrier_sharedInteriorPair}
rules out deriving the direct replacement packages merely from those boundary
order, coloring, and carrier facts.  The same regression now also proves
\path{not_theorem49ClosedWalkAnnulusCollarPositiveProjectedGeometryOn_sharedInteriorPair}
and
\path{not_theorem49SuccessorCycleAnnulusCollarPositiveProjectedGeometryOn_sharedInteriorPair},
with failed universal route-facing forms for the closed-walk and
successor-cycle annulus-collar packages.  Thus even the sharpened source-facing
packages require genuine annulus-collar or witness-cover geometry, not just a
repacking of the available boundary-order, coloring, and carrier facts.
\path{Theorem49PositiveGeometricSourceConstruction.lean} now supplies the first
positive inhabitant of the replacement surface.  It reuses the two-collar
annulus geometry benchmark and proves
\path{counterEmbedding_selectedBoundaryInteriorEdgeEndpointVertices_nonempty}:
the intermediate interior edge contributes endpoints which are not incident to
the selected ambient boundary and therefore survive purification.  With the
new exact support lemmas
\path{counterEmbedding_interiorEdgeSupport_eq},
\path{counterEmbedding_interiorEdgeEndpointSupport_eq},
\path{counterEmbedding_selectedBoundaryEndpointSupport_eq}, and
\path{counterEmbedding_endpointSupport_disjoint_selectedBoundarySupport}, Lean
now records that this survival is precisely endpoint-support disjointness for
the shared intermediate edge.  The same calculation gives the named local
witness \path{counterEmbedding_hasUnblockedInteriorEndpoint}, matching the
first-order carrier obligation used by the replacement route.  With the
existing annulus-collar conversions, Lean then constructs both
\path{counterEmbedding_collarLayerPositiveProjectedGeometryOn} and
\path{counterEmbedding_heightOrderedPositiveProjectedGeometryOn}, together with
the graph-level packages
\path{counterGraph_collarLayerPositiveProjectedGeometry} and
\path{counterGraph_heightOrderedPositiveProjectedGeometry}.  The same benchmark
now carries an explicit constant nonzero Tait coloring
\path{counterGraphTaitEdgeColoring}, and Lean proves
\path{counterEmbedding_boundaryRootNonemptyProjectedSynthesis}; the concrete
two-collar geometry therefore reaches the nonempty projected theorem-4.9
endpoint, not merely the precursor package.  Lean now also records the same
positive benchmark at the raw quotient/error endpoint:
\path{counterEmbedding_collarGeometry_explicitTait_nonemptyCarrier_to_theorem49RawQuotientErrorPackage}
bundles the actual collar geometry, explicit Tait coloring, nonempty purified
carrier, nonempty projected synthesis package, and the raw Definition 4.8
representative plus selected-boundary-kernel error decomposition for every
Kirchhoff chain on this embedding.  The same module now also records the
source-side limitation of this benchmark:
\path{counterGraph_isAcyclic} proves that its graph is a three-edge matching,
\path{not_nonempty_closedWalkAnnulusBoundarySource_counterEmbedding} blocks an
honest closed-walk annulus source on the fixed embedding, and
\path{counterEmbedding_rawQuotientErrorPackage_without_closedWalkAnnulusBoundarySource}
packages the raw quotient/error endpoint together with that obstruction.  The
same benchmark now also records the
repaired-witness calibration
\path{counterEmbedding_previousBoundaryWitness_nonemptyProjectedSynthesis_without_parentWitnessOrientation}:
the repaired previous-boundary witness geometry and explicit endpoint witness
coexist with the nonempty projected endpoint, while the stricter canonical
annulus-construction parent-witness orientation still fails on the same collar.
Thus the replacement endpoint is not vacuous, and it should not be conflated
with the stronger parent-orientation route; the remaining burden is to lift
this positive construction from a benchmark to route-facing boundary-order
hypotheses.  The generic bridge has
also been factored explicitly: an existential annulus collar geometry together
with nonempty purified carrier now reaches
\path{Theorem49BoundaryRootNonemptyProjectedSynthesis} for any supplied Tait
coloring via
\path{exists_theorem49BoundaryRootNonemptyProjectedSynthesis_of_exists_annulusCollarGeometry_and_nonempty_selectedBoundaryInteriorEdgeEndpointVertices_direct}.
The same bridges now also accept the named local endpoint witness
\path{HasUnblockedInteriorEndpoint}: the fixed-embedding constructors
\path{theorem49CollarLayerPositiveProjectedGeometryOn_of_annulusCollarGeometry_and_hasUnblockedInteriorEndpoint}
and
\path{theorem49HeightOrderedPositiveProjectedGeometryOn_of_annulusCollarGeometry_and_hasUnblockedInteriorEndpoint},
their graph-level package constructors, and
\path{exists_theorem49BoundaryRootNonemptyProjectedSynthesis_of_exists_annulusCollarGeometry_and_hasUnblockedInteriorEndpoint_direct}
let a future source proof supply one unblocked interior endpoint directly,
without first repackaging the purified carrier set.
The same endpoint witness can now be obtained from the older endpoint-support
separation condition: raw \path{interiorEdgeEndpointSupport} nonemptiness
together with disjointness from selected-boundary endpoint support implies
\path{HasUnblockedInteriorEndpoint}, and the annulus-collar projected endpoint
has direct closed-walk and successor-cycle source wrappers for this form.
This bridge now reaches the concrete collar geometry layer as well:
\path{PlanarBoundaryAnnulusCollarGeometry.hasUnblockedInteriorEndpoint_of_peelFaces_endpoint_disjoint_selectedBoundarySupport}
applies peeled-face endpoint no-touch on the canonical construction
\path{planarBoundaryAnnulusConstruction_of_annulusCollarGeometry data}
and raw \path{interiorEdgeEndpointSupport} nonemptiness to produce
\path{HasUnblockedInteriorEndpoint}.  Consequently
\path{Theorem49PositiveGeometricSourceReplacement.lean} now exposes direct
finite-collar, height-ordered, and projected-synthesis constructors from
annulus collar geometry plus that peeled-face no-touch obligation, without an
intermediate purified-carrier proof.
The source-facing file
\path{Theorem49PositiveGeometricSourceReplacementRoute.lean} now exposes the
same peeled-face no-touch surface at the honest source layer: closed-walk and
successor-cycle source data, same-embedding annulus collar geometry, canonical
peeled-face endpoint no-touch, and raw
\path{interiorEdgeEndpointSupport} nonemptiness directly inhabit the
route-facing annulus-collar positive packages and reach the nonempty projected
synthesis endpoint for any supplied Tait coloring.
The same source-layer wrappers now also return the graph-level finite-collar
and height-ordered replacement packages directly.  Thus an honest closed-walk
or successor-cycle caller carrying the peeled-face no-touch input can consume
the replacement API without first building the route-facing annulus-collar
package and then lowering it by a separate step.
This has now been pushed one step further in
\path{Theorem49Synthesis.lean}.  The theorem
\path{theorem49BoundaryRootSynthesis_of_planarBoundaryAnnulusConstruction_of_parentWitness}
assembles the full purified theorem-4.9 synthesis package directly from
\path{PlanarBoundaryAnnulusConstruction} plus
\path{ParentWitnessOrientation}: besides the previously isolated span,
separation, kernel, and finrank consequences, it also supplies the corrected
raw-source image identity, projected/raw finite generator representations, raw
Kirchhoff representatives, and the boundary-kernel decomposition inside one
route-level object.  The graph-level theorem
\path{exists_theorem49BoundaryRootSynthesis_of_admitsPlanarBoundaryAnnulusConstructionParentWitnessOrientation}
shows that the remaining constructive burden is now upstream of synthesis
itself: once actual boundary-order geometry yields honest parent-witness
orientation on an annulus construction, Lean already produces the whole
purified theorem-4.9 endpoint without passing back through the older
interior-dual wrapper stack.

This also records a precise blocker in Lean:
\path{selectedBoundaryInteriorEdgeEndpointVertices_ne_interiorEdgeEndpointSupport_of_endpoint_shared}
says that if a face-incidence interior edge and a selected-boundary edge share
a primal endpoint, then raw interior-edge endpoints cannot already satisfy the
boundary-erasure vertex side condition.  Thus edge-level disjointness of
\path{selectedBoundarySupport} and \path{interiorEdgeSupport} is not the same
as the needed vertex-level separation.

This is not just a theoretical concern.  The module
\path{Theorem49InteriorVerticesEndpointObstruction.lean} builds a two-edge
path incidence model in which the selected-boundary support and interior-edge
support are disjoint as edge sets, but the selected-boundary edge and the
interior edge share a primal endpoint.  The theorem
\path{edge_disjointness_does_not_imply_endpoint_separation_sharedEndpoint}
therefore refutes the shortcut from incidence-count edge disjointness to
endpoint separation, and
\path{not_endpointSupport_disjoint_sharedEndpoint} records the same failure as
non-disjointness of the finite endpoint supports.  The theorem
\path{selectedBoundaryInteriorEdgeEndpointVertices_ne_interiorEdgeEndpointSupport_sharedEndpoint}
shows that the raw endpoint carrier is genuinely too large in this model.
The same countermodel now also admits linear endpoint-sharing face-boundary
data: \path{nonempty_orderedFaceBoundaryData_sharedEndpoint} and
\path{nonempty_faceBoundaryRunGeometry_sharedEndpoint} build the ordered-run
shell, while
\path{orderedFaceBoundaryData_does_not_imply_endpointSupport_disjoint_sharedEndpoint}
records that endpoint-support disjointness still fails.  Thus a merely linear
boundary-order/run interface cannot discharge the raw endpoint-separation
burden; the route needs either an explicit endpoint-separation theorem or a
stronger geometric source, such as honest cyclic/closed-walk data tied to the
annular interior.
The file now also records the complementary fact that this particular
countermodel is only a linear-run obstruction.  The theorem
\path{not_nonempty_cyclicOrderedFaceBoundaryData_sharedEndpoint}, together
with \path{not_nonempty_faceBoundaryClosedRunGeometry_sharedEndpoint} and
\path{not_nonempty_faceBoundaryClosedWalkGeometry_sharedEndpoint}, shows that
the singleton interior face prevents cyclic closure or honest closed-walk
data.  The combined theorem
\path{linearRun_endpointSupport_obstruction_not_closedWalk_sharedEndpoint}
therefore localizes the next viable target: cyclic/closed-walk geometry may
still be the right source for the missing endpoint-separation proof, but the
linear boundary-run shell definitely is not.
The next regression closes the tempting cyclic-only shortcut as well.  The
two-triangle diamond model \path{diamondEmbedding} has cyclic closed
face-boundary runs on both triangular faces, witnessed by
\path{nonempty_faceBoundaryClosedRunGeometry_diamond} and the underlying
\path{diamondCyclicOrderedFaceBoundaryData}.  Nevertheless, its shared edge is
an interior edge and shares a primal endpoint with the selected-boundary edge
\path{d01}.  The theorem
\path{cyclicClosedRun_does_not_imply_endpointSupport_disjoint_diamond} records
that cyclic closed-run data alone still does not imply endpoint-support
disjointness.
This obstruction now reaches the honest facial-walk interface as well:
\path{diamondFaceBoundaryPureDartCycleGeometry} supplies oriented dart cycles
for the two triangular faces and
\path{nonempty_faceBoundaryClosedWalkGeometry_diamond} lowers them to actual
facial closed walks, while
\path{faceBoundaryClosedWalkGeometry_does_not_imply_endpointSupport_disjoint_diamond}
still refutes endpoint-support disjointness.  The remaining geometric burden
is therefore not merely cyclic ordering or the existence of facial closed
walks.  The strengthened regression
\path{closedWalkSelectedArc_does_not_imply_endpointSupport_disjoint_diamond}
rotates the same diamond facial walks so that the selected-boundary edges form
a contiguous selected arc on each induced linear face run, and endpoint-support
disjointness still fails.  Thus the missing ingredient is not just closed
walks plus selected arcs: it is an annular no-touch/separation theorem, or an
explicit purification statement, for the intended interior vertices.
The latest regression lifts this warning to the actual closed-walk annulus
boundary-source package.  The model
\path{diamondWithTriangleEmbedding} keeps the diamond endpoint failure in one
selected-boundary component and adds a disjoint triangular selected-boundary
component, so that the two-root shared-endpoint boundary-reachability data
\path{diamondWithTriangleAnnulusBoundaryReachabilityData} is available.  It
also carries honest facial closed walks with selected arcs, packaged as
\path{diamondWithTriangleClosedWalkAnnulusBoundarySource}.  Nevertheless
\path{closedWalkAnnulusBoundarySource_does_not_imply_endpointSupport_disjoint_diamondWithTriangle}
still refutes endpoint-support disjointness.  This is intentionally not a
claim against the later peeled-face no-touch bridge; it says that the
closed-walk annulus boundary source alone cannot replace the missing
endpoint-level separation hypothesis.  That bridge is now explicit at the
generic two-sided annulus-root surface:
\path{Theorem49InteriorVerticesAnnulusConstruction.lean} defines
\path{planarBoundaryAnnulusConstruction_of_planarBoundaryAnnulusRootAdjDistancePeelData}
and proves the corresponding annulus-root lossless-purification theorem from
peeled-face endpoint no-touch, while \path{Theorem49Synthesis.lean} proves
\path{theorem49AnnulusRootNonemptySynthesis_of_planarBoundaryAnnulusRootAdjDistancePeelData_of_peelFaces_endpoint_disjoint_selectedBoundarySupport}.
The closed-walk theorem
\path{theorem49AnnulusRootNonemptySynthesis_of_closedWalkAnnulusBoundarySource_and_interiorDualBoundaryRootAdjDistancePeelData_of_peelFaces_endpoint_disjoint_selectedBoundarySupport}
is now the source-facing specialization: it uses the canonical construction's
peeled-face endpoint no-touch condition to derive global raw-carrier endpoint
disjointness before invoking the non-vacuous synthesis wrapper.  The abstract
obstruction
\path{false_of_interiorDualBoundaryRootAdjDistancePeelData_of_forcing_interiorEdge}
now isolates the same mechanism at the peeled interior-dual layer: if some
interior edge is forced to meet selected-boundary support on every incident
face, then covering that edge would require a peeled face that violates the
edge-disjoint \path{hpeelInterior} field.  The calibration theorem
\path{not_nonempty_interiorDualBoundaryRootAdjDistancePeelData_diamondWithTriangle}
is the concrete specialization of this criterion to the unique interior edge
\path{dt12}.
The direct incompatibility theorem
\path{closedWalkAnnulusBoundarySource_does_not_imply_interiorDualBoundaryRootAdjDistancePeelData_diamondWithTriangle}
packages this in the exact shape of the current positive route: on this
embedding the honest closed-walk source exists, but the peeled interior-dual
package does not.
The same source model also survives the natural source-side strengthening to
cyclic ordered face-arc data.  The closed-walk witness upgrades to
\path{diamondWithTriangleCyclicOrderedFaceArcEmbeddingData}, and the theorem
\path{annulusBoundaryReachability_and_cyclicOrderedFaceArcEmbeddingData_does_not_imply_interiorDualBoundaryRootAdjDistancePeelData_diamondWithTriangle}
shows that even after adjoining the live annulus boundary-reachability data and
the stronger cyclic per-face boundary order with selected-boundary contiguity,
the peeled interior-dual package is still absent.
The same contradiction is now formalized one layer earlier, before any
interior-dual data are introduced.  The theorem
\path{false_of_planarBoundaryAnnulusConstruction_and_hpeelInterior_diamondWithTriangle}
shows that on the same embedding, any annulus construction with the usual
peeled-interior disjointness field already fails: the unique interior edge
\path{dt12} must be covered by some peeled witness face, but every face
containing \path{dt12} also contains selected-boundary support.  Hence even
the weakest current construction shell
\path{PlanarBoundaryAnnulusConstructionBoundarySupportFaceData} is impossible,
as recorded by
\path{not_nonempty_planarBoundaryAnnulusConstructionBoundarySupportFaceData_diamondWithTriangle}
and the summary theorem
\path{closedWalkAnnulusBoundarySource_does_not_imply_planarBoundaryAnnulusConstructionBoundarySupportFaceData_diamondWithTriangle}.
The cyclic ordered face-arc strengthening still does not bypass this earlier
geometric failure either: the theorem
\path{annulusBoundaryReachability_and_cyclicOrderedFaceArcEmbeddingData_does_not_imply_planarBoundaryAnnulusConstructionBoundarySupportFaceData_diamondWithTriangle}
shows that the same embedding has annulus roots and cyclic ordered face-arc
data, yet still does not reach even the weakest peeled-interior construction
shell.
The same obstruction now survives after discarding the cyclic face order and
retaining only the bundled per-face selected-boundary arc witness.
\path{diamondWithTriangleSelectedBoundaryArcGeometry} is extracted from the same
embedding, and the theorems
\path{annulusBoundaryReachability_and_selectedBoundaryArcGeometry_does_not_imply_interiorDualBoundaryRootAdjDistancePeelData_diamondWithTriangle}
and
\path{annulusBoundaryReachability_and_selectedBoundaryArcGeometry_does_not_imply_planarBoundaryAnnulusConstructionBoundarySupportFaceData_diamondWithTriangle}
show that even annulus roots plus the weaker
\path{PlanarBoundarySelectedBoundaryArcGeometry} shell still do not start the
peeled interior-dual route or even its weakest construction shell on this
model.
There is now also a same-embedding regression strictly earlier in the source
ladder.  In \path{PlanarBoundaryFaceBoundaryRunRegression.lean}, the explicit
two-face model \path{doubleStarEmbedding} carries both
\path{doubleStarAnnulusBoundaryReachabilityData} and
\path{doubleStarCyclicOrderedFaceArcEmbeddingData}, but
\path{not_nonempty_planarBoundaryClosedWalkAnnulusBoundarySource_doubleStarEmbedding}
shows that it still admits no honest closed-walk annulus source.  The left
ambient face has the three-star boundary, and any candidate facial closed walk
there transfers to the acyclic graph \path{doubleStarLeftGraph}.  So annulus
boundary reachability plus cyclic ordered face-arc data do not yet recover the
honest facial closed-walk source; this is summarized abstractly by
\path{not_forall_nonempty_planarBoundaryClosedWalkAnnulusBoundarySource_of_nonempty_planarBoundaryAnnulusBoundaryReachabilityData_and_nonempty_planarBoundaryCyclicOrderedFaceArcEmbeddingData}.
The same two-face model also carries the weaker bundled shell
\path{PlanarBoundarySelectedBoundaryArcGeometry}, via
\path{nonempty_planarBoundarySelectedBoundaryArcGeometry_doubleStarEmbedding},
and the theorems
\path{exists_embedding_withPlanarBoundaryAnnulusBoundaryReachabilityData_andPlanarBoundarySelectedBoundaryArcGeometry_withoutClosedWalkAnnulusBoundarySource}
and
\path{not_forall_nonempty_planarBoundaryClosedWalkAnnulusBoundarySource_of_nonempty_planarBoundaryAnnulusBoundaryReachabilityData_and_nonempty_planarBoundarySelectedBoundaryArcGeometry}
show that even annulus roots plus one contiguous selected-boundary arc on each
ambient face still do not recover the honest closed-walk annulus source.
The positive direction is now explicit too: the source file
\path{PlanarBoundaryClosedWalkSource.lean} lowers any honest
\path{PlanarBoundaryClosedWalkAnnulusBoundarySource} directly to both
\path{AdmitsPlanarBoundaryAnnulusBoundaryReachabilityData} together with
\path{AdmitsPlanarBoundaryCyclicOrderedFaceArcEmbeddingData}, and to
\path{AdmitsPlanarBoundaryAnnulusBoundaryReachabilityData} together with
\path{AdmitsPlanarBoundarySelectedBoundaryArcGeometry}.  So the double-star
and related regression theorems should be read as genuine failed converses of
an available positive lowering, not merely as mismatched interfaces.  This is
now graph-level as well, not only same-embedding: the file proves
\path{doubleStarGraph_isAcyclic}, then uses the reusable acyclic blocker to
derive
\path{not_admitsPlanarBoundaryClosedWalkAnnulusBoundarySource_doubleStarGraph}.
Consequently the new theorems
\path{not_forall_admitsPlanarBoundaryClosedWalkAnnulusBoundarySource_of_admitsPlanarBoundaryAnnulusBoundaryReachabilityData_and_cyclicOrderedFaceArcEmbeddingData}
and
\path{not_forall_admitsPlanarBoundaryClosedWalkAnnulusBoundarySource_of_admitsPlanarBoundaryAnnulusBoundaryReachabilityData_and_selectedBoundaryArcGeometry}
show that even after existentially quantifying over embeddings, the paired
annulus-rooted shells still do not force the honest closed-walk annulus
source.  The same double-star witness now blocks rotation-side repairs too:
\path{not_admitsPlanarBoundaryDartCycleEmbeddingData_doubleStarGraph},
\path{not_admitsPlanarBoundaryPureDartCycleEmbeddingData_doubleStarGraph}, and
\path{not_admitsPlanarBoundaryDartSuccessorCycleEmbeddingData_doubleStarGraph}
show that the graph admits no dart-cycle, pure dart-cycle, or successor-cycle
source on any embedding.  The failed-converse wrappers
\path{not_forall_admitsPlanarBoundaryDartCycleEmbeddingData_of_admitsPlanarBoundaryAnnulusBoundaryReachabilityData_and_selectedBoundaryArcGeometry},
\path{not_forall_admitsPlanarBoundaryPureDartCycleEmbeddingData_of_admitsPlanarBoundaryAnnulusBoundaryReachabilityData_and_selectedBoundaryArcGeometry},
and
\path{not_forall_admitsPlanarBoundaryDartSuccessorCycleEmbeddingData_of_admitsPlanarBoundaryAnnulusBoundaryReachabilityData_and_selectedBoundaryArcGeometry}
therefore rule out the obvious repair of replacing closed facial walks by
those stronger rotation-side source packages.  This is now abstracted in the source file itself, not just in the
double-star regression: the generic theorems
\path{admitsPlanarBoundaryAnnulusBoundaryReachabilityData_and_cyclicOrderedFaceArcEmbeddingData_withoutClosedWalkAnnulusBoundarySource_of_isAcyclic_of_nonempty_edgeSet},
\path{admitsPlanarBoundaryAnnulusBoundaryReachabilityData_and_selectedBoundaryArcGeometry_withoutClosedWalkAnnulusBoundarySource_of_isAcyclic_of_nonempty_edgeSet},
and their two existential failed-converse corollaries show that any acyclic
positive-edge witness graph carrying one of these paired shells already blocks
the route to the honest closed-walk annulus source.
This obstruction is now also factored through one reusable acyclic theorem in
the shared walk-embedding layer.  The file
\path{WalkPlanarEmbeddingWithBoundaryAcyclic.lean} proves
\path{not_nonempty_faceBoundaryClosedWalkGeometry_of_isAcyclic_of_nonempty_edgeSet}:
any acyclic graph with at least one edge admits no honest facial closed-walk
geometry on any boundary-aware embedding, because some ambient face must exist
and its nonempty closed trail would be a path, hence nil.  The FourColor
wrapper
\path{not_nonempty_planarBoundaryClosedWalkEmbeddingData_of_isAcyclic_of_nonempty_edgeSet}
exposes the same fact directly on the route-facing source interface.  So the
path, star, and peeled-endpoint-touch source impossibility lemmas are now
instances of one reusable theorem rather than separate local arguments.  This
abstraction now reaches one layer higher as well: the route-facing theorems
\path{not_admitsPlanarBoundaryDartCycleEmbeddingData_of_isAcyclic_of_nonempty_edgeSet},
\path{not_admitsPlanarBoundaryPureDartCycleEmbeddingData_of_isAcyclic_of_nonempty_edgeSet},
\path{not_admitsPlanarBoundaryDartSuccessorCycleEmbeddingData_of_isAcyclic_of_nonempty_edgeSet},
and
\path{not_admitsPlanarBoundaryClosedWalkAnnulusBoundarySource_of_isAcyclic_of_nonempty_edgeSet}
show that every acyclic positive-edge model already blocks the stronger
rotation-side and annulus-source interfaces before any model-specific geometry
is used.
The same witness pattern is now abstracted one step further at the failed-converse
level: \path{not_forall_admitsPlanarBoundaryClosedWalkEmbeddingData_of_exists_isAcyclic_of_nonempty_edgeSet_and_property}
and
\path{not_forall_admitsPlanarBoundaryClosedWalkAnnulusBoundarySource_of_exists_isAcyclic_of_nonempty_edgeSet_and_property}
package the generic reason why any acyclic positive-edge witness carrying a
shell $P$ refutes a universal theorem ``$P$ implies honest closed walks'' or
``$P$ implies the honest annulus source''.  The star/path regression theorems
for cyclic ordered arcs, selected-boundary arcs, cyclic face-boundary vertex
walks, and ordered face arcs now instantiate these generic criteria directly.
This is no longer just a one-off gadget computation: the file now isolates
the reusable obstruction
\path{false_of_planarBoundaryAnnulusConstruction_and_hpeelInterior_of_forcing_interiorEdge},
which says abstractly that if some interior edge is forced to meet
selected-boundary support on every incident face, then the peeled-interior
annulus shell is impossible.
This forcing-edge mechanism is now factored into the smaller route-facing
wrapper file \path{Theorem49ForcingInteriorEdgeObstruction.lean}.  Its witness
structure \path{ForcingInteriorEdgeWitness} packages exactly the charitable
repair burden: one designated interior edge together with a proof that every
incident face already carries selected-boundary support.  From that single
witness, the new generic theorems
\path{not_nonempty_interiorDualBoundaryRootAdjDistancePeelData},
\path{not_nonempty_planarBoundaryAnnulusConstructionBoundarySupportFaceData},
and
\path{closedWalkAnnulusBoundarySource_does_not_imply_planarBoundaryAnnulusConstructionBoundarySupportFaceData}
recover the obstruction directly at the route surface, without reusing the
diamond specialization by hand.  The construction-side positive invariant is
now stated in \path{PlanarBoundaryAnnulusConstructionCore.lean} as
\path{PlanarBoundaryAnnulusConstruction.exists_peelFace_witnessEdge_eq_and_faceBoundary_disjoint_selectedBoundarySupport_of_mem_interiorEdgeSupport}:
any annulus construction whose peeled faces avoid selected-boundary support
must assign every interior edge to a witness peeled face whose entire boundary
avoids that support.  The forcing-edge contradiction is now the direct negation
of this invariant.  This negation has also been isolated as the positive
face-incidence predicate
\path{EveryInteriorEdgeHasBoundaryFreeIncidentFace}: every interior edge has
some incident ambient face whose whole boundary is disjoint from selected
boundary support.  The theorem
\path{not_nonempty_forcingInteriorEdgeWitness_iff_everyInteriorEdgeHasBoundaryFreeIncidentFace}
proves that this predicate is exactly the absence of a
\path{ForcingInteriorEdgeWitness}, and
\path{everyInteriorEdgeHasBoundaryFreeIncidentFace_of_interiorDualBoundaryRootAdjDistancePeelData}
now proves that the interior-dual boundary-root adjacency-distance package
itself supplies this predicate: its cover field assigns every interior edge to
a peeled witness face, and its peeled-interior field makes that face disjoint
from selected-boundary support.  The corollary
\path{not_nonempty_forcingInteriorEdgeWitness_of_interiorDualBoundaryRootAdjDistancePeelData}
therefore records the source-side necessary condition for deriving
\path{InteriorDualBoundaryRootAdjDistancePeelData}: forcing interior edges must
already have been eliminated before any later annulus-construction or selector
compatibility step is relevant.  Separately,
\path{everyInteriorEdgeHasBoundaryFreeIncidentFace_of_planarBoundaryAnnulusConstructionBoundarySupportFaceData}
shows that the weak boundary-support construction shell supplies it.  The
same file now also exposes the data-valued companion
\path{BoundaryFreeIncidentFaceSelector}, whose nonemptiness is equivalent to
\path{EveryInteriorEdgeHasBoundaryFreeIncidentFace}; the noncomputable
extraction
\path{boundaryFreeIncidentFaceSelector_of_planarBoundaryAnnulusConstructionBoundarySupportFaceData}
turns the weak construction shell into an actual per-interior-edge choice of a
selected-boundary-free incident face.  This selector is deliberately only a
local face-choice layer: it does not yet assert compatibility among the chosen
faces for distinct interior edges.  The selector/forcing relationship is now
also available without the intermediate predicate:
\path{nonempty_boundaryFreeIncidentFaceSelector_iff_not_nonempty_forcingInteriorEdgeWitness}
and
\path{not_nonempty_boundaryFreeIncidentFaceSelector_iff_nonempty_forcingInteriorEdgeWitness}
state that a boundary-free selector exists exactly when no forcing interior edge
exists.  That compatibility layer is now partially
formalized in \path{Theorem49BoundaryFreeSelectorConstruction.lean}.  The
definition \path{BoundaryFreeIncidentFaceSelector.peelFaces} takes the finite
image of the selected faces, and
\path{BoundaryFreeIncidentFaceSelector.selectedWitnessEdge} inverts the
selector on that image.  If the selector is injective on interior edges and
every non-witness remainder edge is either ambient-boundary support or the
selected witness of a strictly deeper selected face, then
\path{BoundaryFreeIncidentFaceSelector.toPlanarBoundaryAnnulusConstruction}
builds a genuine \path{PlanarBoundaryAnnulusConstruction}.  The companion
lemma
\path{BoundaryFreeIncidentFaceSelector.toPlanarBoundaryAnnulusConstruction_hpeelInterior}
proves that the constructed peeled faces are disjoint from selected-boundary
support.  The strict-descent compatibility condition also now has its own
selector-level cycle test: \path{BoundaryFreeIncidentFaceSelector.remainderDependency}
records when the selected witness of one peeled face appears as a non-witness,
non-boundary remainder edge of another, and
\path{BoundaryFreeIncidentFaceSelector.faceDistance_lt_of_remainderDependency}
turns each such dependency into a strict face-distance inequality.  The theorem
\path{BoundaryFreeIncidentFaceSelector.false_of_remainderDependency_cycle}
therefore rules out every nonempty finite closed selector-dependency chain
under the same strict-distance hypothesis consumed by the construction bridge.
The new comparison theorem
\path{BoundaryFreeIncidentFaceSelector.toPlanarBoundaryAnnulusConstruction_remainderDependency_iff}
shows that after lowering a strict-descent selector to a
\path{PlanarBoundaryAnnulusConstruction}, this selector dependency relation is
exactly the construction-level
\path{PlanarBoundaryAnnulusConstruction.remainderDependency}; the strict
face-distance field on the construction side is recovered from selector
\path{hrest}.  Thus selector-side cycle diagnostics transfer to the central
annulus construction without a second dependency argument.  The companion
\path{BoundaryFreeIncidentFaceSelector.toPlanarBoundaryAnnulusConstruction_no_remainderDependency_iff}
shows that terminal no-dependency is preserved exactly by the same lowering.
The maximum-distance strengthening
\path{BoundaryFreeIncidentFaceSelector.false_of_total_remainderDependency}
rules out a nonempty finite selected-face set in which every selected face has
an outgoing non-boundary remainder dependency, and
\path{BoundaryFreeIncidentFaceSelector.exists_face_without_remainderDependency}
extracts a terminal selected face with no such outgoing dependency.  This is a
more usable diagnostic than an explicit cycle enumeration: any successful
source-to-selector construction must turn that terminal face into a real
ambient-boundary break, or provide a separate well-founded orientation.
The selector terminal condition now has the same local characterization as the
construction terminal condition:
\path{BoundaryFreeIncidentFaceSelector.no_remainderDependency_iff_boundary_remainders}
states that, on a selected face, no outgoing selector dependency is equivalent
to every non-witness remainder edge lying on the outer-or-inner ambient annulus
boundary.
The theorem
\path{BoundaryFreeIncidentFaceSelector.exists_terminal_face_boundary_remainders}
now states that boundary break directly: for some selected face, every
non-witness remainder edge is already in the outer-or-inner ambient annulus
boundary.  This is the local form of the selector construction burden, separated
from the later non-peeled-face and frontier obligations.  The variants
\path{BoundaryFreeIncidentFaceSelector.peelFaces_nonempty_of_mem_interiorEdgeSupport},
\path{BoundaryFreeIncidentFaceSelector.exists_terminal_face_boundary_remainders_of_mem_interiorEdgeSupport},
and
\path{BoundaryFreeIncidentFaceSelector.exists_terminal_face_boundary_remainders_of_interiorEdgeSupport_nonempty}
remove the separate selected-face nonemptiness premise as soon as a live
interior edge is present.
The same maximum-distance principle is now available at the central construction
surface as
\path{PlanarBoundaryAnnulusConstruction.exists_terminal_peelFace_boundary_remainders}:
any nonempty BFS-stratified annulus construction has a maximum-distance peeled
face whose every non-witness remainder edge is already on the ambient outer or
inner boundary split.  This theorem now factors through the construction-level
strict dependency relation
\path{PlanarBoundaryAnnulusConstruction.remainderDependency}.  Non-boundary
remainders yield such dependencies by
\path{PlanarBoundaryAnnulusConstruction.exists_remainderDependency_of_nonboundary_remainder},
closed finite dependency chains are ruled out by
\path{PlanarBoundaryAnnulusConstruction.false_of_remainderDependency_cycle},
and the total-outgoing form
\path{PlanarBoundaryAnnulusConstruction.false_of_total_remainderDependency}
produces
\path{PlanarBoundaryAnnulusConstruction.exists_terminal_peelFace_no_remainderDependency}.
The local terminal condition is now characterized exactly by
\path{PlanarBoundaryAnnulusConstruction.no_remainderDependency_iff_boundary_remainders}:
for a peeled face, no outgoing construction dependency is equivalent to all
non-witness remainders lying on the ambient annulus boundary.  Its forward
direction is also available as
\path{PlanarBoundaryAnnulusConstruction.boundary_remainders_of_no_remainderDependency}.
Thus the terminal boundary break is now a named finite strict-descent
obstruction on the reusable annulus-construction surface.  The nonemptiness side
condition has now been pushed back
to the construction cover itself:
\path{PlanarBoundaryAnnulusConstruction.peelFaces_nonempty_of_mem_interiorEdgeSupport}
shows that any live interior edge covered by the construction forces a peeled
face, and
\path{PlanarBoundaryAnnulusConstruction.exists_terminal_peelFace_boundary_remainders_of_mem_interiorEdgeSupport}
packages the resulting terminal boundary break directly.  Thus the terminal
boundary break is not just a selector artifact; it is a direct consequence of
the construction's strict-descent \path{hrest} field as soon as the construction
actually covers a live interior edge.
\path{Theorem49PlanarBoundaryAnnulusConstructionOrientationObstruction.lean}
now also records the corresponding negative calibration:
\path{not_forall_live_planarBoundaryAnnulusConstruction_terminalNoDependency_implies_parentWitnessOrientation}
shows that live interior-edge support together with a terminal peeled face with
no outgoing construction-level \path{remainderDependency} still does not force
\path{PlanarBoundaryAnnulusConstruction.ParentWitnessOrientation}.  The
terminal dependency break is therefore a diagnostic for unbroken strict
remainder cycles, not the missing well-founded parent-witness orientation
itself.
The same finite counterconstruction also lowers to
\path{PlanarBoundaryAnnulusConstructionPositiveFrontierData} and
\path{PlanarBoundaryAnnulusConstructionBoundarySupportFaceData} while still
refuting parent orientation, as recorded by
\path{not_forall_planarBoundaryAnnulusConstructionPositiveFrontierData_implies_parentWitnessOrientation}
and
\path{not_forall_planarBoundaryAnnulusConstructionBoundarySupportFaceData_implies_parentWitnessOrientation}.
Thus selected-boundary contact and positive current-frontier bookkeeping do not
by themselves supply the missing well-founded parent witness.
The construction-level obstruction is sharper than merely pointing at the
zero stratum:
\path{not_forall_planarBoundaryAnnulusConstruction_positive_faceDistance_implies_parentWitnessOrientation}
uses a variant in which every peeled face has positive stored distance, but the
bad witness still has no incident face of strictly smaller distance.  At the
same time, this positivity condition cannot be imposed as a repair on the
positive-frontier shell itself:
\path{PlanarBoundaryAnnulusConstructionPositiveFrontierData.exists_peelFace_faceDistance_zero}
and
\path{PlanarBoundaryAnnulusConstructionPositiveFrontierData.not_forall_peelFaces_positive}
show that positive-frontier data necessarily contains a zero-distance peeled
stratum.  The construction core now factors this as a stratum-level clash:
\path{PlanarBoundaryAnnulusConstruction.not_forall_stratumFaces_nonempty_of_parentWitness}
shows that parent orientation makes the zero stratum empty, while the current
positive-frontier axioms inhabit every stratum.  The core construction file now
records the stronger incompatibility as
\path{PlanarBoundaryAnnulusConstructionPositiveFrontierData.not_parentWitnessOrientation},
with corresponding face-partition and boundary-support variants.  Thus the
live parent-orientation burden is genuine lower-distance incidence, not scalar
positivity of the stored distances or the current zero-based positive-frontier
shell.
\path{PlanarBoundaryAnnulusConstruction.lean} now makes the source-side
mismatch explicit as
\path{not_exists_planarBoundaryAnnulusConstructionPositiveFrontierData_with_rootDistanceConstruction_of_boundaryReachabilityData_and_interiorDualBoundaryRootAdjDistancePeelData}:
the parent-oriented root-distance construction produced from boundary
reachability plus interior-dual boundary-root adjacency-distance data cannot be
wrapped in the current positive-frontier shell.  Thus the root-distance
parent-orientation route and the zero-based frontier bookkeeping route are
different construction surfaces, not interchangeable views of one object.  The
same construction file now records the exact reindexing:
\path{rootDistance_faceDistance_succ_filter_eq_shifted_faceDistance_filter_of_boundaryReachabilityData_and_interiorDualBoundaryRootAdjDistancePeelData}
says that root-distance layer \(n+1\) on peeled faces is exactly shifted
construction layer \(n\), and
\path{rootDistance_faceDistance_one_filter_eq_shifted_faceDistance_zero_filter_of_boundaryReachabilityData_and_interiorDualBoundaryRootAdjDistancePeelData}
specializes this to root-distance layer \(1\) versus shifted layer \(0\).
\path{rootDistance_faceDistance_one_filter_nonempty_iff_not_shifted_parentWitnessOrientation_of_boundaryReachabilityData_and_interiorDualBoundaryRootAdjDistancePeelData}
then sharpens the interpretation: nonempty root-distance layer \(1\) is exactly
failure of parent-witness orientation for the shifted construction.
The stronger bridge
\path{BoundaryFreeIncidentFaceSelector.toPlanarBoundaryAnnulusConstructionBoundarySupportFaceData}
then reaches
\path{PlanarBoundaryAnnulusConstructionBoundarySupportFaceData} once the
remaining annulus-side hypotheses are supplied explicitly: every non-peeled face
touches selected-boundary support, outer and inner boundary faces are separated,
and the positive current outer frontiers are nonempty and pairwise disjoint.
Thus the selector supplies the construction cover and peeled-face
selected-boundary avoidance, while the real unresolved geometric burden remains
visible as face/frontier data rather than hidden in the witness-choice layer.
The theorem
\path{BoundaryFreeIncidentFaceSelector.not_nonempty_forcingInteriorEdgeWitness_of_boundarySupportFaceData_conditions}
packages this stronger bridge with the forcing-edge obstruction.
The positive endpoint of this selector lane is now recorded in
\path{Theorem49BoundaryFreeSelectorPositiveRoute.lean}, and it now factors
through a reusable construction-level bridge in
\path{Theorem49PositiveGeometricSourceReplacement.lean}.  The theorem
\path{theorem49HeightOrderedPositiveProjectedGeometryOn_of_planarBoundaryAnnulusConstruction}
shows that any \path{PlanarBoundaryAnnulusConstruction}, once paired with a
live purified carrier on the same embedding, already gives
\path{Theorem49HeightOrderedPositiveProjectedGeometryOn}; the companion
\path{theorem49BoundaryRootNonemptyProjectedSynthesis_of_planarBoundaryAnnulusConstruction}
adds a Tait coloring and reaches the corrected nonempty projected synthesis
endpoint.  The selector theorem
\path{BoundaryFreeIncidentFaceSelector.theorem49HeightOrderedPositiveProjectedGeometryOn_of_strictDescent}
shows that a boundary-free selector, an annulus boundary split, injectivity over
interior edges, the strict-descent remainder condition, and a nonempty purified
selected-boundary carrier supply such a construction-level bridge.  With a
supplied Tait coloring, the companion theorem
\path{BoundaryFreeIncidentFaceSelector.theorem49BoundaryRootNonemptyProjectedSynthesis_of_strictDescent}
reaches the nonempty projected boundary-root synthesis endpoint.  This does not
derive the selector or its compatibility from source data; it isolates the
precise annulus-construction positive target that a source-facing
boundary-order argument must now inhabit.  The same construction and live
carrier hypotheses also expose the terminal boundary-break face through
\path{theorem49TerminalPeelFaceBoundaryRemainders_of_planarBoundaryAnnulusConstruction},
so the positive endpoint and the local maximum-distance boundary break now live
on the same construction surface.  The construction surface also has the
matching local-endpoint entry points
\path{theorem49HeightOrderedPositiveProjectedGeometryOn_of_planarBoundaryAnnulusConstruction_and_hasUnblockedInteriorEndpoint},
\path{theorem49BoundaryRootNonemptyProjectedSynthesis_of_planarBoundaryAnnulusConstruction_and_hasUnblockedInteriorEndpoint},
and
\path{theorem49TerminalPeelFaceBoundaryRemainders_of_planarBoundaryAnnulusConstruction_and_hasUnblockedInteriorEndpoint},
so a source-side \path{HasUnblockedInteriorEndpoint} witness can be consumed
without first being repackaged as raw carrier nonemptiness.  The selector
positive route now uses this local endpoint language as well:
\path{BoundaryFreeIncidentFaceSelector.theorem49HeightOrderedPositiveProjectedGeometryOn_of_strictDescent_and_hasUnblockedInteriorEndpoint}
and
\path{BoundaryFreeIncidentFaceSelector.theorem49BoundaryRootNonemptyProjectedSynthesis_of_strictDescent_and_hasUnblockedInteriorEndpoint}
consume the endpoint witness directly, while
\path{BoundaryFreeIncidentFaceSelector.exists_terminal_face_boundary_remainders_of_strictDescent_and_hasUnblockedInteriorEndpoint}
exposes the terminal boundary-break face from the same witness.  The small
bridge
\path{interiorEdgeSupport_nonempty_of_hasUnblockedInteriorEndpoint}
is now the only conversion needed to use selector terminal-break theorems that
expect live interior-edge support.  The sharpened selector diagnostic
\path{BoundaryFreeIncidentFaceSelector.exists_terminal_face_boundary_remainders_and_disjoint_of_strictDescent_and_hasUnblockedInteriorEndpoint}
keeps the two pieces on the same selected face: every non-witness remainder is
ambient-boundary support, and the whole face boundary is still disjoint from
selected-boundary support.  The same
forcing obstruction has now been calibrated on the wheel-with-inner-triangle
source benchmark as well.  In
\path{Theorem49ForcingInteriorEdgeObstructionRegression.lean},
\path{not_nonempty_boundaryFreeIncidentFaceSelector_wheelWithInnerTriangle}
shows that the honest-source, Tait-colorable wheel model with nonempty purified
carrier cannot even carry a boundary-free incident-face selector; this now
follows through the direct selector/forcing duality rather than a separate
face-choice argument.  The theorem
\path{wheelWithInnerTriangle_closedWalkSource_tait_hasUnblockedEndpoint_without_interiorDualBoundaryRootAdjDistancePeelData}
puts the same obstruction in the exact raw source-IDBRAD bridge shape: source,
coloring, and `HasUnblockedInteriorEndpoint' coexist, but the generic
interior-dual datum does not.  The theorem
\path{not_forall_some_boundaryFreeSelector_or_acyclicEndpoint_of_closedWalkSource_tait_carrier_wheel}
then refutes the source-level universal claim that honest closed-walk annulus
source data, a Tait coloring, and a nonempty purified carrier force either such
a selector or any of the current acyclic annulus endpoints: well-founded
witness-cover data, height-ordered witness-cover data, finite collar-layer
data, or weak annulus-collar geometry.
The obstruction file also exposes the universal
failed-converse form
\path{not_forall_nonempty_planarBoundaryAnnulusConstructionBoundarySupportFaceData_of_exists_embedding_with_property_and_forcingInteriorEdgeWitness}
and its honest-source specialization
\path{not_forall_nonempty_planarBoundaryAnnulusConstructionBoundarySupportFaceData_of_exists_embedding_closedWalkAnnulusBoundarySource_and_forcingInteriorEdgeWitness},
so a single same-embedding forcing edge blocks any theorem deriving the weakest
peeled-interior construction shell from the source property alone.  The same
abstract witness now also covers the
weaker paired source shell
\path{annulusBoundaryReachability_and_selectedBoundaryArcGeometry_does_not_imply_interiorDualBoundaryRootAdjDistancePeelData}
and its construction-side companion
\path{annulusBoundaryReachability_and_selectedBoundaryArcGeometry_does_not_imply_planarBoundaryAnnulusConstructionBoundarySupportFaceData},
as well as the reduced construction companion
\path{annulusBoundaryReachability_and_selectedBoundaryArcGeometry_does_not_imply_planarBoundaryAnnulusConstructionPositiveFrontierData},
so the selected-boundary-arc weakening on the diamond witness is no longer a
separate local argument.  The accompanying regression file
\path{Theorem49ForcingInteriorEdgeObstructionRegression.lean} reconstructs the
earlier diamond-with-triangle failures, including the selected-boundary-arc
weakening, the reduced positive-frontier shell, and the concrete
\path{PlanarBoundaryAnnulusConstructionFaceLayerData} shell that feeds the
annulus decomposition route, from this abstract witness alone.  The same file
now also lifts these statements one level up to
graph-level same-embedding existential form: on \path{diamondWithTriangleGraph}
there exists an embedding carrying the honest source shell or the weaker
annulus-rooted shells together with a forcing interior-edge witness, but no
embedding of that same graph can carry the witness together with the peeled
interior-dual package, the reduced positive-frontier shell, or the concrete
construction face-layer shell.  The obstruction file now also states the
generic exact-shell converse failures
\path{not_forall_nonempty_planarBoundaryAnnulusConstructionFaceLayerData_of_exists_embedding_closedWalkAnnulusBoundarySource_and_oneCollarAnnulusCurrentBoundaryData_and_taitEdgeColoring_and_hasUnblockedInteriorEndpoint_and_v23ResidualBoundaryInitialState_and_forcingInteriorEdgeWitness}
and
\path{not_forall_nonempty_planarBoundaryAnnulusConstructionFaceLayerData_of_exists_embedding_boundaryReachabilityData_and_dartSuccessorCycleEmbeddingData_and_selectedBoundaryArc_and_oneCollarAnnulusCurrentBoundaryData_and_taitEdgeColoring_and_hasUnblockedInteriorEndpoint_and_v23ResidualBoundaryInitialState_and_forcingInteriorEdgeWitness}:
any explicit same-embedding exact one-collar/v23 witness carrying a forcing
interior edge already refutes a universal derivation of the concrete
construction face-layer shell.  The shared-interior-pair exact-shell failures
are therefore one concrete instance of a generic route-level false burden.
The obstruction file also states the
exact local burden already on the one-collar v23/v23.5 shell:
\path{not_forall_not_nonempty_forcingInteriorEdgeWitness_of_closedWalkAnnulusBoundarySource_and_oneCollarAnnulusCurrentBoundaryData_and_v23ResidualBoundaryInitialState_and_taitEdgeColoring_and_hasUnblockedInteriorEndpoint_sharedInteriorPair}
and the successor-cycle counterpart
\path{not_forall_not_nonempty_forcingInteriorEdgeWitness_of_boundaryReachabilityData_and_dartSuccessorCycleEmbeddingData_and_selectedBoundaryArc_and_oneCollarAnnulusCurrentBoundaryData_and_v23ResidualBoundaryInitialState_and_taitEdgeColoring_and_hasUnblockedInteriorEndpoint_sharedInteriorPair}
show that even this exact one-collar source package still does not universally
exclude forcing interior edges.  So any sound repair on that shell must add or
derive a precise local no-forcing hypothesis before it can reach the
downstream construction surface.  The obstruction file now also states the
forward repaired burden directly on that downstream shell:
\path{everyInteriorEdgeHasBoundaryFreeIncidentFace_of_planarBoundaryAnnulusConstructionFaceLayerData},
\path{boundaryFreeIncidentFaceSelector_of_planarBoundaryAnnulusConstructionFaceLayerData}, and
\path{not_nonempty_forcingInteriorEdgeWitness_of_planarBoundaryAnnulusConstructionFaceLayerData}
show that any theorem reaching
\path{PlanarBoundaryAnnulusConstructionFaceLayerData} must already build the
local boundary-free selector, equivalently exclude forcing interior edges,
before any later annulus-decomposition lowering.
The repaired positive direction is now stated on that exact shell too.
\path{Theorem49ResidualBoundaryRoute.lean} proves
\path{theorem49BoundaryRootNonemptyProjectedSynthesis_of_closedWalkAnnulusBoundarySource_and_oneCollarAnnulusCurrentBoundaryData_and_v23ResidualBoundaryInitialState_and_planarBoundaryAnnulusRootParentPeelData_and_hasUnblockedInteriorEndpoint}
and the actual boundary-order counterpart
\path{theorem49BoundaryRootNonemptyProjectedSynthesis_of_boundaryReachabilityData_and_dartSuccessorCycleEmbeddingData_and_selectedBoundaryArc_and_oneCollarAnnulusCurrentBoundaryData_and_v23ResidualBoundaryInitialState_and_planarBoundaryAnnulusRootParentPeelData_and_hasUnblockedInteriorEndpoint}.
So the current strongest sound v23/v23.5 repair is explicit: if the exact
one-collar shell already carries annulus-root parent peel data on the same
annulus boundary split, then the local endpoint witness closes the route all
the way to \path{Theorem49BoundaryRootNonemptyProjectedSynthesis}.  The
remaining geometric burden is therefore the same-split parent-peel
certificate itself, not additional projection symmetry.
The constructive selector/construction lane now reaches one real shell beyond
that point.  \path{Theorem49BoundaryFreeSelectorConstruction.lean} proves
\path{planarBoundaryAnnulusConstructionBoundarySupportFaceData_of_boundaryReachabilityData_and_closedWalkEmbeddingData_and_selectedBoundaryArcOnFace_and_planarBoundaryAnnulusRootParentPeelData_and_selectorDeficitSelectedBoundary}
and the successor-cycle variant
\path{planarBoundaryAnnulusConstructionBoundarySupportFaceData_of_boundaryReachabilityData_and_dartSuccessorCycleEmbeddingData_and_selectedBoundaryArc_and_planarBoundaryAnnulusRootParentPeelData_and_selectorDeficitSelectedBoundary}.
So same-split annulus-root parent peel data already lowers to the real
\path{PlanarBoundaryAnnulusConstructionBoundarySupportFaceData} shell once the
exact local gap is discharged: peeled faces missed by the selector image must
still meet selected-boundary support, and the positive
\path{currentOuterBoundary} layers of the resulting construction must be
nonempty and pairwise disjoint.  This identifies the remaining
parent-peel-to-construction burden precisely.  It is not a generic
non-peeled-face condition, but the smaller selector-deficit peeled-face
requirement.  Lean now also proves
\path{selectorDeficitSelectedBoundary_of_planarBoundaryAnnulusRootParentPeelData_and_nonPeelSelectedBoundary},
so whenever an older route already gives the stronger hypothesis that every
non-peeled face of the same parent-peel construction touches
selected-boundary support, this repaired selector-deficit premise follows
immediately.  The same exact one-collar/v23 shell now also has generic
forcing-edge converse failures already at that stronger construction surface
and the next two intermediate construction shells:
\path{not_forall_nonempty_planarBoundaryAnnulusConstructionBoundarySupportFaceData_of_exists_embedding_closedWalkAnnulusBoundarySource_and_oneCollarAnnulusCurrentBoundaryData_and_taitEdgeColoring_and_hasUnblockedInteriorEndpoint_and_v23ResidualBoundaryInitialState_and_forcingInteriorEdgeWitness}
and
\path{not_forall_nonempty_planarBoundaryAnnulusConstructionBoundarySupportFaceData_of_exists_embedding_boundaryReachabilityData_and_dartSuccessorCycleEmbeddingData_and_selectedBoundaryArc_and_oneCollarAnnulusCurrentBoundaryData_and_taitEdgeColoring_and_hasUnblockedInteriorEndpoint_and_v23ResidualBoundaryInitialState_and_forcingInteriorEdgeWitness}.
\path{not_forall_nonempty_planarBoundaryAnnulusConstructionFacePartitionData_of_exists_embedding_closedWalkAnnulusBoundarySource_and_oneCollarAnnulusCurrentBoundaryData_and_taitEdgeColoring_and_hasUnblockedInteriorEndpoint_and_v23ResidualBoundaryInitialState_and_forcingInteriorEdgeWitness},
\path{not_forall_nonempty_planarBoundaryAnnulusConstructionFacePartitionData_of_exists_embedding_boundaryReachabilityData_and_dartSuccessorCycleEmbeddingData_and_selectedBoundaryArc_and_oneCollarAnnulusCurrentBoundaryData_and_taitEdgeColoring_and_hasUnblockedInteriorEndpoint_and_v23ResidualBoundaryInitialState_and_forcingInteriorEdgeWitness},
\path{not_forall_nonempty_planarBoundaryAnnulusConstructionPositiveFrontierData_of_exists_embedding_closedWalkAnnulusBoundarySource_and_oneCollarAnnulusCurrentBoundaryData_and_taitEdgeColoring_and_hasUnblockedInteriorEndpoint_and_v23ResidualBoundaryInitialState_and_forcingInteriorEdgeWitness},
and
\path{not_forall_nonempty_planarBoundaryAnnulusConstructionPositiveFrontierData_of_exists_embedding_boundaryReachabilityData_and_dartSuccessorCycleEmbeddingData_and_selectedBoundaryArc_and_oneCollarAnnulusCurrentBoundaryData_and_taitEdgeColoring_and_hasUnblockedInteriorEndpoint_and_v23ResidualBoundaryInitialState_and_forcingInteriorEdgeWitness}
now show the same obstruction one layer later inside the construction lane.
So any explicit same-embedding exact-shell witness carrying a forcing interior
edge already refutes universal derivations of
\path{PlanarBoundaryAnnulusConstructionBoundarySupportFaceData},
\path{PlanarBoundaryAnnulusConstructionFacePartitionData}, and
\path{PlanarBoundaryAnnulusConstructionPositiveFrontierData}; the live gap
therefore persists before the final lowering to face-layer data.
That same repaired exact-shell package now reaches the downstream
construction face-layer shell used by annulus decomposition.
\path{Theorem49ResidualBoundaryRoute.lean} proves
\path{nonempty_planarBoundaryAnnulusConstructionFaceLayerData_of_closedWalkAnnulusBoundarySource_and_oneCollarAnnulusCurrentBoundaryData_and_v23ResidualBoundaryInitialState_and_planarBoundaryAnnulusRootParentPeelData_and_selectorDeficitSelectedBoundary}
and
\path{nonempty_planarBoundaryAnnulusConstructionFaceLayerData_of_boundaryReachabilityData_and_dartSuccessorCycleEmbeddingData_and_selectedBoundaryArc_and_oneCollarAnnulusCurrentBoundaryData_and_v23ResidualBoundaryInitialState_and_planarBoundaryAnnulusRootParentPeelData_and_selectorDeficitSelectedBoundary}.
So once the exact one-collar shell carries same-split parent peel data
together with the selector-deficit selected-boundary condition and the
positive \path{currentOuterBoundary} side conditions, the route no longer
stops at selected-boundary support data: it already produces real
\path{PlanarBoundaryAnnulusConstructionFaceLayerData} on that same embedding.
So the remaining issue is not any further lowering below parent-peel itself,
but whether the current exact shell forces that package at all.
That same repaired exact-shell package now also reaches a real annulus
decomposition on the same embedding.
\path{Theorem49ResidualBoundaryRoute.lean} proves
\path{nonempty_planarBoundaryAnnulusDecomposition_of_closedWalkAnnulusBoundarySource_and_oneCollarAnnulusCurrentBoundaryData_and_v23ResidualBoundaryInitialState_and_planarBoundaryAnnulusRootParentPeelData_and_selectorDeficitSelectedBoundary}
and
\path{nonempty_planarBoundaryAnnulusDecomposition_of_boundaryReachabilityData_and_dartSuccessorCycleEmbeddingData_and_selectedBoundaryArc_and_oneCollarAnnulusCurrentBoundaryData_and_v23ResidualBoundaryInitialState_and_planarBoundaryAnnulusRootParentPeelData_and_selectorDeficitSelectedBoundary}.
So once the exact shell carries the same-split parent-peel and selector-deficit
package, the route already passes through annulus decomposition itself.  The
remaining burden is therefore not another decomposition-side shell, but
whatever stronger geometry or witness data is needed downstream of that
repaired local package.  Lean now also proves
\path{nonempty_planarBoundaryAnnulusDecomposition_of_planarBoundaryAnnulusConstructionFaceLayerData},
\path{theorem49BoundaryRootSynthesis_of_planarBoundaryAnnulusConstructionFaceLayerData_and_resultingDecompositionWitnessAssignment},
\path{theorem49BoundaryRootSynthesis_of_planarBoundaryAnnulusConstructionFaceLayerData_and_inducedDecompositionWitnessAssignment},
\path{nonempty_planarBoundaryAnnulusDecomposition_of_planarBoundaryAnnulusConstructionBoundarySupportFaceData},
and
\path{theorem49BoundaryRootSynthesis_of_planarBoundaryAnnulusConstructionBoundarySupportFaceData_and_inducedDecompositionWitnessAssignment},
\path{theorem49BoundaryRootSynthesis_of_planarBoundaryAnnulusConstructionBoundarySupportFaceData_and_resultingDecompositionWitnessAssignment},
so once the route reaches \path{PlanarBoundaryAnnulusConstructionFaceLayerData},
this decomposition step is automatic and the full Theorem~4.9 endpoint reduces
exactly to witness ownership on the exact annulus decomposition canonically
induced by that shell.  The stronger
selected-boundary \path{PlanarBoundaryAnnulusConstructionBoundarySupportFaceData}
shell is now only an upstream route into that face-layer package, not part of
the downstream closure theorem itself.  The same file now also proves the
graph-level route wrappers
\path{exists_theorem49BoundaryRootSynthesis_of_exists_closedWalkAnnulusBoundarySource_and_planarBoundaryAnnulusConstructionFaceLayerData_and_resultingDecompositionWitnessAssignment_direct},
the stronger selected-boundary-support analogue, and the matching
successor-cycle pair.  So once the honest-source or successor-cycle route
reaches either same-embedding construction shell and the resulting
decomposition carries witness ownership, full
\path{Theorem49BoundaryRootSynthesis} already follows directly at graph level;
no further repackaging step remains below that shell.
\path{Theorem49PositiveGeometricSourceReplacement.lean}
now also proves
\path{theorem49BoundaryRootNonemptyProjectedSynthesis_of_planarBoundaryAnnulusConstructionFaceLayerData},
\path{theorem49BoundaryRawQuotientErrorPackage_of_planarBoundaryAnnulusConstructionFaceLayerData},
\path{theorem49BoundaryRootNonemptyProjectedSynthesis_of_planarBoundaryAnnulusConstructionBoundarySupportFaceData},
and
\path{theorem49BoundaryRawQuotientErrorPackage_of_planarBoundaryAnnulusConstructionBoundarySupportFaceData},
so once either construction shell is paired with a surviving purified carrier
on the same embedding, the current projected endpoint and its raw
quotient/error package are automatic as well.  The remaining burden above
those shells is therefore to preserve the purified carrier, or else derive
stronger earlier geometry, not to find another replacement-side lowering.
\path{Theorem49PositiveGeometricSourceReplacementRoute.lean} now also proves
the graph-package methods
\path{Theorem49ClosedWalkAnnulusConstructionFaceLayerPositiveProjectedGeometry.boundaryRootSynthesis},
\path{Theorem49ClosedWalkAnnulusConstructionFaceLayerPositiveProjectedGeometry.exists_boundaryRootSynthesis},
and the matching successor-cycle pair.  So once those named
construction-face-layer route packages themselves carry witness ownership on
the same derived annulus decomposition, full
\path{Theorem49BoundaryRootSynthesis} follows immediately; there is no
remaining repackaging gap between those graph-level route objects and the
Theorem~4.9 endpoint.
\path{Theorem49PositiveGeometricSourceReplacementRoute.lean} now also packages
the stronger selected-boundary-contact shell itself as the fixed-embedding predicates
\path{Theorem49ClosedWalkAnnulusConstructionBoundarySupportPositiveProjectedGeometryOn}
and
\path{Theorem49SuccessorCycleAnnulusConstructionBoundarySupportPositiveProjectedGeometryOn},
together with the graph-level structures
\path{Theorem49ClosedWalkAnnulusConstructionBoundarySupportPositiveProjectedGeometry}
and
\path{Theorem49SuccessorCycleAnnulusConstructionBoundarySupportPositiveProjectedGeometry}.
Those graph packages already expose \path{boundaryRootSynthesis} and
\path{exists_boundaryRootSynthesis} on the induced annulus decomposition, so
callers no longer need to first repackage this shell as construction-face-layer
route data just to state the direct Theorem~4.9 closure.  The same file still
packages the derived face-layer closure as the fixed-embedding predicates
\path{Theorem49ClosedWalkAnnulusConstructionFaceLayerPositiveProjectedGeometryOn}
and
\path{Theorem49SuccessorCycleAnnulusConstructionFaceLayerPositiveProjectedGeometryOn},
together with the graph-level structures
\path{Theorem49ClosedWalkAnnulusConstructionFaceLayerPositiveProjectedGeometry}
and
\path{Theorem49SuccessorCycleAnnulusConstructionFaceLayerPositiveProjectedGeometry}.
The stronger selected-boundary-contact shell still lowers directly into those
derived face-layer route surfaces via
\path{PlanarBoundaryAnnulusConstructionBoundarySupportFaceData.toPlanarBoundaryAnnulusConstructionFaceLayerData}.
The same file now also proves
\path{not_theorem49ClosedWalkAnnulusConstructionFaceLayerPositiveProjectedGeometryOn_of_facewiseAtMostOneInteriorEdge},
\path{not_theorem49SuccessorCycleAnnulusConstructionFaceLayerPositiveProjectedGeometryOn_of_facewiseAtMostOneInteriorEdge},
and
\path{exists_two_distinct_interior_edges_on_faceBoundary_of_successorCycleAnnulusConstructionFaceLayerPositiveProjectedGeometryOn}.
So the stronger source-facing construction face-layer shell itself already sits
outside the canonical at-most-one branch, and the successor-cycle shell now
directly exposes the same local two-interior-edge obstruction there too.
So once the honest-source or successor-cycle route reaches same-embedding
\path{PlanarBoundaryAnnulusConstructionFaceLayerData} and still carries
\path{HasUnblockedInteriorEndpoint}, the current projected Theorem~4.9
endpoint, the raw quotient/error package, and the replacement-positive
height/collar packages already follow directly, without any further
collar-geometry or well-founded-witness detour.
That diagnosis is now repaired again.
\path{Theorem49BoundaryFreeSelectorConstruction.lean} proves
\path{parentWitnessOrientation_of_boundaryData_and_interiorDualBoundaryRootParentPeelData},
\path{parentWitnessOrientation_of_planarBoundaryAnnulusRootParentPeelData}, and
\path{not_exists_planarBoundaryAnnulusConstructionBoundarySupportFaceData_with_planarBoundaryAnnulusRootParentPeelConstruction}.
So the exact selector-driven parent-peel construction is already
\path{PlanarBoundaryAnnulusConstruction.ParentWitnessOrientation} on the same
embedding, while any same-construction selected-boundary-support shell still
forbids parent orientation.  The same construction file now isolates the
resulting contradiction at the actual shell that creates it:
\path{false_of_boundaryReachabilityData_and_closedWalkEmbeddingData_and_selectedBoundaryArcOnFace_and_planarBoundaryAnnulusRootParentPeelData_and_selectorDeficitSelectedBoundary_and_currentOuterBoundary}
and
\path{false_of_boundaryReachabilityData_and_dartSuccessorCycleEmbeddingData_and_selectedBoundaryArc_and_planarBoundaryAnnulusRootParentPeelData_and_selectorDeficitSelectedBoundary_and_currentOuterBoundary}.
The same construction file now also packages this as raw boundary-order graph-level
nonexistence theorems:
\path{not_exists_boundaryReachabilityData_and_closedWalkEmbeddingData_and_selectedBoundaryArcOnFace_and_planarBoundaryAnnulusRootParentPeelData_and_selectorDeficitSelectedBoundary_and_currentOuterBoundary}
and
\path{not_exists_boundaryReachabilityData_and_dartSuccessorCycleEmbeddingData_and_selectedBoundaryArc_and_planarBoundaryAnnulusRootParentPeelData_and_selectorDeficitSelectedBoundary_and_currentOuterBoundary}.
It also now gives the fixed-embedding local nonexistence forms
\path{not_exists_planarBoundaryAnnulusRootParentPeelData_and_selectorDeficitSelectedBoundary_and_currentOuterBoundary_of_boundaryReachabilityData_and_closedWalkEmbeddingData_and_selectedBoundaryArcOnFace}
and
\path{not_exists_planarBoundaryAnnulusRootParentPeelData_and_selectorDeficitSelectedBoundary_and_currentOuterBoundary_of_boundaryReachabilityData_and_dartSuccessorCycleEmbeddingData_and_selectedBoundaryArc},
which are the direct reusable statements downstream route arguments should consume.
So the repaired selector-deficit plus positive-current-boundary package is already globally
impossible before any exact one-collar or exact-v23 residual hypotheses are added.
\path{Theorem49ResidualBoundaryRoute.lean} then re-exports this at the repaired
exact one-collar shell itself with
\path{false_of_closedWalkAnnulusBoundarySource_and_oneCollarAnnulusCurrentBoundaryData_and_v23ResidualBoundaryInitialState_and_planarBoundaryAnnulusRootParentPeelData_and_selectorDeficitSelectedBoundary_and_currentOuterBoundary}
and
\path{false_of_boundaryReachabilityData_and_dartSuccessorCycleEmbeddingData_and_selectedBoundaryArc_and_oneCollarAnnulusCurrentBoundaryData_and_v23ResidualBoundaryInitialState_and_planarBoundaryAnnulusRootParentPeelData_and_selectorDeficitSelectedBoundary_and_currentOuterBoundary}.
So the old positive \path{currentOuterBoundary} side conditions are not merely
another local burden to derive on the same parent-peel construction: they are
incompatible with that construction.  The earlier exact-shell face-layer,
decomposition, and synthesis theorems with those hypotheses remain valid but
vacuous.  The same route file now also states that the stronger older
non-peeled-face selected-boundary-contact repair collapses for the same
reason:
\path{false_of_closedWalkAnnulusBoundarySource_and_oneCollarAnnulusCurrentBoundaryData_and_v23ResidualBoundaryInitialState_and_planarBoundaryAnnulusRootParentPeelData_and_nonPeelSelectedBoundary_and_currentOuterBoundary}
and
\path{false_of_boundaryReachabilityData_and_dartSuccessorCycleEmbeddingData_and_selectedBoundaryArc_and_oneCollarAnnulusCurrentBoundaryData_and_v23ResidualBoundaryInitialState_and_planarBoundaryAnnulusRootParentPeelData_and_nonPeelSelectedBoundary_and_currentOuterBoundary}.
So the stronger \path{nonPeelSelectedBoundary} branch is not merely reducible
to the selector-deficit branch inside later synthesis theorems; it is already
directly inconsistent with the same old positive current-boundary side
conditions.  \path{Theorem49ResidualBoundaryRoute.lean} now also packages this as
direct graph-level nonexistence theorems:
\path{not_exists_closedWalkAnnulusBoundarySource_and_oneCollarAnnulusCurrentBoundaryData_and_v23ResidualBoundaryInitialState_and_planarBoundaryAnnulusRootParentPeelData_and_selectorDeficitSelectedBoundary_and_currentOuterBoundary}
and
\path{not_exists_boundaryReachabilityData_and_dartSuccessorCycleEmbeddingData_and_selectedBoundaryArc_and_oneCollarAnnulusCurrentBoundaryData_and_v23ResidualBoundaryInitialState_and_planarBoundaryAnnulusRootParentPeelData_and_selectorDeficitSelectedBoundary_and_currentOuterBoundary}.
So this repaired selector-deficit plus positive-current-boundary shell is not
merely unforced by the live route: even the stricter exact-shell witness form
is globally impossible, as an immediate corollary of the boundary-order
contradiction above.  The matching stronger non-peeled-face graph-level
nonexistence theorems are now present too:
\path{not_exists_closedWalkAnnulusBoundarySource_and_oneCollarAnnulusCurrentBoundaryData_and_v23ResidualBoundaryInitialState_and_planarBoundaryAnnulusRootParentPeelData_and_nonPeelSelectedBoundary_and_currentOuterBoundary}
and
\path{not_exists_boundaryReachabilityData_and_dartSuccessorCycleEmbeddingData_and_selectedBoundaryArc_and_oneCollarAnnulusCurrentBoundaryData_and_v23ResidualBoundaryInitialState_and_planarBoundaryAnnulusRootParentPeelData_and_nonPeelSelectedBoundary_and_currentOuterBoundary}.
So even that stronger repair branch is globally impossible on the exact shell.
\path{Theorem49ResidualBoundaryObstruction.lean} now upgrades this to
direct exact-shell failed universals:
\path{not_forall_exists_planarBoundaryAnnulusRootParentPeelData_and_selectorDeficitSelectedBoundary_and_currentOuterBoundary_of_closedWalkAnnulusBoundarySource_and_oneCollarAnnulusCurrentBoundaryData_and_v23ResidualBoundaryInitialState_and_taitEdgeColoring_and_hasUnblockedInteriorEndpoint_sharedInteriorPair}
and
\path{not_forall_exists_planarBoundaryAnnulusRootParentPeelData_and_selectorDeficitSelectedBoundary_and_currentOuterBoundary_of_boundaryReachabilityData_and_dartSuccessorCycleEmbeddingData_and_selectedBoundaryArc_and_oneCollarAnnulusCurrentBoundaryData_and_v23ResidualBoundaryInitialState_and_taitEdgeColoring_and_hasUnblockedInteriorEndpoint_sharedInteriorPair},
with the matching forcing-edge abstract converse failures.  It now also exposes
the exact-shell counterexamples themselves as
\path{exists_embedding_closedWalkAnnulusBoundarySource_and_oneCollarAnnulusCurrentBoundaryData_and_taitEdgeColoring_and_hasUnblockedInteriorEndpoint_and_v23ResidualBoundaryInitialState_without_planarBoundaryAnnulusRootParentPeelData_and_selectorDeficitSelectedBoundary_and_currentOuterBoundary_sharedInteriorPair}
and
\path{exists_embedding_boundaryReachabilityData_and_dartSuccessorCycleEmbeddingData_and_selectedBoundaryArc_and_oneCollarAnnulusCurrentBoundaryData_and_taitEdgeColoring_and_hasUnblockedInteriorEndpoint_and_v23ResidualBoundaryInitialState_without_planarBoundaryAnnulusRootParentPeelData_and_selectorDeficitSelectedBoundary_and_currentOuterBoundary_sharedInteriorPair}.
So even the full repaired selector-deficit plus positive-current-boundary
package is not forced by the live exact shell on the shared-interior-pair
benchmark.  The same obstruction file now upgrades the stronger older
non-peeled-face branch in the same way:
\path{not_forall_exists_planarBoundaryAnnulusRootParentPeelData_and_nonPeelSelectedBoundary_and_currentOuterBoundary_of_closedWalkAnnulusBoundarySource_and_oneCollarAnnulusCurrentBoundaryData_and_v23ResidualBoundaryInitialState_and_taitEdgeColoring_and_hasUnblockedInteriorEndpoint_sharedInteriorPair}
and
\path{not_forall_exists_planarBoundaryAnnulusRootParentPeelData_and_nonPeelSelectedBoundary_and_currentOuterBoundary_of_boundaryReachabilityData_and_dartSuccessorCycleEmbeddingData_and_selectedBoundaryArc_and_oneCollarAnnulusCurrentBoundaryData_and_v23ResidualBoundaryInitialState_and_taitEdgeColoring_and_hasUnblockedInteriorEndpoint_sharedInteriorPair},
with the matching exact-shell counterexamples
\path{exists_embedding_closedWalkAnnulusBoundarySource_and_oneCollarAnnulusCurrentBoundaryData_and_taitEdgeColoring_and_hasUnblockedInteriorEndpoint_and_v23ResidualBoundaryInitialState_without_planarBoundaryAnnulusRootParentPeelData_and_nonPeelSelectedBoundary_and_currentOuterBoundary_sharedInteriorPair}
and
\path{exists_embedding_boundaryReachabilityData_and_dartSuccessorCycleEmbeddingData_and_selectedBoundaryArc_and_oneCollarAnnulusCurrentBoundaryData_and_taitEdgeColoring_and_hasUnblockedInteriorEndpoint_and_v23ResidualBoundaryInitialState_without_planarBoundaryAnnulusRootParentPeelData_and_nonPeelSelectedBoundary_and_currentOuterBoundary_sharedInteriorPair}.
So even that stronger older repair is not forced by the live exact shell on the
same benchmark.  The live route must therefore avoid that positive-frontier
shell entirely and continue through the already parent-oriented construction /
well-founded-parent lane.
\path{Theorem49ResidualBoundaryRoute.lean} now also
proves the direct exact-shell endpoint theorems
\path{theorem49BoundaryRootSynthesis_of_closedWalkAnnulusBoundarySource_and_oneCollarAnnulusCurrentBoundaryData_and_v23ResidualBoundaryInitialState_and_planarBoundaryAnnulusRootParentPeelData}
and
\path{theorem49BoundaryRootSynthesis_of_boundaryReachabilityData_and_dartSuccessorCycleEmbeddingData_and_selectedBoundaryArc_and_oneCollarAnnulusCurrentBoundaryData_and_v23ResidualBoundaryInitialState_and_planarBoundaryAnnulusRootParentPeelData}.
So once the exact one-collar shell carries same-split annulus-root parent peel
data at all, the full Theorem~4.9 synthesis endpoint is already immediate on
that embedding; the selector-deficit and positive-current-boundary detour is
not just unnecessary but inconsistent on that same construction.  The live
burden is therefore sharper once again: derive the actual
\path{PlanarBoundaryAnnulusRootParentPeelData} package from the residual shell,
rather than any construction-shell refinement beyond it.  The same file now
also proves the even earlier exact-shell theorems
\path{theorem49BoundaryRootSynthesis_of_closedWalkAnnulusBoundarySource_and_oneCollarAnnulusCurrentBoundaryData_and_v23ResidualBoundaryInitialState_and_boundaryFaceRootsCanonicalParentSharedEdgeCover}
and
\path{theorem49BoundaryRootSynthesis_of_boundaryReachabilityData_and_dartSuccessorCycleEmbeddingData_and_selectedBoundaryArc_and_oneCollarAnnulusCurrentBoundaryData_and_v23ResidualBoundaryInitialState_and_boundaryFaceRootsCanonicalParentSharedEdgeCover},
together with the matching graph-level direct existence forms.  So once the
exact one-collar shell already carries the raw source-fixed canonical-parent
shared-edge cover on the extracted boundary-face roots, full
\path{Theorem49BoundaryRootSynthesis} is already immediate on that embedding,
with no \path{HasUnblockedInteriorEndpoint} and no separate parent-peel
repackaging step.  The live exact-shell burden is therefore sharper again:
force that concrete source-fixed shared-edge cover itself, or prove that the
repaired shell cannot support it.  The same file now also proves the exact-
shell selector-side lowerings
\path{theorem49HeightOrderedPositiveProjectedGeometryOn_of_closedWalkAnnulusBoundarySource_and_oneCollarAnnulusCurrentBoundaryData_and_v23ResidualBoundaryInitialState_and_boundaryFaceRootsCanonicalParentSharedEdgeCover_via_boundaryFreeSelector}
and
\path{theorem49HeightOrderedPositiveProjectedGeometryOn_of_boundaryReachabilityData_and_dartSuccessorCycleEmbeddingData_and_selectedBoundaryArc_and_oneCollarAnnulusCurrentBoundaryData_and_v23ResidualBoundaryInitialState_and_boundaryFaceRootsCanonicalParentSharedEdgeCover_via_boundaryFreeSelector},
together with the matching graph-level existential forms.  So once that same
exact shell carries the raw source-fixed canonical-parent shared-edge cover
together with \path{HasUnblockedInteriorEndpoint}, it already lands on the
named \path{Theorem49HeightOrderedPositiveProjectedGeometryOn} surface before
any residual-boundary repackaging.  The same file now also lowers the sharper
rooted-cover shell plus \path{HasUnblockedInteriorEndpoint} to the named
closed-walk and successor-cycle root-distance positive route surfaces, and
from there to the projected nonempty and raw quotient/error theorem-4.9
endpoints on the same embedding.  So the exact one-collar/v23 rooted-cover
shell now enters the modular replacement-positive pipeline itself, not only as
a direct full-synthesis shortcut.  The same file now also proves the direct
exact-shell height-ordered lowerings
\path{theorem49HeightOrderedPositiveProjectedGeometryOn_of_closedWalkAnnulusBoundarySource_and_oneCollarAnnulusCurrentBoundaryData_and_v23ResidualBoundaryInitialState_and_boundaryFaceRootsCanonicalRootedSharedEdgeFaceMembershipCover_and_hasUnblockedInteriorEndpoint}
and
\path{theorem49HeightOrderedPositiveProjectedGeometryOn_of_boundaryReachabilityData_and_dartSuccessorCycleEmbeddingData_and_selectedBoundaryArc_and_oneCollarAnnulusCurrentBoundaryData_and_v23ResidualBoundaryInitialState_and_boundaryFaceRootsCanonicalRootedSharedEdgeFaceMembershipCover_and_hasUnblockedInteriorEndpoint},
together with the matching graph-level existential forms.  So once that same
exact shell carries the concrete rooted shared-edge face-membership cover
together with \path{HasUnblockedInteriorEndpoint}, it already lands on the
named \path{Theorem49HeightOrderedPositiveProjectedGeometryOn} surface itself,
not just on the intermediate annulus root-distance positive package.  The same
file now also proves the direct exact-shell finite-collar lowerings
\path{theorem49CollarLayerPositiveProjectedGeometryOn_of_closedWalkAnnulusBoundarySource_and_oneCollarAnnulusCurrentBoundaryData_and_v23ResidualBoundaryInitialState_and_boundaryFaceRootsCanonicalRootedSharedEdgeFaceMembershipCover_and_hasUnblockedInteriorEndpoint}
and
\path{theorem49CollarLayerPositiveProjectedGeometryOn_of_boundaryReachabilityData_and_dartSuccessorCycleEmbeddingData_and_selectedBoundaryArc_and_oneCollarAnnulusCurrentBoundaryData_and_v23ResidualBoundaryInitialState_and_boundaryFaceRootsCanonicalRootedSharedEdgeFaceMembershipCover_and_hasUnblockedInteriorEndpoint},
together with the matching graph-level existential forms.  So that same exact
shell now also lands directly on the named
\path{Theorem49CollarLayerPositiveProjectedGeometryOn} surface itself, not
only after passing through the height-ordered route surface or the later
residual-boundary positive interface.  The same
file now also lowers those named
root-distance positive route surfaces to
\path{Theorem49ResidualBoundaryPositiveProjectedGeometryOn} and the associated
current-boundary remainder witness interface, so the exact one-collar/v23
rooted-cover shell reaches the residual-boundary positive lane by composition
too.  \path{Theorem49ResidualBoundaryPeeling.lean} now also proves the same
residual bridge directly from the rooted shared-edge face-membership cover
itself on both the honest closed-walk and successor-cycle shells, and
\path{Theorem49ResidualBoundaryRoute.lean} adds the matching exact
one-collar/v23 wrappers.  So the exact rooted-cover shell no longer needs an
explicit detour through a named root-distance package just to expose residual-
positive geometry or a current-boundary remainder witness; those surfaces are
available there directly as well.  The same route file now also packages these
direct exact-shell residual consequences at graph level: from existential
one-collar/v23 rooted-cover data plus \path{HasUnblockedInteriorEndpoint}, it
now directly derives existential residual-positive geometry and existential
current-boundary remainder witnesses on both the honest closed-walk and
successor-cycle shells.  So downstream route consumers no longer need to
rebuild that existential residual plumbing by hand once they start from the
live exact shell.  The same exact rooted-cover shell now also has matching
graph-level theorem-4.9 endpoint packaging: existential one-collar/v23 rooted-
cover data plus \path{HasUnblockedInteriorEndpoint} directly yields
existential \path{Theorem49BoundaryRootNonemptyProjectedSynthesis} and
existential \path{Theorem49BoundaryRawQuotientErrorPackage} on both the honest
closed-walk and successor-cycle shells.  So the exact rooted-cover lane now
exports both the residual interface and the theorem-4.9 endpoint interface
without further existential repackaging.  At the more generic residual layer-
face witness level, the peeling and route files now also expose the raw
quotient/error theorem-4.9 endpoint directly: same-embedding witness data plus
\path{HasUnblockedInteriorEndpoint} yields
\path{Theorem49BoundaryRawQuotientErrorPackage}, and the honest closed-walk /
successor-cycle existential shells now package that endpoint at graph level
too.  So the residual witness layer itself now exports the full theorem-4.9
endpoint pair, not only the root-synthesis and projected-nonempty half.  The
selector shell now matches that route-facing endpoint coverage too.  Starting
from honest closed-walk or successor-cycle source data on a fixed embedding,
same-embedding \path{ResidualBoundarySelectorData} now exports
\path{Theorem49BoundaryRootSynthesis},
\path{Theorem49BoundaryRootNonemptyProjectedSynthesis}, and
\path{Theorem49BoundaryRawQuotientErrorPackage}; and the route file now
packages all three at graph level under the corresponding existential closed-
walk / successor-cycle selector shells.  So downstream work can now use the
selector lane as a full theorem-4.9 endpoint interface rather than only as a
source of root synthesis.  The residual interface is now exposed at that same
boundary-order layer too.  Honest closed-walk or successor-cycle source data
paired with either same-embedding residual witness data or same-embedding
\path{ResidualBoundarySelectorData}, together with
\path{HasUnblockedInteriorEndpoint}, now directly exports
\path{Theorem49ResidualBoundaryPositiveProjectedGeometryOn} and the
corresponding current-boundary remainder witness on the same embedding; and the
route file packages both consequences at graph level for all four shells.  So
both the residual witness lane and the selector lane now expose not only
theorem-4.9 endpoints but also the upstream residual-positive interface used
elsewhere in the live route.  The same graph-level residual export is now available one
layer more generally too: once an exact one-collar/v23 shell already carries
same-embedding \path{InteriorDualBoundaryRootAdjDistancePeelData} together with
\path{HasUnblockedInteriorEndpoint}, the route file now directly packages the
same existential residual-positive geometry and current-boundary remainder
witnesses on both the honest closed-walk and successor-cycle shells.  So any
later proof of the generic adj-distance package on a live exact-shell witness
can enter the residual interface without further existential repackaging.  The
same file now also proves the
exact-shell theorems
\path{theorem49BoundaryRootSynthesis_of_closedWalkAnnulusBoundarySource_and_oneCollarAnnulusCurrentBoundaryData_and_v23ResidualBoundaryInitialState_and_interiorDualBoundaryRootAdjDistancePeelData}
and
\path{theorem49BoundaryRootSynthesis_of_boundaryReachabilityData_and_dartSuccessorCycleEmbeddingData_and_selectedBoundaryArc_and_oneCollarAnnulusCurrentBoundaryData_and_v23ResidualBoundaryInitialState_and_interiorDualBoundaryRootAdjDistancePeelData}.
So on the exact one-collar v23/v23.5 shell, even the annulus-root parent-peel
package is no longer the minimal positive burden: it already suffices to
derive generic \path{InteriorDualBoundaryRootAdjDistancePeelData} on the same
embedding, or any equivalent cycle-breaking parent witness strong enough to
construct it.  The same file now also proves the exact-shell route-surface
lowerings
\path{theorem49ClosedWalkAnnulusRootParentPositiveProjectedGeometryOn_of_closedWalkAnnulusBoundarySource_and_oneCollarAnnulusCurrentBoundaryData_and_v23ResidualBoundaryInitialState_and_planarBoundaryAnnulusRootParentPeelData_and_hasUnblockedInteriorEndpoint},
\path{theorem49SuccessorCycleAnnulusRootParentPositiveProjectedGeometryOn_of_boundaryReachabilityData_and_dartSuccessorCycleEmbeddingData_and_selectedBoundaryArc_and_oneCollarAnnulusCurrentBoundaryData_and_v23ResidualBoundaryInitialState_and_planarBoundaryAnnulusRootParentPeelData_and_hasUnblockedInteriorEndpoint},
\path{theorem49ClosedWalkAnnulusRootAdjDistancePositiveProjectedGeometryOn_of_closedWalkAnnulusBoundarySource_and_oneCollarAnnulusCurrentBoundaryData_and_v23ResidualBoundaryInitialState_and_interiorDualBoundaryRootAdjDistancePeelData_and_hasUnblockedInteriorEndpoint},
and
\path{theorem49SuccessorCycleAnnulusRootAdjDistancePositiveProjectedGeometryOn_of_boundaryReachabilityData_and_dartSuccessorCycleEmbeddingData_and_selectedBoundaryArc_and_oneCollarAnnulusCurrentBoundaryData_and_v23ResidualBoundaryInitialState_and_interiorDualBoundaryRootAdjDistancePeelData_and_hasUnblockedInteriorEndpoint}.
So the live exact shell no longer merely has endpoint corollaries: once it
carries the parent-peel or generic adj-distance package together with
\path{HasUnblockedInteriorEndpoint}, it already lands on the named
root-parent or root-distance positive route surfaces used elsewhere in the
replacement pipeline.  The same file now also packages the corresponding
graph-level existential forms through
\path{nonempty_theorem49ClosedWalkAnnulusRootParentPositiveProjectedGeometry_of_exists_closedWalkAnnulusBoundarySource_and_oneCollarAnnulusCurrentBoundaryData_and_v23ResidualBoundaryInitialState_and_planarBoundaryAnnulusRootParentPeelData_and_hasUnblockedInteriorEndpoint},
its successor-cycle analogue, and the matching root-distance pair.  So
downstream graph-level route consumers can now use the live exact shell
directly, without rebuilding those positive packages by hand.
\path{Theorem49PositiveGeometricSourceReplacementRoute.lean} now also gives
those four graph-level positive route packages their own
\path{boundaryRootSynthesis} and \path{exists_boundaryRootSynthesis} methods,
so once the exact shell has lowered into the named root-parent or
root-distance surfaces, no further manual descent through collar-layer
packages is needed to recover full Theorem~4.9 synthesis.
\path{Theorem49PositiveGeometricSourceReplacementJointObstruction.lean} now
adds a durable joint burden theorem on these same four named surfaces: each of
the closed-walk/successor-cycle root-parent and root-distance positive route
packages simultaneously forces a face boundary with two distinct interior
edges and excludes forcing interior-edge witnesses on the same embedding.  So
the constructive parent-oriented lane is no longer tracked only by endpoint
closure; it now carries an explicit combined local profile that any live
derivation into that lane must satisfy.  The same file now
closes one step lower still: it proves the direct exact-shell graph-level
endpoint and raw quotient/error wrappers
\path{exists_theorem49BoundaryRootNonemptyProjectedSynthesis_of_exists_closedWalkAnnulusBoundarySource_and_oneCollarAnnulusCurrentBoundaryData_and_v23ResidualBoundaryInitialState_and_planarBoundaryAnnulusRootParentPeelData_and_hasUnblockedInteriorEndpoint_direct},
its successor-cycle analogue, the matching generic adj-distance pair, and the
four corresponding \path{exists_theorem49BoundaryRawQuotientErrorPackage...}
theorems.  So once the exact shell already carries parent-peel or generic
adj-distance data together with \path{HasUnblockedInteriorEndpoint}, the route
no longer needs even an explicit graph-level positive-surface unpacking step
before reaching the current projected endpoint or raw quotient/error package.
The same file now also proves the direct graph-level full-synthesis closures
\path{exists_theorem49BoundaryRootSynthesis_of_exists_closedWalkAnnulusBoundarySource_and_oneCollarAnnulusCurrentBoundaryData_and_v23ResidualBoundaryInitialState_and_planarBoundaryAnnulusRootParentPeelData_direct},
its successor-cycle analogue, and the matching generic adj-distance pair.  So
once the live exact shell already carries same-embedding parent-peel data or
generic adj-distance data at all, the route now reaches full
\path{Theorem49BoundaryRootSynthesis} directly; even the projected-endpoint or
positive-surface packaging steps are no longer downstream burdens below that
shell.
The residual interface now closes one step lower too.
\path{Theorem49ResidualBoundaryPeeling.lean} proves
\path{theorem49BoundaryRootSynthesis_of_residualBoundarySelectorData},
\path{theorem49BoundaryRootSynthesis_of_residualBoundaryLayerFacePeelWitnessData},
and
\path{theorem49BoundaryRootSynthesis_of_residualBoundaryPositiveProjectedGeometryOn},
while \path{Theorem49ResidualBoundaryRoute.lean} proves
\path{theorem49BoundaryRootSynthesis_of_closedWalkAnnulusBoundarySource_and_residualBoundarySelectorData}
and
\path{theorem49BoundaryRootSynthesis_of_closedWalkAnnulusBoundarySource_and_residualBoundaryLayerFacePeelWitnessData}.
So once the route reaches residual/current-boundary selector data itself, the
full Theorem~4.9 synthesis endpoint is already immediate on that embedding;
the lower residual witness package and \path{HasUnblockedInteriorEndpoint} are
no longer downstream burdens beyond that surface.
\path{Theorem49ResidualBoundaryPeeling.lean} now also proves
\path{nonempty_residualBoundarySelectorData_of_closedWalkAnnulusBoundarySource_and_interiorDualBoundaryRootAdjDistancePeelData}
and
\path{nonempty_residualBoundarySelectorData_of_boundaryReachabilityData_and_dartSuccessorCycleEmbeddingData_and_selectedBoundaryArc_and_interiorDualBoundaryRootAdjDistancePeelData}.
So the residual selector surface is no longer the first unresolved layer above
the live parent-oriented shell: generic
\path{InteriorDualBoundaryRootAdjDistancePeelData} already lowers to that
surface on the same embedding.  The live constructive burden therefore
sharpens again to deriving that adjacency-distance package, or any equivalent
cycle-breaking parent witness strong enough to build residual/current-boundary
selector data.  \path{Theorem49ResidualBoundaryPeeling.lean} now also proves
\path{theorem49ResidualBoundaryPositiveProjectedGeometryOn_of_boundaryData_and_interiorDualBoundaryRootAdjDistancePeelData_and_hasUnblockedInteriorEndpoint}
and
\path{theorem49CurrentBoundaryRemainders_of_boundaryData_and_interiorDualBoundaryRootAdjDistancePeelData_and_hasUnblockedInteriorEndpoint},
while \path{Theorem49ResidualBoundaryRoute.lean} adds the exact one-collar
closed-walk and successor-cycle wrappers.  So once the same embedding carries
generic \path{InteriorDualBoundaryRootAdjDistancePeelData} together with
\path{HasUnblockedInteriorEndpoint}, the route already recovers both the
residual/current-boundary positive projected geometry surface and a concrete
layer-face current-boundary remainder witness there.  These lower residual
surfaces therefore are not any additional constructive burden below the
adjacency-distance package; they are already downstream consequences of it.
The same route file now also exposes those residual/current-boundary
consequences one layer higher on the exact same-split parent-peel shell
itself, together with matching graph-level existential closures.  So once the
live exact shell reaches \path{PlanarBoundaryAnnulusRootParentPeelData}
together with \path{HasUnblockedInteriorEndpoint}, the residual interface is
already present there as well; downstream users no longer need to unpack the
intermediate annulus-parent positive package or rebuild existential residual
witnesses by hand.  \path{Theorem49ResidualBoundaryObstruction.lean} now calibrates that same live
burden directly on the exact one-collar/v23 shared-interior-pair shell.  It
proves
\path{sharedInteriorPair_closedWalkSource_tait_hasUnblockedInteriorEndpoint_and_oneCollarAnnulusCurrentBoundaryData_and_v23ResidualBoundaryInitialState_without_interiorDualBoundaryRootAdjDistancePeelData}
and
\path{exists_embedding_boundaryReachabilityData_and_dartSuccessorCycleEmbeddingData_and_selectedBoundaryArc_and_oneCollarAnnulusCurrentBoundaryData_and_taitEdgeColoring_and_hasUnblockedInteriorEndpoint_and_v23ResidualBoundaryInitialState_without_interiorDualBoundaryRootAdjDistancePeelData_sharedInteriorPair},
together with the closed-walk and successor-cycle failed-universal forms.  So
the current live adjacency-distance burden is not merely unproved on that
benchmark: even exact one-collar current-boundary data, a real Tait coloring,
\path{HasUnblockedInteriorEndpoint}, and the exact v23 residual initial state
still do not force any same-embedding generic
\path{InteriorDualBoundaryRootAdjDistancePeelData} there.  More sharply, the
same exact shell also fails to force weak annulus-collar geometry itself: the
new one-collar theorems
\path{exists_embedding_closedWalkAnnulusBoundarySource_and_oneCollarAnnulusCurrentBoundaryData_and_taitEdgeColoring_and_hasUnblockedInteriorEndpoint_and_v23ResidualBoundaryInitialState_without_planarBoundaryAnnulusCollarGeometry_sharedInteriorPair},
\path{not_forall_nonempty_planarBoundaryAnnulusCollarGeometry_of_closedWalkAnnulusBoundarySource_and_oneCollarAnnulusCurrentBoundaryData_and_v23ResidualBoundaryInitialState_and_taitEdgeColoring_and_hasUnblockedInteriorEndpoint_sharedInteriorPair},
\path{exists_embedding_boundaryReachabilityData_and_dartSuccessorCycleEmbeddingData_and_selectedBoundaryArc_and_oneCollarAnnulusCurrentBoundaryData_and_taitEdgeColoring_and_hasUnblockedInteriorEndpoint_and_v23ResidualBoundaryInitialState_without_planarBoundaryAnnulusCollarGeometry_sharedInteriorPair},
and
\path{not_forall_nonempty_planarBoundaryAnnulusCollarGeometry_of_boundaryReachabilityData_and_dartSuccessorCycleEmbeddingData_and_selectedBoundaryArc_and_oneCollarAnnulusCurrentBoundaryData_and_v23ResidualBoundaryInitialState_and_taitEdgeColoring_and_hasUnblockedInteriorEndpoint_sharedInteriorPair}
show that even this single route-facing same-embedding geometry shell is not
recovered there.  The same exact shell also fails at the central acyclic
positive burden itself: the one-collar theorems
\path{exists_embedding_closedWalkAnnulusBoundarySource_and_oneCollarAnnulusCurrentBoundaryData_and_taitEdgeColoring_and_hasUnblockedInteriorEndpoint_and_v23ResidualBoundaryInitialState_without_planarBoundaryWellFoundedFacePeelWitnessData_sharedInteriorPair},
\path{not_forall_nonempty_planarBoundaryWellFoundedFacePeelWitnessData_of_closedWalkAnnulusBoundarySource_and_oneCollarAnnulusCurrentBoundaryData_and_v23ResidualBoundaryInitialState_and_taitEdgeColoring_and_hasUnblockedInteriorEndpoint_sharedInteriorPair},
\path{exists_embedding_boundaryReachabilityData_and_dartSuccessorCycleEmbeddingData_and_selectedBoundaryArc_and_oneCollarAnnulusCurrentBoundaryData_and_taitEdgeColoring_and_hasUnblockedInteriorEndpoint_and_v23ResidualBoundaryInitialState_without_planarBoundaryWellFoundedFacePeelWitnessData_sharedInteriorPair},
and
\path{not_forall_nonempty_planarBoundaryWellFoundedFacePeelWitnessData_of_boundaryReachabilityData_and_dartSuccessorCycleEmbeddingData_and_selectedBoundaryArc_and_oneCollarAnnulusCurrentBoundaryData_and_v23ResidualBoundaryInitialState_and_taitEdgeColoring_and_hasUnblockedInteriorEndpoint_sharedInteriorPair}
show that even same-embedding well-founded parent-peeling is not recovered
there.  More primitively, the
same exact shell also fails at the earlier child-membership face-membership
shell itself: the new one-collar theorems
\path{exists_embedding_closedWalkAnnulusBoundarySource_and_oneCollarAnnulusCurrentBoundaryData_and_taitEdgeColoring_and_hasUnblockedInteriorEndpoint_and_v23ResidualBoundaryInitialState_without_boundaryFaceRootsCanonicalParentSharedEdgeFaceMembershipCover_sharedInteriorPair},
\path{not_forall_some_closedWalkAnnulusBoundarySourceBoundaryFaceRootsCanonicalParentSharedEdgeFaceMembershipCover_of_closedWalkAnnulusBoundarySource_and_oneCollarAnnulusCurrentBoundaryData_and_v23ResidualBoundaryInitialState_and_taitEdgeColoring_and_hasUnblockedInteriorEndpoint_sharedInteriorPair},
\path{exists_embedding_boundaryReachabilityData_and_dartSuccessorCycleEmbeddingData_and_selectedBoundaryArc_and_oneCollarAnnulusCurrentBoundaryData_and_taitEdgeColoring_and_hasUnblockedInteriorEndpoint_and_v23ResidualBoundaryInitialState_without_boundaryFaceRootsCanonicalParentSharedEdgeFaceMembershipCover_sharedInteriorPair},
and
\path{not_forall_some_successorCycleBoundaryFaceRootsCanonicalParentSharedEdgeFaceMembershipCover_of_boundaryReachabilityData_and_dartSuccessorCycleEmbeddingData_and_selectedBoundaryArc_and_oneCollarAnnulusCurrentBoundaryData_and_v23ResidualBoundaryInitialState_and_taitEdgeColoring_and_hasUnblockedInteriorEndpoint_sharedInteriorPair}
show that even exact one-collar current-boundary data preserving the extracted
boundary split still does not recover that stronger source-preserving shell on
the same embedding.  The same file now also lifts this face-membership
obstruction to generic exact-shell converse failures:
\path{not_forall_some_closedWalkAnnulusBoundarySourceBoundaryFaceRootsCanonicalParentSharedEdgeFaceMembershipCover_of_exists_embedding_closedWalkAnnulusBoundarySource_and_taitEdgeColoring_and_hasUnblockedInteriorEndpoint_and_v23ResidualBoundaryInitialState},
\path{not_forall_some_closedWalkAnnulusBoundarySourceBoundaryFaceRootsCanonicalParentSharedEdgeFaceMembershipCover_of_exists_embedding_closedWalkAnnulusBoundarySource_and_oneCollarAnnulusCurrentBoundaryData_and_taitEdgeColoring_and_hasUnblockedInteriorEndpoint_and_v23ResidualBoundaryInitialState},
\path{not_forall_some_successorCycleBoundaryFaceRootsCanonicalParentSharedEdgeFaceMembershipCover_of_exists_embedding_boundaryReachabilityData_and_dartSuccessorCycleEmbeddingData_and_selectedBoundaryArc_and_taitEdgeColoring_and_hasUnblockedInteriorEndpoint_and_v23ResidualBoundaryInitialState},
and
\path{not_forall_some_successorCycleBoundaryFaceRootsCanonicalParentSharedEdgeFaceMembershipCover_of_exists_embedding_boundaryReachabilityData_and_dartSuccessorCycleEmbeddingData_and_selectedBoundaryArc_and_oneCollarAnnulusCurrentBoundaryData_and_taitEdgeColoring_and_hasUnblockedInteriorEndpoint_and_v23ResidualBoundaryInitialState}.
So any same-embedding exact-shell witness with
\path{HasUnblockedInteriorEndpoint} already refutes a universal derivation of
that stricter face-membership shell; the shared-interior-pair examples are just
concrete instantiations.  More primitively, the same exact shell also fails at the
earlier raw-cover and degenerate boundary-face-root shells themselves: the new one-collar theorems
\path{exists_embedding_closedWalkAnnulusBoundarySource_and_oneCollarAnnulusCurrentBoundaryData_and_taitEdgeColoring_and_hasUnblockedInteriorEndpoint_and_v23ResidualBoundaryInitialState_without_boundaryFaceRootsCanonicalParentSharedEdgeCover_sharedInteriorPair},
\path{not_forall_some_closedWalkAnnulusBoundarySourceBoundaryFaceRootsCanonicalParentSharedEdgeCover_of_closedWalkAnnulusBoundarySource_and_oneCollarAnnulusCurrentBoundaryData_and_v23ResidualBoundaryInitialState_and_taitEdgeColoring_and_hasUnblockedInteriorEndpoint_sharedInteriorPair},
\path{exists_embedding_boundaryReachabilityData_and_dartSuccessorCycleEmbeddingData_and_selectedBoundaryArc_and_oneCollarAnnulusCurrentBoundaryData_and_taitEdgeColoring_and_hasUnblockedInteriorEndpoint_and_v23ResidualBoundaryInitialState_without_boundaryFaceRootsCanonicalParentSharedEdgeCover_sharedInteriorPair},
\path{not_forall_some_successorCycleBoundaryFaceRootsCanonicalParentSharedEdgeCover_of_boundaryReachabilityData_and_dartSuccessorCycleEmbeddingData_and_selectedBoundaryArc_and_oneCollarAnnulusCurrentBoundaryData_and_v23ResidualBoundaryInitialState_and_taitEdgeColoring_and_hasUnblockedInteriorEndpoint_sharedInteriorPair},
and the matching \path{NoInteriorEdges} quartet show that even exact
one-collar current-boundary data preserving the extracted boundary split still
does not recover those earlier route shells on the same embedding.  The same
file now also lifts the raw-cover half of this diagnosis to generic exact-shell
converse failures:
\path{not_forall_some_closedWalkAnnulusBoundarySourceBoundaryFaceRootsCanonicalParentSharedEdgeCover_of_exists_embedding_closedWalkAnnulusBoundarySource_and_taitEdgeColoring_and_hasUnblockedInteriorEndpoint_and_v23ResidualBoundaryInitialState},
\path{not_forall_some_closedWalkAnnulusBoundarySourceBoundaryFaceRootsCanonicalParentSharedEdgeCover_of_exists_embedding_closedWalkAnnulusBoundarySource_and_oneCollarAnnulusCurrentBoundaryData_and_taitEdgeColoring_and_hasUnblockedInteriorEndpoint_and_v23ResidualBoundaryInitialState},
\path{not_forall_some_successorCycleBoundaryFaceRootsCanonicalParentSharedEdgeCover_of_exists_embedding_boundaryReachabilityData_and_dartSuccessorCycleEmbeddingData_and_selectedBoundaryArc_and_taitEdgeColoring_and_hasUnblockedInteriorEndpoint_and_v23ResidualBoundaryInitialState},
and
\path{not_forall_some_successorCycleBoundaryFaceRootsCanonicalParentSharedEdgeCover_of_exists_embedding_boundaryReachabilityData_and_dartSuccessorCycleEmbeddingData_and_selectedBoundaryArc_and_oneCollarAnnulusCurrentBoundaryData_and_taitEdgeColoring_and_hasUnblockedInteriorEndpoint_and_v23ResidualBoundaryInitialState}.
So any same-embedding exact-shell witness with
\path{HasUnblockedInteriorEndpoint} already refutes a universal derivation of
that raw canonical-parent shell; the shared-interior-pair examples are just
concrete instantiations.  The same exact shell also fails to repair the entire
currently known same-embedding geometry burden in generic form:
\path{not_forall_some_knownSameEmbeddingGeometry_of_exists_embedding_closedWalkAnnulusBoundarySource_and_oneCollarAnnulusCurrentBoundaryData_and_taitEdgeColoring_and_hasUnblockedInteriorEndpoint_and_v23ResidualBoundaryInitialState}
and
\path{not_forall_some_knownSameEmbeddingGeometry_of_exists_embedding_boundaryReachabilityData_and_dartSuccessorCycleEmbeddingData_and_selectedBoundaryArc_and_oneCollarAnnulusCurrentBoundaryData_and_taitEdgeColoring_and_hasUnblockedInteriorEndpoint_and_v23ResidualBoundaryInitialState}
show that any explicit exact-shell witness carrying simultaneous failure of the
currently known same-embedding geometry surfaces already refutes a universal
derivation of that whole disjunctive burden.  The same
file now
also proves
\path{not_nonempty_planarBoundaryAnnulusRootAdjDistancePeelData_sharedInteriorPair},
\path{sharedInteriorPair_closedWalkSource_tait_hasUnblockedInteriorEndpoint_and_oneCollarAnnulusCurrentBoundaryData_and_v23ResidualBoundaryInitialState_without_planarBoundaryAnnulusRootAdjDistancePeelData},
\path{exists_embedding_boundaryReachabilityData_and_dartSuccessorCycleEmbeddingData_and_selectedBoundaryArc_and_oneCollarAnnulusCurrentBoundaryData_and_taitEdgeColoring_and_hasUnblockedInteriorEndpoint_and_v23ResidualBoundaryInitialState_without_planarBoundaryAnnulusRootAdjDistancePeelData_sharedInteriorPair},
and the matching closed-walk / successor-cycle failed-universal theorems.  The
same obstruction is now lifted to reusable forcing-edge converse failures too:
\path{not_forall_nonempty_planarBoundaryAnnulusRootAdjDistancePeelData_of_exists_embedding_closedWalkAnnulusBoundarySource_and_oneCollarAnnulusCurrentBoundaryData_and_taitEdgeColoring_and_hasUnblockedInteriorEndpoint_and_v23ResidualBoundaryInitialState_and_forcingInteriorEdgeWitness}
and
\path{not_forall_nonempty_planarBoundaryAnnulusRootAdjDistancePeelData_of_exists_embedding_boundaryReachabilityData_and_dartSuccessorCycleEmbeddingData_and_selectedBoundaryArc_and_oneCollarAnnulusCurrentBoundaryData_and_taitEdgeColoring_and_hasUnblockedInteriorEndpoint_and_v23ResidualBoundaryInitialState_and_forcingInteriorEdgeWitness}.
So the same repaired exact one-collar/v23 shell does not merely fail to force
generic interior-dual peel data: it already fails to force the source-preserving
two-sided annulus-root adjacency-distance route package on the same embedding.
The same file now
also proves
\path{not_nonempty_planarBoundaryAnnulusRootParentPeelData_sharedInteriorPair},
\path{sharedInteriorPair_closedWalkSource_tait_hasUnblockedInteriorEndpoint_and_oneCollarAnnulusCurrentBoundaryData_and_v23ResidualBoundaryInitialState_without_planarBoundaryAnnulusRootParentPeelData},
\path{exists_embedding_boundaryReachabilityData_and_dartSuccessorCycleEmbeddingData_and_selectedBoundaryArc_and_oneCollarAnnulusCurrentBoundaryData_and_taitEdgeColoring_and_hasUnblockedInteriorEndpoint_and_v23ResidualBoundaryInitialState_without_planarBoundaryAnnulusRootParentPeelData_sharedInteriorPair},
and the matching closed-walk / successor-cycle failed-universal theorems.  So
the same repaired exact one-collar/v23 shell does not merely fail at the
stronger generic adjacency-distance package: it already fails to force the
actual same-embedding \path{PlanarBoundaryAnnulusRootParentPeelData} route
target itself.  The live diagnosis is therefore sharper than ``derive
parent-peel from the residual shell'': that implication is false on the
current shell, so the route must either strengthen the upstream same-embedding
hypotheses or turn the residual shell into an impossibility theorem instead of
continuing to lower downstream.  The same obstruction is now lifted from the
benchmark instance to reusable route-level failed-universal theorems too:
\path{not_forall_nonempty_planarBoundaryAnnulusRootParentPeelData_of_exists_embedding_closedWalkAnnulusBoundarySource_and_oneCollarAnnulusCurrentBoundaryData_and_taitEdgeColoring_and_hasUnblockedInteriorEndpoint_and_v23ResidualBoundaryInitialState_and_forcingInteriorEdgeWitness}
and
\path{not_forall_nonempty_planarBoundaryAnnulusRootParentPeelData_of_exists_embedding_boundaryReachabilityData_and_dartSuccessorCycleEmbeddingData_and_selectedBoundaryArc_and_oneCollarAnnulusCurrentBoundaryData_and_taitEdgeColoring_and_hasUnblockedInteriorEndpoint_and_v23ResidualBoundaryInitialState_and_forcingInteriorEdgeWitness}.
So any same-embedding exact-shell witness carrying a forcing interior edge now
formally refutes the parent-peel converse; the obstruction is no longer stated
only as a single shared-interior-pair counterexample.
\path{Theorem49InteriorVerticesEndpointObstruction.lean} now exposes the
opposite concrete calibration:
\path{nonempty_interiorDualBoundaryRootAdjDistancePeelData_peeledEndpointTouch}
and
\path{admitsPlanarBoundaryInteriorDualBoundaryRootAdjDistancePeelData_withoutClosedWalkAnnulusBoundarySource_peeledEndpointTouchGraph}
show that generic boundary-root adjacency-distance data is itself realizable
on a concrete benchmark, but only off the live source-preserving route.  The
endpoint-touch graph is acyclic, so it cannot satisfy the honest closed-walk
annulus source at all.  Thus the shared-interior-pair exact-shell failure is a
source-compatible obstruction on the live route, not an absolute impossibility
of \path{InteriorDualBoundaryRootAdjDistancePeelData} itself.
\path{Theorem49ForcingInteriorEdgeObstructionRegression.lean} now strengthens
that diagnosis on an honest source-compatible benchmark too.  The chained-
diamonds source already inhabits the stronger local package
\path{exists_closedWalkSource_with_atMostOneInteriorEdgePerFace_chainedDiamondsWithTriangleGraph},
which is enough elsewhere in the development to derive canonical witness choice
and weak collar geometry.  But the new theorems
\path{not_nonempty_interiorDualBoundaryRootAdjDistancePeelData_chainedDiamondsWithTriangle_via_forcingInteriorEdgeWitness}
and
\path{not_forall_nonempty_interiorDualBoundaryRootAdjDistancePeelData_of_closedWalkSource_with_atMostOneInteriorEdgePerFace_chainedDiamondsWithTriangleGraph}
show that even this stronger honest closed-walk source package still does not
force generic boundary-root adjacency-distance data on the same embedding.  So
the adj-distance obstruction is not peculiar to the repaired exact
one-collar/v23 shell; it already appears higher up on a source-compatible local
source surface.  The same file now also proves
\path{not_nonempty_planarBoundaryAnnulusRootParentPeelData_chainedDiamondsWithTriangle_via_forcingInteriorEdgeWitness}
and
\path{not_forall_nonempty_planarBoundaryAnnulusRootParentPeelData_of_closedWalkSource_with_atMostOneInteriorEdgePerFace_chainedDiamondsWithTriangleGraph},
so even that stronger honest-source package still does not force the current
parent-oriented annulus-root target on the same embedding.  This sharpens the
live burden again: the route now needs genuinely stronger source-side structure
than the chained-diamonds at-most-one package, not just more downstream
lowering from it.
That same repaired exact-shell package now also has an honest full
Theorem~4.9 synthesis endpoint once the exact construction induced by that
package carries witness ownership.  \path{Theorem49ResidualBoundaryRoute.lean}
proves
\path{theorem49BoundaryRootSynthesis_of_closedWalkAnnulusBoundarySource_and_oneCollarAnnulusCurrentBoundaryData_and_v23ResidualBoundaryInitialState_and_planarBoundaryAnnulusRootParentPeelData_and_selectorDeficitSelectedBoundary_and_inducedDecompositionWitnessAssignment},
\path{theorem49BoundaryRootSynthesis_of_boundaryReachabilityData_and_dartSuccessorCycleEmbeddingData_and_selectedBoundaryArc_and_oneCollarAnnulusCurrentBoundaryData_and_v23ResidualBoundaryInitialState_and_planarBoundaryAnnulusRootParentPeelData_and_selectorDeficitSelectedBoundary_and_inducedDecompositionWitnessAssignment},
and also the older chosen-decomposition forms
\path{theorem49BoundaryRootSynthesis_of_closedWalkAnnulusBoundarySource_and_oneCollarAnnulusCurrentBoundaryData_and_v23ResidualBoundaryInitialState_and_planarBoundaryAnnulusRootParentPeelData_and_selectorDeficitSelectedBoundary_and_resultingDecompositionWitnessAssignment}
and
\path{theorem49BoundaryRootSynthesis_of_boundaryReachabilityData_and_dartSuccessorCycleEmbeddingData_and_selectedBoundaryArc_and_oneCollarAnnulusCurrentBoundaryData_and_v23ResidualBoundaryInitialState_and_planarBoundaryAnnulusRootParentPeelData_and_selectorDeficitSelectedBoundary_and_resultingDecompositionWitnessAssignment}.
So once the exact shell carries the same-split parent-peel and
selector-deficit package, no further Theorem~4.9 algebra remains after the
derived decomposition acquires witness ownership.  More sharply, the live
downstream burden is now to build \path{PlanarBoundaryAnnulusWitnessAssignment}
on the exact same-embedding annulus decomposition canonically induced by that
repaired parent-peel construction, rather than on an arbitrary chosen
decomposition or through another selector wrapper.  The new theorem
\path{theorem49BoundaryRootSynthesis_of_boundaryReachabilityData_and_dartSuccessorCycleEmbeddingData_and_selectedBoundaryArc_and_oneCollarAnnulusCurrentBoundaryData_and_v23ResidualBoundaryInitialState_and_planarBoundaryAnnulusRootParentPeelData_and_nonPeelSelectedBoundary_and_resultingDecompositionWitnessAssignment}
packages this downward derivation on the live successor-cycle shell itself:
under the older stronger non-peeled-face contact hypothesis, the
selector-deficit burden is discharged internally and the remaining exact-shell
work is to reach that boundary-support shell and then supply witness
ownership.
The live wheel benchmark now shows that this burden is genuinely new.
\path{Theorem49ForcingInteriorEdgeObstructionRegression.lean} proves
\path{not_exists_planarBoundaryAnnulusDecomposition_and_witnessAssignment_wheelWithInnerTriangle},
\path{exists_embedding_boundaryReachabilityData_and_dartSuccessorCycleEmbeddingData_and_selectedBoundaryArc_and_oneCollarAnnulusCurrentBoundaryData_and_taitEdgeColoring_and_nonempty_selectedBoundaryInteriorCarrier_and_v23ResidualBoundaryInitialState_without_any_planarBoundaryAnnulusWitnessAssignment_wheelWithInnerTriangle},
and
\path{not_forall_exists_planarBoundaryAnnulusDecomposition_and_witnessAssignment_of_boundaryReachabilityData_and_dartSuccessorCycleEmbeddingData_and_selectedBoundaryArc_and_oneCollarAnnulusCurrentBoundaryData_and_v23ResidualBoundaryInitialState_and_taitEdgeColoring_and_nonempty_selectedBoundaryInteriorCarrier_wheelWithInnerTriangle}.
So even the actual successor-cycle one-collar exact-v23 shell on the live
wheel carrier still does not force any same-embedding annulus decomposition
carrying \path{PlanarBoundaryAnnulusWitnessAssignment}.  The missing repaired
burden is therefore honestly witness ownership itself, not any further
decomposition-side wrapper.
That same repaired local package now already discharges the forcing-edge
obstruction on the exact one-collar v23/v23.5 shell itself.
\path{Theorem49ResidualBoundaryRoute.lean} proves
\path{not_nonempty_forcingInteriorEdgeWitness_of_closedWalkAnnulusBoundarySource_and_oneCollarAnnulusCurrentBoundaryData_and_v23ResidualBoundaryInitialState_and_planarBoundaryAnnulusRootParentPeelData_and_selectorDeficitSelectedBoundary}
and
\path{not_nonempty_forcingInteriorEdgeWitness_of_boundaryReachabilityData_and_dartSuccessorCycleEmbeddingData_and_selectedBoundaryArc_and_oneCollarAnnulusCurrentBoundaryData_and_v23ResidualBoundaryInitialState_and_planarBoundaryAnnulusRootParentPeelData_and_selectorDeficitSelectedBoundary}.
So once the exact shell carries same-split parent peel data together with the
selector-deficit selected-boundary hypothesis and the positive
\path{currentOuterBoundary} side conditions, forcing interior edges are
excluded on that same embedding.  The remaining exact-shell burden is
therefore no longer an unspecified no-forcing theorem, but precisely to derive
this explicit selector-deficit construction package from the actual residual
shell.  Lean now also records that this burden can be stated one layer more
directly on the actual boundary-order selector package itself:
\path{not_nonempty_forcingInteriorEdgeWitness_of_boundaryReachabilityData_and_closedWalkEmbeddingData_and_selectedBoundaryArcOnFace_and_boundaryFreeIncidentFaceSelector_and_boundarySupportFaceData_conditions}
and
\path{theorem49BoundaryRootSynthesis_of_boundaryReachabilityData_and_closedWalkEmbeddingData_and_selectedBoundaryArcOnFace_and_boundaryFreeIncidentFaceSelector_and_boundarySupportFaceData_conditions_and_inducedDecompositionWitnessAssignment}.
The same route is now also named directly on the actual successor-cycle shell
by
\path{not_nonempty_forcingInteriorEdgeWitness_of_boundaryReachabilityData_and_dartSuccessorCycleEmbeddingData_and_selectedBoundaryArc_and_boundaryFreeIncidentFaceSelector_and_boundarySupportFaceData_conditions},
\path{theorem49BoundaryRootSynthesis_of_boundaryReachabilityData_and_dartSuccessorCycleEmbeddingData_and_selectedBoundaryArc_and_boundaryFreeIncidentFaceSelector_and_boundarySupportFaceData_conditions_and_inducedDecompositionWitnessAssignment},
and the graph-level closure
\path{exists_theorem49BoundaryRootSynthesis_of_exists_boundaryReachabilityData_and_dartSuccessorCycleEmbeddingData_and_selectedBoundaryArc_and_boundaryFreeIncidentFaceSelector_and_boundarySupportFaceData_conditions_and_inducedDecompositionWitnessAssignment_direct}.
So once boundary reachability together with facewise selected-boundary arcs
already supports a boundary-free selector whose nonselected faces meet
selected-boundary support and whose induced positive
\path{currentOuterBoundary} layers are nonempty and pairwise disjoint, both the
forcing-edge obstruction and the remaining Theorem~4.9 algebra disappear on
that same embedding.  More sharply, the remaining burden is witness ownership
on the exact selector-induced annulus decomposition, not merely on an
unspecified chosen decomposition.  Thus the live exact-shell burden is better
stated directly at this selector-side local package, and then at that induced
decomposition, rather than at another intermediate construction shell.  Once
that exact successor-cycle selector package is available at graph level, no
further route-level repackaging step remains before full
\path{Theorem49BoundaryRootSynthesis}.
This forcing-edge route now also blocks the current two-sided annulus-root lane
itself.  \path{Theorem49ForcingInteriorEdgeObstruction.lean} now proves
\path{closedWalkAnnulusBoundarySource_does_not_imply_planarBoundaryAnnulusRootAdjDistancePeelData}
and the same-embedding existential form
\path{not_exists_embedding_closedWalkAnnulusBoundarySource_and_forcingInteriorEdgeWitness_and_planarBoundaryAnnulusRootAdjDistancePeelData},
with the diamond-with-triangle regression instantiated by
\path{closedWalkAnnulusBoundarySource_does_not_imply_planarBoundaryAnnulusRootAdjDistancePeelData_diamondWithTriangle_via_forcingInteriorEdgeWitness}.
So a forcing interior edge does not merely obstruct later annulus construction
or synthesis shells; it already prevents any two-sided annulus-root
repackaging on that same honest-source embedding.
This has now been pushed one step earlier in the boundary-order pipeline as
well.  The same file proves
\path{annulusBoundaryReachability_and_cyclicOrderedFaceArcEmbeddingData_does_not_imply_planarBoundaryAnnulusRootAdjDistancePeelData}
and
\path{annulusBoundaryReachability_and_selectedBoundaryArcGeometry_does_not_imply_planarBoundaryAnnulusRootAdjDistancePeelData},
plus the corresponding same-embedding existential obstructions.  The regression
file instantiates both on the diamond-with-triangle witness.  So even the
stronger cyclic boundary-order shells still do not rule out the forcing-edge
pathology needed to reach two-sided annulus-root data.
This forcing-edge calibration is now also packaged at the right theorem
surface.  \path{Theorem49ForcingInteriorEdgeObstruction.lean} proves the
generic failed-converse theorem
\path{not_forall_nonempty_planarBoundaryAnnulusRootAdjDistancePeelData_of_exists_embedding_with_property_and_forcingInteriorEdgeWitness}:
any explicit embedding carrying a property $P$ together with a forcing
interior-edge witness refutes a universal same-embedding theorem from $P$ to
\path{PlanarBoundaryAnnulusRootAdjDistancePeelData}.  The same file records the
route-facing specializations where $P$ is the honest closed-walk annulus
source shell, the reachability-plus-cyclic shell, and the
reachability-plus-selected-arc shell; the regression file instantiates all
three on \path{diamondWithTriangleGraph}.  So the forcing-edge blocker is no
longer only a list of local witness failures.  It is now a reusable
same-embedding no-universal-converse theorem surface for the current
annulus-root route.
The missing burden is now stated even more directly.  The same file proves the
generic theorem
\path{not_forall_not_nonempty_forcingInteriorEdgeWitness_of_exists_embedding_with_property_and_forcingInteriorEdgeWitness}:
any explicit embedding carrying a property $P$ together with a forcing
interior-edge witness refutes a universal same-embedding theorem saying that
$P$ excludes forcing interior edges.  The route-facing specializations cover
the honest closed-walk annulus source shell, the reachability-plus-cyclic
shell, and the reachability-plus-selected-arc shell, with the regression file
again instantiating all three on \path{diamondWithTriangleGraph}.  So the
remaining geometric burden is explicit in Lean: current source-side boundary
order data does not yet rule out the forcing-edge pathology that blocks the
surviving annulus-root lane.
The attempted face-local strengthening is now also refuted.  The obstruction
file proves
\path{not_forall_not_nonempty_forcingInteriorEdgeWitness_of_exists_embedding_reachability_and_closedWalkEmbedding_and_atMostOneInteriorEdge_and_forcingInteriorEdgeWitness},
and the regression file instantiates it through
\path{exists_embedding_reachability_and_closedWalkEmbedding_and_atMostOneInteriorEdge_and_forcingInteriorEdgeWitness_diamondWithTriangleGraph}
and
\path{not_forall_not_nonempty_forcingInteriorEdgeWitness_of_reachability_and_closedWalkEmbedding_and_atMostOneInteriorEdge_diamondWithTriangle}.
Thus, on \path{diamondWithTriangleEmbedding}, annulus boundary reachability,
honest facial closed walks, and the condition that every face has at most one
interior boundary edge coexist with the forcing interior edge \path{dt12}.  The
positive exclusion burden therefore cannot be met by this at-most-one
condition alone; it needs a genuinely stronger collar/peel ownership
principle or a different geometric exclusion hypothesis.
That calibration is now sharpened on the honest source shell itself.
\path{PlanarBoundaryClosedWalkSource.lean} now proves that on a fixed embedding
\path{Nonempty (PlanarBoundaryClosedWalkAnnulusBoundarySource emb)} already
lowers to both \path{Nonempty (PlanarBoundaryCyclicOrderedFaceArcEmbeddingData emb)}
and \path{Nonempty (PlanarBoundarySelectedBoundaryArcGeometry emb)}.
\path{Theorem49ForcingInteriorEdgeObstruction.lean} then adds
\path{closedWalkAnnulusBoundarySource_and_cyclicOrderedFaceArcEmbeddingData_does_not_imply_planarBoundaryAnnulusRootAdjDistancePeelData}
and its selected-arc analogue, while the regression file realizes both on
\path{diamondWithTriangleEmbedding}.  So adjoining the stronger cyclic or
bundled selected-boundary-arc shells on top of the honest closed-walk source
still does not rescue the annulus-root route from the forcing-edge blocker.
The same strengthened honest-source surfaces are now blocked at graph level
too.  The obstruction file adds
\path{not_exists_embedding_closedWalkAnnulusBoundarySource_and_cyclicOrderedFaceArcEmbeddingData_and_forcingInteriorEdgeWitness_and_planarBoundaryAnnulusRootAdjDistancePeelData}
and the selected-arc analogue, with the regression file instantiating both on
\path{diamondWithTriangleGraph}.  So even the conjunction ``honest source +
honest cyclic boundary-order shell'' is now formally separated from the
surviving annulus-root lane by the forcing-edge witness.
So the honest closed-walk annulus-boundary source still does not supply the
peeled-interior geometric burden needed to start the current annulus
construction route.
A sharper regression now reaches the peeled edge-disjoint data itself.  The
minimal model \path{peeledEndpointTouchEmbedding} admits
\path{peeledEndpointTouchOuterBoundaryRootAdjDistancePeelData}: its peeled face
contains only the shared interior edge, so the current
\path{hpeelInterior} edge-disjoint field is satisfied.  But that interior edge
shares an endpoint with a selected-boundary edge on the root face, and
\path{outerBoundaryRootAdjDistancePeelData_does_not_imply_endpointSupport_disjoint_peeledEndpointTouch}
therefore refutes endpoint-support disjointness even under the current
outer-boundary-rooted adjacency-distance peel package.  The extra peeled-face
endpoint no-touch hypothesis is not cosmetic; it is strictly stronger than the
current edge-disjoint annulus data.  The direct theorem
\path{peelFaces_endpoint_disjoint_selectedBoundarySupport_fails_peeledEndpointTouch}
checks the exact construction-level obligation: the peeled singleton face and
the selected boundary support share a primal endpoint, so the no-touch bridge
cannot be recovered from the present peeled package by unpacking definitions.
The theorem
\path{planarBoundaryAnnulusConstruction_edge_disjoint_peelFaces_does_not_imply_endpoint_noTouch_peeledEndpointTouch}
packages the same separation at the bare annulus-construction surface: its
peeled faces are edge-disjoint from the selected boundary support, but the
endpoint no-touch condition still fails.  The negated universal theorem
\path{not_forall_planarBoundaryAnnulusConstruction_edge_disjoint_implies_endpoint_noTouch_peeledEndpointTouch}
rules out deriving the endpoint no-touch condition from annulus-construction
data plus peeled-face edge disjointness, even on this fixed embedding.
The no-touch bridge is now also factored through the aggregate finite carrier
\path{peelFaceBoundaryEndpointSupport}.  The lemma
\path{interiorEdgeEndpointSupport_subset_peelFaceBoundaryEndpointSupport}
says that every raw interior-edge endpoint lies on this peeled-face boundary
endpoint carrier, and
\path{peelFacesEndpointDisjointSelectedBoundarySupport_iff} identifies
aggregate carrier disjointness with the previous per-peeled-face no-touch
schema.  The endpoint-touch model refutes this aggregate repair as well:
\path{planarBoundaryAnnulusConstruction_edge_disjoint_peelFaces_does_not_imply_aggregate_endpoint_noTouch_peeledEndpointTouch}
and
\path{not_forall_planarBoundaryAnnulusConstruction_edge_disjoint_implies_aggregate_endpoint_noTouch_peeledEndpointTouch}
show that peeled-face edge disjointness does not imply disjointness of the
peeled-face endpoint carrier from selected-boundary endpoints.
At the outer-boundary-rooted route surface, the file also proves
\path{outerBoundaryRootAdjDistancePeelData_does_not_imply_withEndpointSeparation_peeledEndpointTouch}:
the same embedding admits the current peeled route package, but admits no
\path{PlanarBoundaryOuterBoundaryRootAdjDistancePeelDataWithEndpointSeparation}
package because raw endpoint support disjointness is false.
The companion theorem
\path{selectedBoundaryInteriorEdgeEndpointVertices_eq_empty_peeledEndpointTouch}
records the positive alternative on the same model: the selected-boundary
purified endpoint carrier is exactly empty.  The raw interior edge contributes
vertices \(1\) and \(2\), but \path{pet01} removes vertex \(1\) and
\path{pet23} removes vertex \(2\).
The summary theorem
\path{peeledEndpointTouch_rawCarrierRepairObstructionSummary} packages the
strongest current conclusion in one Lean statement.  This should be read as a
high-confidence crux-block for one precise repair strategy: deriving the raw
\path{interiorEdgeEndpointSupport} carrier, or the peeled-face endpoint
no-touch premise needed for it, from the current outer-boundary-rooted peeled
annulus data and peeled-face edge disjointness.  It should not be overstated
as a high-confidence block of every possible Theorem~4.9 continuation.  Two
honest continuations remain open in the formalization: add and prove a
genuinely stronger geometric endpoint no-touch hypothesis, or keep the
selected-boundary purification explicit, as in
\path{Theorem49OuterBoundaryRootSynthesis}.
The theorem \path{peeledEndpointTouch_emptyCarrier_synthesis} makes the fork
even more explicit on the same fixed model: for any Tait coloring, the current
synthesis package is available from
\path{peeledEndpointTouchOuterBoundaryRootAdjDistancePeelData}, but its
selected-boundary endpoint carrier is exactly \(\emptyset\).  Thus the
positive synthesis theorem and the raw-carrier obstruction are not
contradictory; they prove that this data package supports the purified
statement while carrying no non-vacuous raw endpoint information in the
endpoint-touch model.
The Lean surface now separates this explicitly:
\path{Theorem49OuterBoundaryRootNonemptySynthesis} is the same synthesis
package plus the additional proof that the purified carrier is nonempty.  The
endpoint-touch model has nonempty raw carrier, certified by
\path{interiorEdgeEndpointSupport_nonempty_peeledEndpointTouch}, but
\path{not_theorem49OuterBoundaryRootNonemptySynthesis_peeledEndpointTouch}
shows that it cannot satisfy this non-vacuous selected-boundary synthesis
surface.  This blocks the tempting interpretation that the purified endpoint
theorem has preserved a genuine annular interior vertex in this model.
The same issue now survives on the newer two-sided annulus-root synthesis lane
too.  The endpoint-touch model carries
\path{peeledEndpointTouchAnnulusRootAdjDistancePeelData}, and
\path{peeledEndpointTouch_emptyCarrier_annulusSynthesis} shows that the
ordinary \path{Theorem49AnnulusRootSynthesis} package is still available there
for any Tait coloring, but only over the empty purified carrier.  Likewise
\path{not_theorem49AnnulusRootNonemptySynthesis_peeledEndpointTouch} and
\path{peeledEndpointTouch_annulusNonvacuousPurifiedSynthesisCruxObstructionSummary}
show that abandoning the all-outer-root condition does not by itself repair the
non-vacuity problem: the selected-boundary purification can still collapse the
endpoint carrier exactly on the two-sided annulus-root surface.
This collapse already survives one further step upstream on the direct
boundary height-ordered witness-cover lane.  The endpoint-touch model now
carries
\path{peeledEndpointTouchPlanarBoundaryHeightOrderedFacePeelWitnessData}, and
\path{planarBoundaryHeightOrderedFacePeelWitnessData_does_not_imply_nonempty_selectedBoundaryInteriorEdgeEndpointVertices_peeledEndpointTouch}
shows that even this earlier planarity-side witness endpoint can coexist with a
nonempty raw interior-edge endpoint carrier and an empty purified
selected-boundary carrier.  The packaged theorem
\path{peeledEndpointTouch_heightOrderedWitnessCarrierCruxObstructionSummary}
records that the graph already admits
\path{AdmitsPlanarBoundaryHeightOrderedFacePeelWitnessData} while the purified
carrier is still exactly \(\emptyset\), and
\path{admitsPlanarBoundaryHeightOrderedFacePeelWitnessData_withoutClosedWalkAnnulusBoundarySource_peeledEndpointTouchGraph}
shows that this earlier witness-cover endpoint still does not recover the
honest closed-walk annulus source on the same acyclic graph.
This collapse already survives one further step upstream on the direct
well-founded witness-cover lane.  The endpoint-touch model now carries
\path{peeledEndpointTouchPlanarBoundaryWellFoundedFacePeelWitnessData}, and
\path{planarBoundaryWellFoundedFacePeelWitnessData_does_not_imply_nonempty_selectedBoundaryInteriorEdgeEndpointVertices_peeledEndpointTouch}
shows that even this earlier planarity-side witness endpoint can coexist with a
nonempty raw interior-edge endpoint carrier and an empty purified
selected-boundary carrier.  The packaged theorem
\path{peeledEndpointTouch_wellFoundedWitnessCarrierCruxObstructionSummary}
records that the graph already admits
\path{AdmitsPlanarBoundaryWellFoundedFacePeelWitnessData} while the purified
carrier is still exactly \(\emptyset\), and
\path{admitsPlanarBoundaryWellFoundedFacePeelWitnessData_withoutClosedWalkAnnulusBoundarySource_peeledEndpointTouchGraph}
shows that this earlier witness-cover endpoint still does not recover the
honest closed-walk annulus source on the same acyclic graph.
This collapse also survives on the direct attached-certificate
surface.  The endpoint-touch model now carries the explicit graph-level
certificate
\path{peeledEndpointTouchPlanarBoundaryAttachedRootedFacePeelCertificate}, but
\path{peeledEndpointTouch_attachedCertificate_emptyCarrier} shows that its own
purified selected-boundary endpoint carrier is still exactly \(\emptyset\).
The theorem
\path{planarBoundaryAttachedRootedFacePeelCertificate_does_not_imply_nonempty_selectedBoundaryInteriorEdgeEndpointVertices_peeledEndpointTouch}
packages the same separation in route-facing form: the raw interior-edge
endpoint carrier is nonempty, yet the attached-certificate endpoint alone does
not force a nonempty purified carrier.  And
\path{admitsPlanarBoundaryAttachedRootedFacePeelCertificate_withoutClosedWalkAnnulusBoundarySource_peeledEndpointTouchGraph}
shows that even the graph-level attached-certificate endpoint still does not
recover the honest closed-walk annulus source on this acyclic model.
The crux-directed summary theorem
\path{peeledEndpointTouch_nonvacuousPurifiedSynthesisCruxObstructionSummary}
packages this in the form most relevant to the proof route: the current peeled
annulus data exists and the raw interior-edge endpoint carrier is nonempty, but
the purified carrier is empty; consequently every non-vacuous purified
synthesis package is impossible on this model, even though the ordinary
purified synthesis package remains available for any Tait coloring.  This is a
robust obstruction to any continuation that silently reads the selected-boundary
purified endpoint theorem as a theorem about a nonempty annular interior
carrier.
The next crux certificate blocks a broader carrier workaround.  The theorem
\path{peeledEndpointTouch_graphEndpointSubcarrierCruxObstructionSummary}
shows that, in the same model, selected-boundary purification empties not only
the interior-edge endpoint carrier but also the canonical all-graph-endpoint
carrier \path{selectedBoundaryGraphInteriorVertices}.  More generally, every
finite candidate carrier contained in \path{graphEdgeEndpointSupport} purifies
to \(\emptyset\).  Thus this countermodel is not repaired by replacing the raw
interior-edge endpoint carrier with an arbitrary subcarrier of graph-edge
endpoints; any non-vacuous route must supply genuinely new geometry that
produces vertices outside this selected-boundary-covered endpoint set, or must
change the proof strategy.
This has now been sharpened to the arbitrary-carrier level for the same fixed
model.  The theorem
\path{peeledEndpointTouch_arbitraryCarrierCruxObstructionSummary} proves that
\path{graphEdgeEndpointSupport} is all of \(\mathrm{Fin}\,4\), and that the
selected-boundary endpoint support is also all of \(\mathrm{Fin}\,4\).  Hence
for every finite carrier \path{vertices : Finset (Fin 4)}, the purified set
\path{selectedBoundaryInteriorVertices peeledEndpointTouchEmbedding vertices}
is \(\emptyset\).  This closes carrier-only repairs inside the endpoint-touch
model: a non-vacuous continuation must either add a hypothesis excluding this
model, such as a genuine endpoint no-touch condition, or abandon this local
carrier strategy.
This endpoint-touch model is now also calibrated against the honest
boundary-order source.  Because \path{peeledEndpointTouchGraph} is just the
three-edge path, it is acyclic; the theorem
\path{not_nonempty_planarBoundaryClosedWalkEmbeddingData_peeledEndpointTouchGraph}
shows that no embedding of it can carry honest facial closed-walk geometry,
hence
\path{not_admitsPlanarBoundaryClosedWalkAnnulusBoundarySource_peeledEndpointTouchGraph}.
At the same time
\path{admitsPlanarBoundaryOuterBoundaryRootAdjDistancePeelData_peeledEndpointTouchGraph}
still holds, and
\path{admitsPlanarBoundaryOuterBoundaryRootAdjDistancePeelData_withoutClosedWalkAnnulusBoundarySource_peeledEndpointTouchGraph},
\path{admitsPlanarBoundaryHeightOrderedFacePeelWitnessData_withoutClosedWalkAnnulusBoundarySource_peeledEndpointTouchGraph},
\path{admitsPlanarBoundaryWellFoundedFacePeelWitnessData_withoutClosedWalkAnnulusBoundarySource_peeledEndpointTouchGraph},
and
\path{admitsPlanarBoundaryAttachedRootedFacePeelCertificate_withoutClosedWalkAnnulusBoundarySource_peeledEndpointTouchGraph}
make the calibration explicit: the endpoint-touch countermodel blocks later
peeled-route and certificate consequences, but it is not itself a same-model
counterexample to the honest closed-walk source route.

The raw carrier is now connected as far as the current route can honestly take
it.  The lemma
\path{interiorEdgeEndpointSupport_avoids_selectedBoundarySupport_of_no_endpoint_shared}
turns endpoint-level separation into the boundary-erasure side condition, and
\path{theorem49BoundaryZeroKirchhoffSubspace_interiorEdgeEndpointSupport_image_and_separation_of_planarBoundaryOuterBoundaryRootAdjDistancePeelData}
then gives the boundary-erasure image identity and projected-generator
separation theorem for the unpurified
\path{interiorEdgeEndpointSupport}.  The graph-level wrapper
\path{exists_theorem49BoundaryZeroKirchhoffSubspace_interiorEdgeEndpointSupport_image_and_separation_of_admitsPlanarBoundaryOuterBoundaryRootAdjDistancePeelData}
exposes the same statement from admitted outer-boundary-rooted annulus data,
parameterized by the missing endpoint-separation proof.

The same route is now stated in the cleaner finite-support language.  The
theorem
\path{selectedBoundaryInteriorEdgeEndpointVertices_eq_interiorEdgeEndpointSupport_iff_endpointSupport_disjoint}
identifies the raw-carrier condition with disjointness between
\path{interiorEdgeEndpointSupport} and the endpoint support of
\path{selectedBoundarySupport}.  The route theorem
\path{theorem49BoundaryZeroKirchhoffSubspace_interiorEdgeEndpointSupport_image_and_separation_of_planarBoundaryOuterBoundaryRootAdjDistancePeelData_of_endpointSupport_disjoint}
uses this endpoint-support disjointness directly, with graph-level wrapper
\path{exists_theorem49BoundaryZeroKirchhoffSubspace_interiorEdgeEndpointSupport_image_and_separation_of_admitsPlanarBoundaryOuterBoundaryRootAdjDistancePeelData_of_endpointSupport_disjoint}.

The endpoint-separation obligation is now also present as a named annulus-route
data surface rather than only as an external parameter.  The structure
\path{PlanarBoundaryOuterBoundaryRootAdjDistancePeelDataWithEndpointSeparation}
extends the outer-boundary-rooted annulus data by exactly the finite
endpoint-support disjointness needed for the raw
\path{interiorEdgeEndpointSupport} carrier.  The theorem
\path{theorem49BoundaryZeroKirchhoffSubspace_interiorEdgeEndpointSupport_image_and_separation_of_planarBoundaryOuterBoundaryRootAdjDistancePeelDataWithEndpointSeparation}
then gives the raw-carrier image identity and projected-generator separation
from this packaged route data, with a matching graph-level wrapper
\path{exists_theorem49BoundaryZeroKirchhoffSubspace_interiorEdgeEndpointSupport_image_and_separation_of_admitsPlanarBoundaryOuterBoundaryRootAdjDistancePeelDataWithEndpointSeparation}.
This package now also yields the first honest positive non-vacuity repair on
the endpoint side.  The theorem
\path{PlanarBoundaryOuterBoundaryRootAdjDistancePeelDataWithEndpointSeparation.selectedBoundaryInteriorEdgeEndpointVertices_eq_interiorEdgeEndpointSupport}
proves that the packaged endpoint-separation field makes selected-boundary
purification lossless on the raw interior-edge endpoint carrier.  Therefore
\path{theorem49OuterBoundaryRootNonemptySynthesis_of_planarBoundaryOuterBoundaryRootAdjDistancePeelDataWithEndpointSeparation}
shows that once the same data also has a nonempty raw
\path{interiorEdgeEndpointSupport}, the current outer-boundary-rooted route
really does close to the non-vacuous theorem-4.9 synthesis endpoint.  The
graph-level wrapper
\path{exists_theorem49OuterBoundaryRootNonemptySynthesis_of_exists_planarBoundaryOuterBoundaryRootAdjDistancePeelDataWithEndpointSeparation_with_nonempty_interiorEdgeEndpointSupport}
makes the repaired target explicit.

There is now a first construction-level no-touch bridge into this endpoint
condition.  In
\path{Theorem49InteriorVerticesAnnulusConstruction.lean}, the theorem
\path{PlanarBoundaryAnnulusConstruction.endpointSupport_disjoint_of_peelFaces_endpoint_disjoint_selectedBoundarySupport}
uses the annulus construction's cover of
\path{interiorEdgeSupport} by peeled-face witness edges: if every peeled
face boundary is endpoint-disjoint from the selected-boundary support, then
\path{interiorEdgeEndpointSupport} is endpoint-disjoint from the
selected-boundary endpoint support.  The theorem
\path{PlanarBoundaryOuterBoundaryRootAdjDistancePeelData.selectedBoundaryInteriorEdgeEndpointVertices_eq_interiorEdgeEndpointSupport_of_peelFaces_endpoint_disjoint_selectedBoundarySupport}
lifts the same fact to the outer-boundary-rooted annulus data surface, and
\path{exists_selectedBoundaryInteriorEdgeEndpointVertices_eq_interiorEdgeEndpointSupport_of_admitsPlanarBoundaryOuterBoundaryRootAdjDistancePeelData_of_peelFaces_endpoint_disjoint_selectedBoundarySupport}
records the graph-level lossless-purification consequence.  The theorem
\path{PlanarBoundaryOuterBoundaryRootAdjDistancePeelData.withEndpointSeparationOfPeelFacesEndpointDisjoint}
then packages this construction-level no-touch proof as
\path{PlanarBoundaryOuterBoundaryRootAdjDistancePeelDataWithEndpointSeparation},
with graph-level constructor
\path{admitsPlanarBoundaryOuterBoundaryRootAdjDistancePeelDataWithEndpointSeparation_of_admitsPlanarBoundaryOuterBoundaryRootAdjDistancePeelData_of_peelFaces_endpoint_disjoint_selectedBoundarySupport}.
Finally,
\path{theorem49BoundaryZeroKirchhoffSubspace_interiorEdgeEndpointSupport_image_and_separation_of_planarBoundaryOuterBoundaryRootAdjDistancePeelData_of_peelFaces_endpoint_disjoint_selectedBoundarySupport}
therefore gives the raw-carrier \(W_0(H)\) image identity and separation
endpoint directly from a peeled-face no-touch obligation, and the graph-level
wrapper
\path{exists_theorem49BoundaryZeroKirchhoffSubspace_interiorEdgeEndpointSupport_image_and_separation_of_admitsPlanarBoundaryOuterBoundaryRootAdjDistancePeelData_of_peelFaces_endpoint_disjoint_selectedBoundarySupport}
exposes the same endpoint from admitted annulus data plus the per-datum
no-touch proof.  This does not yet
prove the manuscript's geometry; it identifies the concrete geometric
condition that would discharge the raw carrier without falling back to
selected-boundary purification.

The positive route has now been pushed one layer further.  In
\path{Theorem49Synthesis.lean}, the theorem
\path{theorem49OuterBoundaryRootNonemptySynthesis_of_planarBoundaryOuterBoundaryRootAdjDistancePeelData_of_peelFaces_endpoint_disjoint_selectedBoundarySupport}
shows that the same peeled-face endpoint no-touch hypothesis, together with a
nonempty raw \path{interiorEdgeEndpointSupport}, already yields the
non-vacuous outer-boundary-root Theorem~4.9 synthesis package.  Then
\path{PlanarBoundaryClosedWalkSourceRoute.lean} adds the honest source-facing
wrapper
\path{theorem49OuterBoundaryRootNonemptySynthesis_of_closedWalkAnnulusBoundarySource_and_interiorDualBoundaryRootAdjDistancePeelData_of_rootsOuterAmbientBoundary_and_peelFaces_endpoint_disjoint_selectedBoundarySupport}
and its graph-level existential form, so the current positive burden is now
precise at the closed-walk source surface: prove root localization to the
outer ambient boundary, peeled-face endpoint no-touch, and raw-carrier
nonemptiness on one embedding.

That positive bridge is now calibrated honestly as well.  The same file proves
\path{not_exists_closedWalkAnnulusBoundarySource_and_interiorDualBoundaryRootAdjDistancePeelData_and_rootsOuterAmbientBoundary_and_peelFaces_endpoint_disjoint_selectedBoundarySupport_and_nonempty_interiorEdgeEndpointSupport}
and
\path{not_exists_closedWalkAnnulusBoundarySource_and_theorem49OuterBoundaryRootNonemptySynthesis_sameBoundaryData}.
So on the honest same-boundary split extracted from the closed-walk source,
the new outer-root non-vacuity theorem is formally vacuous: the required
all-roots-on-the-extracted-outer-boundary witness already contradicts the
closed-walk source's boundary separation, and therefore the resulting
outer-boundary-root non-vacuous synthesis endpoint cannot coexist with that
same split either.

The honest positive route therefore shifts to the two-sided annulus-root
surface.  \path{Theorem49Synthesis.lean} now proves
\path{theorem49AnnulusRootNonemptySynthesis_of_planarBoundaryAnnulusRootAdjDistancePeelData_of_endpointSupport_disjoint}
and the graph-level wrapper
\path{exists_theorem49AnnulusRootNonemptySynthesis_of_exists_planarBoundaryAnnulusRootAdjDistancePeelData_with_endpointSupport_disjoint_and_nonempty_interiorEdgeEndpointSupport}:
for annulus-root data, non-vacuity already follows from nonempty raw
\path{interiorEdgeEndpointSupport} plus endpoint-support disjointness from the
selected boundary.  Then \path{PlanarBoundaryClosedWalkSourceRoute.lean} adds
the honest source-facing theorem
\path{theorem49AnnulusRootNonemptySynthesis_of_closedWalkAnnulusBoundarySource_and_interiorDualBoundaryRootAdjDistancePeelData_of_endpointSupport_disjoint}
and its graph-level existential form.  This is the first non-vacuous theorem
on a boundary-root surface that survives on the honest closed-walk source shell
without any all-roots-on-the-outer-boundary hypothesis.

That annulus-root repair theorem is now calibrated on its own route surface as
well.  \path{Theorem49InteriorVerticesEndpointObstruction.lean} proves
\path{annulusRootAdjDistancePeelData_does_not_imply_endpointSupport_disjoint_peeledEndpointTouch}
and the fixed-embedding universal form
\path{not_forall_annulusRootAdjDistancePeelData_implies_endpointSupport_disjoint_peeledEndpointTouch}.
So the new two-sided annulus-root bridge is an honest conditional repair, not
something already forced by arbitrary current annulus-root data.

The same missing burden now has an honest geometric calibration too.
\path{Theorem49PlanarBoundaryAnnulusHonestGeometryRegression.lean} proves
\path{closedWalkAnnulusBoundarySource_and_annulusCollarGeometry_and_attachedRootedFacePeelCertificate_with_nonempty_interiorEdgeEndpointSupport_without_endpointSupport_disjoint_diamondWithTriangle}:
on the genuine cyclic source model \path{diamondWithTriangleEmbedding}, the
honest closed-walk annulus source, a weak one-collar annulus geometry, the
attached rooted face-peel certificate endpoint, and a nonempty raw
\path{interiorEdgeEndpointSupport} all coexist, yet endpoint-support
disjointness from the selected boundary still fails.  So the present
endpoint-separation repair lies strictly later than the current honest
annulus-geometry and certificate lane as well; it is not already forced by
reaching those earlier geometric surfaces.  The same file now strengthens this
to the repaired geometry surface itself:
\path{closedWalkAnnulusBoundarySource_and_previousBoundaryWitnessGeometry_and_wellFoundedFacePeelWitnessData_with_nonempty_interiorEdgeEndpointSupport_without_endpointSupport_disjoint_diamondWithTriangle}.
So even the explicit previous-boundary witness placement needed to feed the
boundary-aware well-founded peel endpoint still does not force the endpoint
separation required by the non-vacuous annulus-root repair.  More sharply, the
same honest source model already satisfies
\path{closedWalkAnnulusBoundarySource_and_previousBoundaryWitnessGeometry_and_wellFoundedFacePeelWitnessData_with_nonempty_interiorEdgeEndpointSupport_and_empty_selectedBoundaryInteriorEdgeEndpointVertices_diamondWithTriangle}:
the raw endpoint carrier is nonempty, but the purified carrier
\path{selectedBoundaryInteriorEdgeEndpointVertices} is exactly empty.  So the
current non-vacuous repair is not merely missing a disjointness lemma; it is
missing a real geometric exclusion principle on an honest cyclic source model.
This now persists one layer higher as well: the same file explicitly builds the
collar-layer witness package
\path{diamondWithTriangleOneCollarLayerFacePeelWitnessData}, lowers it to the
height-ordered peel witness
\path{diamondWithTriangleOneCollarHeightOrderedFacePeelWitnessData}, and proves
\path{closedWalkAnnulusBoundarySource_and_previousBoundaryWitnessGeometry_and_heightOrderedFacePeelWitnessData_with_nonempty_interiorEdgeEndpointSupport_and_empty_selectedBoundaryInteriorEdgeEndpointVertices_diamondWithTriangle}.
So even the explicit layered peel surface still does not force a nonempty
purified endpoint carrier on an honest cyclic source model.

The remaining honest burden is therefore sharper: either prove the manuscript's
intended annular interior vertices avoid selected-boundary endpoints, or make
the purification step explicit in the paper-facing statement.

Independently of that non-vacuity calibration, the positive Theorem~4.9
synthesis layer has now been pushed earlier and made route-independent above
the boundary-aware certificate surface.  In
\path{Theorem49Synthesis.lean}, the theorem
\path{theorem49BoundaryRootSynthesis_of_boundaryZeroAnnihilatorTrivial}
packages the full purified synthesis endpoint from the single algebraic
hypothesis that every selected-boundary-zero annihilator of the Definition~4.8
family vanishes.  The same file then proves
\path{theorem49BoundaryRootSynthesis_of_planarBoundaryAttachedRootedFacePeelCertificate},
\path{theorem49BoundaryRootSynthesis_of_planarBoundaryWellFoundedFacePeelWitnessData},
and the graph-level wrappers
\path{exists_theorem49BoundaryRootSynthesis_of_admitsPlanarBoundaryAttachedRootedFacePeelCertificate}
and
\path{exists_theorem49BoundaryRootSynthesis_of_admitsPlanarBoundaryWellFoundedFacePeelWitnessData}.
The older interior-dual and annulus-construction synthesis theorems now factor
through this same generic bridge.  So once any honest route reaches the
attached certificate or well-founded peel witness surface, no further
Theorem~4.9 algebra remains; the live burden is entirely upstream geometric
construction of that certificate data.  More sharply, the same file now proves
\path{theorem49BoundaryRootSynthesis_of_planarBoundaryHeightOrderedFacePeelWitnessData},
\path{theorem49BoundaryRootSynthesis_of_planarBoundaryCollarLayerFacePeelWitnessData},
\path{theorem49BoundaryRootSynthesis_of_planarBoundaryLocalRemainderBoundaryCollarLayerFacePeelWitnessData},
\path{theorem49BoundaryRootSynthesis_of_planarBoundaryAnnulusWitnessAssignment},
\path{exists_theorem49BoundaryRootSynthesis_of_admitsPlanarBoundaryHeightOrderedFacePeelWitnessData},
\path{exists_theorem49BoundaryRootSynthesis_of_admitsPlanarBoundaryCollarLayerFacePeelWitnessData},
\path{exists_theorem49BoundaryRootSynthesis_of_admitsPlanarBoundaryLocalRemainderBoundaryCollarLayerFacePeelWitnessData},
and
\path{exists_theorem49BoundaryRootSynthesis_of_admitsPlanarBoundaryAnnulusWitnessAssignment},
so theorem-4.9 synthesis already starts at the boundary-aware height-ordered
witness-cover surface and even at the weaker finite collar-layer witness-cover
surface, hence also at the explicit local remainder annulus surface and at the
split ``annulus decomposition + witness assignment'' surface.
\path{Theorem49PositiveGeometricSourceReplacement.lean} now also proves
\path{theorem49BoundaryRootSynthesis_of_collarLayerPositiveProjectedGeometryOn},
\path{Theorem49HeightOrderedPositiveProjectedGeometry.boundaryRootSynthesis},
\path{Theorem49HeightOrderedPositiveProjectedGeometry.exists_boundaryRootSynthesis},
and the matching collar-layer graph-package pair.  So once the route already
reaches either of those named positive witness-cover packages, full
\path{Theorem49BoundaryRootSynthesis} is immediate; the projected endpoint and
raw quotient/error package are now explicit projections of that stronger
endpoint rather than the top of those surfaces.
\path{admitsPlanarBoundaryAttachedRootedFacePeelCertificate_of_admitsPlanarBoundaryAnnulusConstructionParentWitnessOrientation},
so the root-distance parent-witness surface already lands on the certificate
surface explicitly.  The honest closed-walk and successor-cycle source
corollaries in \path{PlanarBoundaryClosedWalkSourceRoute.lean} now factor their
direct certificate and synthesis endpoints through that route instead of
dropping straight from interior-dual admits to synthesis.  The same source file
now also proves
\path{exists_theorem49BoundaryRootSynthesis_of_exists_closedWalkAnnulusBoundarySource_and_heightOrderedFacePeelWitnessData_direct},
\path{exists_theorem49BoundaryRootSynthesis_of_exists_closedWalkAnnulusBoundarySource_and_collarLayerFacePeelWitnessData_direct},
\path{exists_theorem49BoundaryRootSynthesis_of_exists_closedWalkAnnulusBoundarySource_and_annulusWitnessAssignment_direct}
and
\path{exists_theorem49BoundaryRootSynthesis_of_exists_boundaryReachabilityData_and_dartSuccessorCycleEmbeddingData_and_selectedBoundaryArc_and_heightOrderedFacePeelWitnessData_direct},
\path{exists_theorem49BoundaryRootSynthesis_of_exists_boundaryReachabilityData_and_dartSuccessorCycleEmbeddingData_and_selectedBoundaryArc_and_collarLayerFacePeelWitnessData_direct},
\path{exists_theorem49BoundaryRootSynthesis_of_exists_boundaryReachabilityData_and_dartSuccessorCycleEmbeddingData_and_selectedBoundaryArc_and_annulusWitnessAssignment_direct},
so the honest source shell already reaches the full synthesis endpoint as soon
as the same embedding carries either boundary-aware height-sorted witness-cover
data, the weaker finite collar-layer witness-cover package, or the later pure
annulus decomposition together with witness-edge ownership.
The same route file now also proves the matching direct Step~2 vanishing
corollaries
\path{zero_on_allEdges_of_exists_closedWalkAnnulusBoundarySource_and_heightOrderedFacePeelWitnessData_direct_of_annihilatesKempeClosureGeneratorFamily},
\path{zero_on_allEdges_of_exists_closedWalkAnnulusBoundarySource_and_collarLayerFacePeelWitnessData_direct_of_annihilatesKempeClosureGeneratorFamily},
\path{zero_on_allEdges_of_exists_boundaryReachabilityData_and_dartSuccessorCycleEmbeddingData_and_selectedBoundaryArc_and_heightOrderedFacePeelWitnessData_direct_of_annihilatesKempeClosureGeneratorFamily},
and
\path{zero_on_allEdges_of_exists_boundaryReachabilityData_and_dartSuccessorCycleEmbeddingData_and_selectedBoundaryArc_and_collarLayerFacePeelWitnessData_direct_of_annihilatesKempeClosureGeneratorFamily}.
So the positive algebraic boundary now matches the positive synthesis boundary:
once an honest source shell reaches boundary-aware height-sorted witness-cover
data, or even just the finite collar-layer witness-cover package, no later
annulus-geometry packaging is needed either for Step~2 vanishing or for the
purified theorem-4.9 endpoint.
\path{theorem49BoundaryRootSynthesis_of_planarBoundaryAnnulusCollarGeometry},
\path{theorem49BoundaryRootSynthesis_of_planarBoundaryAnnulusPreviousBoundaryWitnessGeometry},
and the graph-level wrappers
\path{exists_theorem49BoundaryRootSynthesis_of_admitsPlanarBoundaryAnnulusCollarGeometry},
\path{exists_theorem49BoundaryRootSynthesis_of_admitsPlanarBoundaryAnnulusPreviousBoundaryWitnessGeometry},
and
\path{exists_theorem49BoundaryRootSynthesis_of_exists_planarBoundaryAnnulusCollarGeometry_and_witnessOnCurrentBoundary}.
So for synthesis itself, the burden has now moved even earlier than the weak
bundled collar geometry surface: once the honest route can produce the
boundary-aware height-ordered witness-cover surface, or even the weaker finite
collar-layer witness-cover package, the boundary-aware certificate and the full
current Theorem~4.9 endpoint are automatic.  The live positive burden is now
strictly upstream of explicit remainder bookkeeping and of annulus witness
ownership themselves.
That upstream geometric burden is now exposed directly in Lean.
\path{Theorem49PlanarBoundaryAnnulusGeometry.lean} adds
\path{PlanarBoundaryAnnulusCollarGeometry.toPlanarBoundaryCollarLayerFacePeelWitnessData},
\path{PlanarBoundaryAnnulusCollarGeometry.toPlanarBoundaryHeightOrderedFacePeelWitnessData},
\path{admitsPlanarBoundaryCollarLayerFacePeelWitnessData_of_admitsPlanarBoundaryAnnulusCollarGeometry},
and
\path{admitsPlanarBoundaryHeightOrderedFacePeelWitnessData_of_admitsPlanarBoundaryAnnulusCollarGeometry},
so weak annulus collar geometry canonically constructs the current
witness-cover surfaces.  The same file also adds
\path{PlanarBoundaryAnnulusPreviousBoundaryWitnessGeometry.toPlanarBoundaryHeightOrderedFacePeelWitnessData},
\path{admitsPlanarBoundaryHeightOrderedFacePeelWitnessData_of_admitsPlanarBoundaryAnnulusPreviousBoundaryWitnessGeometry},
and
\path{admitsPlanarBoundaryWellFoundedFacePeelWitnessData_of_exists_planarBoundaryAnnulusCollarGeometry_and_witnessOnCurrentBoundary},
so the repaired witness-on-current-boundary hypothesis lands directly on the
well-founded peel lane.
\path{PlanarBoundaryClosedWalkSourceRoute.lean} then exposes these bridges on
the honest source shells through
\path{admitsPlanarBoundaryCollarLayerFacePeelWitnessData_of_exists_closedWalkAnnulusBoundarySource_and_annulusCollarGeometry_direct},
\path{admitsPlanarBoundaryHeightOrderedFacePeelWitnessData_of_exists_closedWalkAnnulusBoundarySource_and_annulusCollarGeometry_direct},
\path{admitsPlanarBoundaryWellFoundedFacePeelWitnessData_of_exists_closedWalkAnnulusBoundarySource_and_annulusCollarGeometry_and_witnessOnCurrentBoundary_direct},
and the successor-cycle companions
\path{admitsPlanarBoundaryCollarLayerFacePeelWitnessData_of_exists_boundaryReachabilityData_and_dartSuccessorCycleEmbeddingData_and_selectedBoundaryArc_and_annulusCollarGeometry_direct},
\path{admitsPlanarBoundaryHeightOrderedFacePeelWitnessData_of_exists_boundaryReachabilityData_and_dartSuccessorCycleEmbeddingData_and_selectedBoundaryArc_and_annulusCollarGeometry_direct},
and
\path{admitsPlanarBoundaryWellFoundedFacePeelWitnessData_of_exists_boundaryReachabilityData_and_dartSuccessorCycleEmbeddingData_and_selectedBoundaryArc_and_annulusCollarGeometry_and_witnessOnCurrentBoundary_direct}.
So the remaining live burden is now stated cleanly as a geometric one: derive
same-embedding annulus collar geometry from honest source hypotheses, and, if
one wants the repaired well-founded lane, also derive the current-boundary
witness placement.
Lean now packages one concrete sufficient criterion for that first geometric
step on the canonical one-collar route.
\path{Theorem49PlanarBoundaryAnnulusWitness.lean} defines
\path{PlanarBoundaryCanonicalWitnessChoice boundaryData}: for each ambient face
of a fixed annulus boundary split, choose one boundary edge as witness so that
all interior edges are chosen by some face and every non-witness boundary edge
already lies on the outer or inner ambient annulus boundary.
\path{PlanarBoundaryCanonicalWitnessChoice.toPlanarBoundaryAnnulusWitnessAssignment}
then upgrades the canonical decomposition
\path{planarBoundaryAnnulusDecomposition_of_boundaryData boundaryData} to a full
annulus witness assignment on the same embedding.
\path{PlanarBoundaryClosedWalkSourceRoute.lean} packages the source-facing
consequence as
\path{admitsPlanarBoundaryAnnulusCollarGeometry_of_exists_closedWalkAnnulusBoundarySource_and_canonicalWitnessChoice_direct}.
So a surviving honest one-collar route no longer needs an opaque annulus
geometry black box: it is enough to construct this explicit facewise canonical
witness choice on the source boundary split.
This has now been upgraded from graph-level admissibility to actual
one-collar geometry on the same embedding.  In
\path{Theorem49PlanarBoundaryAnnulusGeometry.lean},
\path{PlanarBoundaryCanonicalWitnessChoice.toPlanarBoundaryAnnulusCollarGeometry}
constructs a \path{PlanarBoundaryAnnulusCollarGeometry} directly from any
canonical witness choice on any annulus boundary split.  The companion lemmas
\path{PlanarBoundaryCanonicalWitnessChoice.toPlanarBoundaryAnnulusCollarGeometry_numCollars}
and
\path{PlanarBoundaryCanonicalWitnessChoice.toPlanarBoundaryAnnulusCollarGeometry_boundaryData}
prove that the constructed geometry has exactly one collar and preserves the
original boundary data.  The honest-source wrapper
\path{exists_oneCollarAnnulusCollarGeometry_with_sourceBoundaryData_of_exists_closedWalkAnnulusBoundarySource_and_canonicalWitnessChoice}
therefore proves one-collar geometry not merely on the diamond example but on
the full class of honest closed-walk sources whose extracted boundary split
carries a canonical facewise witness choice.
The canonical witness choice itself now has a checked local sufficient
condition.  \path{PlanarBoundaryCanonicalWitnessChoice.of_atMostOneInteriorEdgePerFace}
chooses, on each face, the unique interior boundary edge when one exists and a
supplied fallback boundary edge otherwise.  Its hypotheses are exactly the
facewise geometry needed for the canonical choice: at most one interior edge on
each face, and every non-interior edge on a face boundary lies on the outer or
inner ambient annulus boundary.  Lean now proves the second clause directly
from the source's annulus boundary split:
\path{PlanarBoundaryAnnulusBoundaryData.mem_outerAmbientBoundary_union_innerAmbientBoundary_of_mem_faceBoundary_of_not_mem_interiorEdgeSupport}
uses the incidence partition of ambient face-boundary edges into selected
boundary support and interior support.  Lean now also discharges the fallback
edge from the honest closed-walk field itself:
\path{PlaneEmbeddingWithBoundary.FaceBoundaryClosedWalk.faceBoundary_nonempty}
shows that an honest nonempty facial boundary walk has a nonempty face-boundary
edge set, and
\path{PlanarBoundaryClosedWalkAnnulusBoundarySource.fallbackEdge} chooses such
an edge with its membership proof.  Consequently the new geometry wrapper
\path{exists_oneCollarAnnulusCollarGeometry_with_sourceBoundaryData_of_exists_closedWalkAnnulusBoundarySource_and_facewiseUniqueInteriorEdge}
constructs same-embedding one-collar collar geometry from an honest source plus
only the equality-style facewise uniqueness of interior boundary edges.  The
source-route cardinal variant
\path{exists_planarBoundaryCanonicalWitnessChoice_of_closedWalkAnnulusBoundarySource_and_facewiseAtMostOneInteriorEdge}
packages the more convenient finite-cardinality at-most-one hypothesis into the
same canonical witness-choice surface, while the older
\path{exists_planarBoundaryCanonicalWitnessChoice_of_closedWalkAnnulusBoundarySource_and_atMostOneInteriorEdgePerFace}
records the fully explicit fallback/count/boundary-rest package.
The geometry file now carries the corresponding fixed-embedding collar result:
\path{exists_oneCollarAnnulusCollarGeometry_of_closedWalkAnnulusBoundarySource_and_facewiseAtMostOneInteriorEdge}
constructs a source-preserving one-collar
\path{PlanarBoundaryAnnulusCollarGeometry} directly from an honest source and
the facewise cardinal at-most-one hypothesis.  No separate fallback or
non-interior boundary-rest hypothesis remains at this layer; those are supplied
by the honest source and its extracted annulus boundary split.  Lean also
exposes the graph-level source-facing endpoint
\path{exists_oneCollarAnnulusCollarGeometry_with_sourceBoundaryData_of_exists_closedWalkAnnulusBoundarySource_and_facewiseAtMostOneInteriorEdgeCardinality},
which starts from an existential honest closed-walk annulus source plus the same
facewise cardinal bound and returns a one-collar collar geometry preserving that
source boundary split on the chosen embedding.
The one-collar branch now also reaches the repaired
previous-boundary-witness surface.  The generic theorem
\path{admitsPlanarBoundaryAnnulusPreviousBoundaryWitnessGeometry_of_exists_oneCollarAnnulusCollarGeometry}
packages the fact that the positive-collar witness-placement obligation is
vacuous when there is only one collar.  The honest-source geometry theorem
\path{admitsPlanarBoundaryAnnulusPreviousBoundaryWitnessGeometry_of_exists_closedWalkAnnulusBoundarySource_and_facewiseUniqueInteriorEdge}
and the cardinal route theorem
\path{admitsPlanarBoundaryAnnulusPreviousBoundaryWitnessGeometry_of_exists_closedWalkAnnulusBoundarySource_and_facewiseAtMostOneInteriorEdge_direct}
therefore reduce the surviving source-tied one-collar burden to the
at-most-one/unique-interior-edge condition alone.  Thus the at-most-one branch
no longer stops at weak collar geometry; it can feed the corrected
well-founded, height-ordered, and collar-layer witness surfaces.
Lean now also exposes this repair at the primitive canonical-choice layer.
\path{PlanarBoundaryCanonicalWitnessChoice.toPlanarBoundaryAnnulusPreviousBoundaryWitnessGeometry}
turns a facewise canonical witness choice on an annulus boundary split directly
into \path{PlanarBoundaryAnnulusPreviousBoundaryWitnessGeometry}; internally it
uses the canonical one-collar geometry, so the positive-collar witness condition
is empty.  The source-facing theorem
\path{exists_planarBoundaryAnnulusPreviousBoundaryWitnessGeometry_of_exists_closedWalkAnnulusBoundarySource_and_canonicalWitnessChoice}
records the corresponding honest closed-walk source package.  This is a
positive statement about the canonical one-collar geometry built from the
source boundary split, not a theorem that an arbitrary same-embedding weak
collar satisfies the repaired placement condition; the latter is still exactly
the explicit \path{WitnessOnCurrentBoundary} obligation.  The route-facing
admission theorem
\path{admitsPlanarBoundaryAnnulusPreviousBoundaryWitnessGeometry_of_exists_closedWalkAnnulusBoundarySource_and_canonicalWitnessChoice_direct}
now uses this repaired canonical-choice geometry theorem directly, rather than
reconstructing the same landing through the intermediate one-collar existence
surface.
The same local criterion now reaches the theorem-4.9 synthesis endpoint
directly.  In \path{PlanarBoundaryClosedWalkSourceRoute.lean},
\path{exists_planarBoundaryCanonicalWitnessChoice_of_closedWalkAnnulusBoundarySource_and_facewiseAtMostOneInteriorEdge}
factors the cardinal at-most-one proof into the source route without caller
fallback or boundary-rest assumptions, and
\path{exists_theorem49BoundaryRootSynthesis_of_exists_closedWalkAnnulusBoundarySource_and_facewiseAtMostOneInteriorEdge_direct}
composes it with the canonical witness-choice route to produce
\path{Theorem49BoundaryRootSynthesis} for any Tait coloring.
The nonempty-carrier version is now isolated at the same route surface in the
same caller-clean cardinal form:
\path{theorem49BoundaryRootNonemptySynthesis_of_closedWalkAnnulusBoundarySource_and_facewiseAtMostOneInteriorEdge_direct}
adds only the explicit hypothesis
\path{(selectedBoundaryInteriorEdgeEndpointVertices emb).Nonempty}, and
\path{exists_theorem49BoundaryRootNonemptySynthesis_of_exists_closedWalkAnnulusBoundarySource_and_facewiseAtMostOneInteriorEdge_with_nonempty_selectedBoundaryInteriorEdgeEndpointVertices_direct}
packages the corresponding graph-level endpoint.  Thus the current positive
nonempty synthesis burden is exactly the honest source, facewise cardinal
at-most-one, and a surviving selected-boundary-purified carrier; no fallback
edge or boundary-rest datum remains caller supplied.
The same source package is now connected to the canonical Definition~4.8
spanning bridge, not only to the projected synthesis endpoint.  The direct
wrappers
\path{exists_theorem49_kempeGeneratorSubspace_spanning_of_exists_closedWalkAnnulusBoundarySource_and_atMostOneInteriorEdgePerFace_heightOrdered_direct}
and
\path{exists_theorem49_kempeGeneratorSubspace_spanning_of_exists_closedWalkAnnulusBoundarySource_and_atMostOneInteriorEdgePerFace_collarLayer_direct}
compose the repaired one-collar source route with the canonical
\path{Theorem49SubmoduleDualityData} constructor from
\path{Theorem49KempeGeneratorSpan.lean}.  Consequently, after the local
at-most-one source package is established, the remaining spanning obligations
are exactly the named target subspace, containment of the raw Definition~4.8
generator span, and the target's selected-boundary-zero interpretation.
The reverse extraction from a stronger one-collar geometry surface is now also
formalized.  In \path{Theorem49PlanarBoundaryAnnulusGeometry.lean},
\path{PlanarBoundaryAnnulusCollarGeometry.toPlanarBoundaryCanonicalWitnessChoice_of_numCollars_eq_one}
shows that any one-collar weak annulus geometry whose boundary split is a fixed
annulus boundary split already contains the corresponding
\path{PlanarBoundaryCanonicalWitnessChoice}; with only one collar there is no
intermediate frontier, so every non-witness edge is forced onto the outer or
inner ambient boundary.  The source-facing wrapper
\path{exists_planarBoundaryCanonicalWitnessChoice_of_exists_closedWalkAnnulusBoundarySource_and_oneCollarAnnulusCollarGeometry_with_sourceBoundaryData}
records the exact honest-source form: one-collar collar geometry preserving the
boundary split extracted from the source supplies the canonical witness-choice
package needed by the direct route theorem.
\path{PlanarBoundaryClosedWalkSourceRoute.lean} now carries this all the way
to the theorem-4.9 synthesis endpoint:
\path{exists_theorem49BoundaryRootSynthesis_of_exists_closedWalkAnnulusBoundarySource_and_canonicalWitnessChoice_direct}
uses the canonical witness choice to build the split annulus witness assignment
consumed by the synthesis layer, and
\path{exists_theorem49BoundaryRootSynthesis_of_exists_closedWalkAnnulusBoundarySource_and_oneCollarAnnulusCollarGeometry_with_sourceBoundaryData}
first extracts that canonical witness choice from same-boundary one-collar
geometry.  Thus the source-facing positive burden can be stated either as a
canonical facewise witness choice or as a stronger one-collar collar geometry
preserving the source boundary split; either package now reaches the current
formal theorem-4.9 endpoint directly.
The source route now pushes these two primitive one-collar surfaces into the
canonical Definition~4.8 spanning lane as well.  The direct wrappers
\path{exists_theorem49_kempeGeneratorSubspace_spanning_of_exists_closedWalkAnnulusBoundarySource_and_canonicalWitnessChoice_heightOrdered_direct}
and
\path{exists_theorem49_kempeGeneratorSubspace_spanning_of_exists_closedWalkAnnulusBoundarySource_and_canonicalWitnessChoice_collarLayer_direct}
compose canonical source witness choice with the repaired previous-boundary
surface and the height-ordered/collar-layer generator-span bridge.  The
corresponding one-collar geometry wrappers
\path{exists_theorem49_kempeGeneratorSubspace_spanning_of_exists_closedWalkAnnulusBoundarySource_and_oneCollarAnnulusCollarGeometry_with_sourceBoundaryData_heightOrdered_direct}
and
\path{exists_theorem49_kempeGeneratorSubspace_spanning_of_exists_closedWalkAnnulusBoundarySource_and_oneCollarAnnulusCollarGeometry_with_sourceBoundaryData_collarLayer_direct}
do the same when the source already carries one-collar geometry preserving its
boundary split.  Thus the theorem-4.9 submodule-duality data no longer has to
be reconstructed separately for any of the source-side one-collar packages:
the remaining spanning inputs are precisely the target subspace, raw-generator
containment, and boundary-zero interpretation.
\path{PlanarBoundaryClosedWalkSourceProjection.lean} records the corrected
boundary-zero version of this source bridge.  Since the raw
Definition~4.8 generator span is not, in general, a subspace of the selected
boundary-zero chains, the new wrappers use the boundary-erased generator
source
\path{projectedKempeClosureGeneratorSubspace}.  Canonical witness choice,
the local at-most-one source package, and same-boundary one-collar geometry
now each imply both the subspace identity
\[
  \mathrm{projectedKempeClosureGeneratorSubspace}
  =
  \mathrm{planarBoundaryZeroSubmodule}
\]
and the explicit-family span version, through either the height-ordered or
finite collar-layer witness surface.  This exposes a concrete source-side
spanning conclusion with no caller-supplied target subspace or raw-generator
containment hypothesis.  The same projected boundary-zero endpoint has now
been lifted to the route-facing successor-cycle source shell.  The lowering
lemmas in the same file convert successor-cycle boundary data plus selected
arc contiguity into the honest closed-walk source package while preserving
canonical witness choice, the local at-most-one package, or same-boundary
one-collar geometry.  The resulting successor-cycle theorems, such as
\path{exists_span_projectedKempeClosureGeneratorFamily_eq_planarBoundaryZeroSubmodule_of_exists_boundaryReachabilityData_and_dartSuccessorCycleEmbeddingData_and_selectedBoundaryArc_and_atMostOneInteriorEdgePerFace_heightOrdered_direct}
and its collar-layer/one-collar/canonical-choice variants, show that the
actual boundary-order source used in the route reaches the corrected projected
span conclusion without reintroducing an arbitrary target subspace.
\path{PlanarBoundaryClosedWalkSourceRoute.lean} also carries the same successor-
cycle lowering through the raw Definition~4.8 submodule-duality bridge:
canonical witness choice, the local at-most-one package, and same-boundary
one-collar geometry each have height-ordered and finite-collar wrappers such as
\path{exists_theorem49_kempeGeneratorSubspace_spanning_of_exists_boundaryReachabilityData_and_dartSuccessorCycleEmbeddingData_and_selectedBoundaryArc_and_canonicalWitnessChoice_heightOrdered_direct}.
These theorems deliberately keep the target subspace, raw-generator containment,
and selected-boundary-zero interpretation as hypotheses, marking the unprojected
\path{Theorem49SubmoduleDualityData} obligations rather than hiding them behind the
projected boundary-zero endpoint.
\path{PlanarBoundaryClosedWalkSourceRoute.lean} now records the corresponding
full synthesis endpoints on the same route-facing source shell.  Successor-cycle
boundary data plus selected-arc contiguity and canonical witness choice reaches
\path{Theorem49BoundaryRootSynthesis} directly, and with nonempty
\path{selectedBoundaryInteriorEdgeEndpointVertices} it now also reaches
\path{Theorem49BoundaryRootNonemptySynthesis}.  The local at-most-one package
has the same synthesis endpoint in the facewise cardinal form
\path{exists_theorem49BoundaryRootSynthesis_of_exists_boundaryReachabilityData_and_dartSuccessorCycleEmbeddingData_and_selectedBoundaryArc_and_facewiseAtMostOneInteriorEdge_direct},
and with the same nonempty carrier reaches the non-vacuous package via
\path{exists_theorem49BoundaryRootNonemptySynthesis_of_exists_boundaryReachabilityData_and_dartSuccessorCycleEmbeddingData_and_selectedBoundaryArc_and_facewiseAtMostOneInteriorEdge_with_nonempty_selectedBoundaryInteriorEdgeEndpointVertices_direct}.
A same-boundary one-collar annulus geometry on the successor-cycle source also
has both the direct synthesis wrapper and, with the nonempty carrier hypothesis,
a nonempty-synthesis wrapper.  Thus the current route-facing positive endpoint
is no longer merely a projected span statement: under the same source-side
cardinal or one-collar hypotheses, the full synthesis package is available
without a separate caller-supplied closed-walk source, fallback edge, or
boundary-rest proof.
This positive construction is now witnessed on a genuine honest source model.
In \path{Theorem49PlanarBoundaryAnnulusHonestGeometryRegression.lean},
\path{diamondWithTriangleClosedWalkCanonicalWitnessChoice} constructs
\path{PlanarBoundaryCanonicalWitnessChoice}
for the boundary split extracted from
\path{diamondWithTriangleClosedWalkAnnulusBoundarySource}: the shared diamond
edge witnesses the two diamond faces and one inner triangle edge witnesses the
separate triangle face.  The file packages the source-level existence theorem
\path{exists_embedding_closedWalkAnnulusBoundarySource_and_canonicalWitnessChoice_diamondWithTriangleGraph}
and the downstream one-collar witness endpoint
\path{diamondWithTriangleGraph_admitsPlanarBoundaryAnnulusWitnessAssignment_via_closedWalkCanonicalWitnessChoice}.
So the canonical-witness burden is no longer only a negatively calibrated
hypothesis; it has a checked finite positive source instance.
The same finite instance is now non-vacuous at the coloring layer as well:
\path{diamondWithTriangleTaitEdgeColoring} gives an explicit proper nonzero
edge coloring of \path{diamondWithTriangleGraph} by the three nonzero
\(F_2^2\) colors, and
\path{diamondWithTriangleTaitEdgeColoring_isTait} verifies the Tait condition.
Consequently
\path{diamondWithTriangleEmbedding_theorem49BoundaryRootSynthesis_for_explicitTaitColoring}
is a concrete source-to-synthesis theorem for an actual named Tait coloring,
not merely a theorem parameterized by a possibly empty coloring class.
The diamond benchmark now also reaches that endpoint through the generic
at-most-one criterion itself:
\path{exists_closedWalkSource_with_atMostOneInteriorEdgePerFace_diamondWithTriangleGraph}
packages the local fallback/count/boundary-rest data,
\path{diamondWithTriangleClosedWalkAnnulusBoundarySource_hasCanonicalWitnessChoice_via_atMostOneInteriorEdgePerFace}
reconstructs canonical witness choice from the route theorem, and
\path{diamondWithTriangleEmbedding_theorem49BoundaryRootSynthesis_for_explicitTaitColoring_via_atMostOneInteriorEdgePerFace}
proves the same explicit-Tait synthesis endpoint through this local criterion.
Lean now also exposes the endpoint as a raw theorem-4.9 quotient/decomposition
package:
\path{diamondWithTriangle_closedWalkSource_explicitTait_atMostOne_to_theorem49RawQuotientErrorPackage}
bundles the honest source, explicit coloring, and at-most-one route with the
raw Definition 4.8 representative and selected-boundary-kernel error
decomposition for every Kirchhoff chain on the same embedding.
The status theorem
\path{diamondWithTriangle_explicitTait_atMostOne_synthesis_has_empty_selectedBoundaryInteriorCarrier}
records that this at-most-one path still does not solve the stronger nonempty
target on the diamond benchmark: selected-boundary purification deletes the
interior-edge endpoint carrier.
Lean now packages the whole currently sufficient same-embedding diamond stack
as
\path{diamondWithTriangle_explicitTait_currentSufficientSameEmbeddingGeometry}:
the honest source, canonical witness choice, explicit Tait coloring,
one-collar geometry, collar-layer, height-ordered, well-founded peel data,
attached rooted certificate, and current theorem-4.9 boundary-root synthesis
all coexist on the same embedding.  The same file also proves the stronger
limitation
\path{diamondWithTriangle_explicitTait_synthesis_has_empty_selectedBoundaryInteriorCarrier}.
Thus this benchmark is genuinely positive for the current synthesis surface
but still fails
\path{Theorem49BoundaryRootNonemptySynthesis}, because its purified
selected-boundary interior carrier is empty.  The next positive construction
burden is therefore a Tait-colorable honest source with a nonempty purified
selected-boundary interior carrier, not another proof of the existing diamond
synthesis endpoint.
Lean now lifts this limitation to the current projected endpoint itself:
\path{not_theorem49BoundaryRootNonemptyProjectedSynthesis_diamondWithTriangle}
shows that no coloring on the diamond-with-triangle embedding can inhabit
\path{Theorem49BoundaryRootNonemptyProjectedSynthesis}, because that package
contains the nonempty synthesis carrier obligation.  The status theorem
\path{diamondWithTriangle_explicitTait_synthesis_blocks_nonemptyProjectedSynthesis}
combines the explicit Tait-colored ordinary synthesis endpoint with this
projected-endpoint failure.  Thus the projected Kempe-span equalities do not
repair an empty selected-boundary interior carrier.
The carrier side of that burden is now an exact endpoint-witness condition:
\path{HasUnblockedInteriorEndpoint} names the local condition, and
\path{selectedBoundaryInteriorEdgeEndpointVertices_nonempty_iff_exists_unblocked_interior_endpoint}
and
\path{hasUnblockedInteriorEndpoint_iff_selectedBoundaryInteriorEdgeEndpointVertices_nonempty}
state that the purified carrier is nonempty exactly when some interior
face-incidence edge has an endpoint incident to no selected ambient-boundary
edge.  This gives future positive examples a small local target for the
non-vacuity part, separate from the annulus-collar or witness-cover geometry.
Endpoint-separated annulus constructions can now prove that local target by
supplying a raw interior endpoint carrier and endpoint-support disjointness,
which is the same condition already used by the annulus-root purification
bridge.
That carrier burden is now separated from the geometry burden by the
shared-interior-pair source.  Lean proves
\path{selectedBoundaryInteriorEdgeEndpointVertices_nonempty_sharedInteriorPair}
using vertex \(1\), which is incident to the interior shared edge
\path{sip01} but to no selected-boundary edge.  Combined with
\path{sharedInteriorPairTaitEdgeColoring_isTait}, this gives a concrete
Tait-colorable honest source with nonempty purified carrier; the strengthened
regression now also packages this carrier as the local endpoint witness
\path{hasUnblockedInteriorEndpoint_sharedInteriorPair}.  The universal
obstruction
\path{not_forall_any_replacementPositiveProjectedGeometryOn_of_boundaryReachabilityData_and_dartSuccessorCycleEmbeddingData_and_selectedBoundaryArc_and_taitEdgeColoring_and_hasUnblockedInteriorEndpoint_sharedInteriorPair}
shows that boundary reachability, successor-cycle selected arcs, a Tait
coloring, and explicit \path{HasUnblockedInteriorEndpoint} data still do not
force any direct or route-facing replacement package.  The obstruction has
also been lifted to the honest closed-walk source surface itself:
\path{not_forall_any_replacementPositiveProjectedGeometryOn_of_closedWalkAnnulusBoundarySource_and_taitEdgeColoring_and_hasUnblockedInteriorEndpoint_sharedInteriorPair}
shows that an honest closed-walk annulus boundary source, a Tait coloring, and
the same explicit endpoint witness still do not produce any of the direct or
route-facing replacement packages.  The companion theorem
\path{not_forall_some_witnessCoverData_of_closedWalkAnnulusBoundarySource_and_taitEdgeColoring_and_hasUnblockedInteriorEndpoint_sharedInteriorPair}
pushes the obstruction below the package boundary: the same source, coloring,
and endpoint hypotheses also do not universally produce either raw
height-ordered or collar-layer witness-cover datum.  The bundled theorem
\path{not_forall_some_currentSufficientSameEmbeddingGeometry_of_closedWalkAnnulusBoundarySource_and_taitEdgeColoring_and_hasUnblockedInteriorEndpoint_sharedInteriorPair}
extends the failure to every current same-embedding endpoint known to be
sufficient for synthesis: height-ordered data, collar-layer data, the raw
attached certificate, well-founded witness data, and annulus-collar geometry.
The exact one-collar current-boundary shell already fails at this lower layer
too:
\path{not_forall_some_witnessCoverData_of_closedWalkAnnulusBoundarySource_and_oneCollarAnnulusCurrentBoundaryData_and_taitEdgeColoring_and_hasUnblockedInteriorEndpoint_sharedInteriorPair}
and its successor-cycle analogue show that even after adding exact one-collar
current-boundary data the shared-interior-pair benchmark still does not force
raw height-ordered/collar-layer witness-cover data.
The same exact one-collar/v23 shell now also has generic converse failures at
the direct-or-route-facing replacement-positive layer:
\path{not_forall_any_replacementPositiveProjectedGeometryOn_of_exists_embedding_closedWalkAnnulusBoundarySource_and_oneCollarAnnulusCurrentBoundaryData_and_taitEdgeColoring_and_hasUnblockedInteriorEndpoint_and_v23ResidualBoundaryInitialState}
and
\path{not_forall_any_replacementPositiveProjectedGeometryOn_of_exists_embedding_boundaryReachabilityData_and_dartSuccessorCycleEmbeddingData_and_selectedBoundaryArc_and_oneCollarAnnulusCurrentBoundaryData_and_taitEdgeColoring_and_hasUnblockedInteriorEndpoint_and_v23ResidualBoundaryInitialState}.
So any explicit same-embedding exact-shell witness carrying simultaneous
failure of the current direct and route-facing replacement-positive endpoints
already refutes a universal derivation of that whole positive disjunction; the
shared-interior-pair shell is just one concrete instance.
The same shared-interior-pair benchmark now also blocks the natural residual
same-embedding interface on that exact one-collar/v23 shell:
\path{exists_embedding_closedWalkAnnulusBoundarySource_and_oneCollarAnnulusCurrentBoundaryData_and_taitEdgeColoring_and_hasUnblockedInteriorEndpoint_and_v23ResidualBoundaryInitialState_without_naturalResidualSameEmbeddingGeometry_sharedInteriorPair},
\path{not_forall_some_naturalResidualSameEmbeddingGeometry_of_closedWalkAnnulusBoundarySource_and_oneCollarAnnulusCurrentBoundaryData_and_v23ResidualBoundaryInitialState_and_taitEdgeColoring_and_hasUnblockedInteriorEndpoint_sharedInteriorPair},
and the successor-cycle analogues show that even after adjoining a real exact-v23 residual
initial state, exact one-collar current-boundary data, a Tait coloring, and
\path{HasUnblockedInteriorEndpoint} still do not force residual-boundary witness data,
residual selector data, or even raw height-ordered/collar-layer witness-cover data on the same
embedding.  The same shared-interior-pair shell now also fails the earliest
selector target directly:
\path{exists_embedding_closedWalkAnnulusBoundarySource_and_oneCollarAnnulusCurrentBoundaryData_and_taitEdgeColoring_and_hasUnblockedInteriorEndpoint_and_v23ResidualBoundaryInitialState_without_boundaryFreeIncidentFaceSelector_sharedInteriorPair},
\path{exists_embedding_boundaryReachabilityData_and_dartSuccessorCycleEmbeddingData_and_selectedBoundaryArc_and_oneCollarAnnulusCurrentBoundaryData_and_taitEdgeColoring_and_hasUnblockedInteriorEndpoint_and_v23ResidualBoundaryInitialState_without_boundaryFreeIncidentFaceSelector_sharedInteriorPair},
\path{not_forall_nonempty_boundaryFreeIncidentFaceSelector_of_closedWalkAnnulusBoundarySource_and_oneCollarAnnulusCurrentBoundaryData_and_v23ResidualBoundaryInitialState_and_taitEdgeColoring_and_hasUnblockedInteriorEndpoint_sharedInteriorPair},
and
\path{not_forall_nonempty_boundaryFreeIncidentFaceSelector_of_boundaryReachabilityData_and_dartSuccessorCycleEmbeddingData_and_selectedBoundaryArc_and_oneCollarAnnulusCurrentBoundaryData_and_v23ResidualBoundaryInitialState_and_taitEdgeColoring_and_hasUnblockedInteriorEndpoint_sharedInteriorPair}.
So on the exact one-collar/v23 shell the first same-embedding selector burden
already fails directly, not only as one branch of a broader bundled
obstruction.  The same exact-shell shared-interior-pair benchmark now also
fails the next construction surface directly:
\path{exists_embedding_closedWalkAnnulusBoundarySource_and_oneCollarAnnulusCurrentBoundaryData_and_taitEdgeColoring_and_hasUnblockedInteriorEndpoint_and_v23ResidualBoundaryInitialState_without_planarBoundaryAnnulusConstructionBoundarySupportFaceData_sharedInteriorPair},
\path{exists_embedding_boundaryReachabilityData_and_dartSuccessorCycleEmbeddingData_and_selectedBoundaryArc_and_oneCollarAnnulusCurrentBoundaryData_and_taitEdgeColoring_and_hasUnblockedInteriorEndpoint_and_v23ResidualBoundaryInitialState_without_planarBoundaryAnnulusConstructionBoundarySupportFaceData_sharedInteriorPair},
\path{not_forall_nonempty_planarBoundaryAnnulusConstructionBoundarySupportFaceData_of_closedWalkAnnulusBoundarySource_and_oneCollarAnnulusCurrentBoundaryData_and_v23ResidualBoundaryInitialState_and_taitEdgeColoring_and_hasUnblockedInteriorEndpoint_sharedInteriorPair},
and
\path{not_forall_nonempty_planarBoundaryAnnulusConstructionBoundarySupportFaceData_of_boundaryReachabilityData_and_dartSuccessorCycleEmbeddingData_and_selectedBoundaryArc_and_oneCollarAnnulusCurrentBoundaryData_and_v23ResidualBoundaryInitialState_and_taitEdgeColoring_and_hasUnblockedInteriorEndpoint_sharedInteriorPair}.
So the exact one-collar/v23 shell now fails not only the earliest selector
package but also the selected-boundary-contact construction package consumed by
the later synthesis route.
The earlier
route package
\path{exists_embedding_boundaryReachabilityData_and_dartSuccessorCycleEmbeddingData_and_selectedBoundaryArc_and_taitEdgeColoring_and_nonempty_selectedBoundaryInteriorCarrier_without_currentSufficientSameEmbeddingGeometry_sharedInteriorPair}
shows at the same time that this still does not provide the witness/certificate
geometry needed for the current synthesis route.
A second carrier-calibration benchmark now avoids leaning on the
shared-interior-pair counterexample itself.  The new module
\path{Theorem49PlanarBoundaryAnnulusWheelWitnessRegression.lean} defines
\path{wheelWithInnerTriangleGraph}, with three interior spokes meeting a
central vertex and a separate triangular selected-boundary component.  Lean
proves \path{wheelWithInnerTriangleTaitEdgeColoring_isTait} and
\path{selectedBoundaryInteriorEdgeEndpointVertices_nonempty_wheelWithInnerTriangle},
with the carrier proof now routed through the same unblocked-endpoint
criterion via the central spoke,
then packages them as
\path{wheelWithInnerTriangle_tait_and_nonempty_selectedBoundaryInteriorCarrier}.
The same module now constructs an honest closed-walk annulus boundary source,
\path{wheelWithInnerTriangleClosedWalkAnnulusBoundarySource}, with nonempty
wrapper
\path{nonempty_wheelWithInnerTriangleClosedWalkAnnulusBoundarySource}.  It is
therefore a source-bearing Tait-plus-carrier benchmark, but Lean proves that
it cannot be promoted to the canonical one-collar route:
\path{exists_two_distinct_interior_edges_on_wheelWithInnerTriangle_boundary}
exhibits a face with two distinct interior spokes,
\path{not_nonempty_planarBoundaryCanonicalWitnessChoice_wheelWithInnerTriangle}
rules out canonical witness choice for every boundary split on this embedding,
and
\path{not_exists_oneCollar_planarBoundaryAnnulusCollarGeometry_wheelWithInnerTriangle}
rules out one-collar collar geometry.  The status theorem
\path{wheelWithInnerTriangle_tait_and_nonempty_carrier_without_canonical_or_oneCollar}
keeps the nonempty-carrier/coloring burden distinct from the stronger
same-embedding witness-ownership burden.
The same benchmark has now been calibrated against the face-local
at-most-one sufficient condition.  Lean proves
\path{wheelWithInnerTriangle_face0_interiorEdgeFilter_card}: on face \(0\),
the interior-edge filter has cardinality \(2\).  Hence
\path{not_wheelWithInnerTriangle_atMostOneInteriorEdgePerFace} blocks the
local premise used by
\path{PlanarBoundaryCanonicalWitnessChoice.of_atMostOneInteriorEdgePerFace},
and
\path{not_exists_closedWalkSource_with_atMostOneInteriorEdgePerFace_wheelWithInnerTriangleEmbedding}
shows that no honest closed-walk source on this fixed embedding can satisfy
the full fallback/at-most-one/boundary-rest package.  The sharper bundle
\path{wheelWithInnerTriangle_tait_and_nonempty_carrier_without_atMostOne_or_canonical_or_oneCollar}
therefore records that Tait colorability and nonempty purified carrier still
do not imply the at-most-one local geometry, canonical witness choice, or
one-collar geometry.
The same fixed wheel embedding now blocks the raw replacement witness-cover
surface as well.  Lean proves
\path{not_nonempty_planarBoundaryHeightOrderedFacePeelWitnessData_wheelWithInnerTriangle}:
covering one central spoke forces the remaining spoke on that triangular face
to be witnessed at a strictly later height, and following the three wheel faces
forces a strict cycle through \path{wit01}, \path{wit02}, and \path{wit03}.
The same dependency cycle now blocks the acyclic parent-peeling surface:
\path{Theorem49PlanarBoundaryPeeling.lean} proves
\path{PlanarBoundaryWellFoundedFacePeelWitnessData.toPlanarBoundaryHeightOrderedFacePeelWitnessData}
using the canonical parent-height rank, and the wheel regression applies this
as
\path{not_nonempty_planarBoundaryWellFoundedFacePeelWitnessData_wheelWithInnerTriangle}.
This blocker is now factored out of the wheel proof itself.  The same peeling
file proves that, under injective witness choices on peeled faces, a
non-boundary remainder edge equal to another peeled face's witness edge forces a
parent dependency in the well-founded package and a strict inequality in the
height-ordered package.  The list-level order lemma
\path{false_of_isChain_height_lt_and_getLast_lt_head} and the two
\path{false_of_witnessRemainderDependency_cycle_of_injective_witness} lemmas
therefore turn any nonempty closed non-boundary remainder-witness dependency
chain into a formal blocker, independent of the particular wheel names.
The wheel regression now uses the same list-level obstruction directly in its
height-ordered no-data proof, applying it to the two possible three-face
cycles around the spokes instead of closing those branches by a bespoke
transitive-irreflexive calculation on natural-number heights.
The well-founded parent side now also has the parent-map analogue
\path{false_of_isChain_parentRel_and_getLast_parent_head}: a nonempty finite
closed chain in \path{ParentRel} contradicts the well-founded parent map
directly.  Accordingly, the well-founded witness-cover finite-cycle blocker
now first extracts actual parent edges and applies this parent-chain theorem,
instead of rebuilding the obstruction locally as a parent-height cycle.
The finite collar-layer version is ruled out by
\path{not_nonempty_planarBoundaryCollarLayerFacePeelWitnessData_wheelWithInnerTriangle},
since collar-layer data lowers to height-ordered witness-cover data.  The
obstruction also reaches weak annulus-collar geometry itself:
\path{not_nonempty_planarBoundaryAnnulusCollarGeometry_wheelWithInnerTriangle}
uses the canonical collar-to-height lowering, and
\path{not_forall_some_acyclicWitnessCover_or_annulusCollarGeometry_of_taitEdgeColoring_and_nonempty_selectedBoundaryInteriorCarrier_wheelWithInnerTriangle}
packages the failed universal from Tait colorability and nonempty purified
carrier to any current acyclic witness-cover or weak collar endpoint.  Thus the
missing positive construction is not merely a nonempty carrier or Tait coloring
condition; it must also break or orient the interior-face dependency cycle before
even the well-founded/height/collar witness surfaces, and hence weak
annulus-collar geometry, can be expected.
The source-bearing form is now checked as well:
\path{wheelWithInnerTriangle_closedWalkSource_tait_and_nonempty_carrier_without_acyclicWitnessCover_or_annulusCollarGeometry}
bundles the honest closed-walk source, Tait coloring, nonempty purified
carrier, and failure of all current acyclic witness-cover and weak
annulus-collar endpoints.  Its universal obstruction
\path{not_forall_some_acyclicWitnessCover_or_annulusCollarGeometry_of_closedWalkAnnulusBoundarySource_and_taitEdgeColoring_and_nonempty_selectedBoundaryInteriorCarrier_wheelWithInnerTriangle}
shows that adding the honest closed-walk annulus source to the wheel
hypotheses still does not construct the positive geometry; the missing input
has to break or orient the dependency cycle.
The same wheel model now also has an explicit forcing-edge calibration.
\path{Theorem49ForcingInteriorEdgeObstructionRegression.lean} proves
\path{wheelWithInnerTriangleForcingInteriorEdgeWitness}: the central spoke
\path{wit01} is an interior edge, and every face incident to it already
contains selected-boundary support, namely \path{wit12} on one wheel face and
\path{wit31} on the other.  The generic forcing-edge obstruction then yields
\path{not_nonempty_planarBoundaryAnnulusRootAdjDistancePeelData_wheelWithInnerTriangle_via_forcingInteriorEdgeWitness}
and
\path{not_nonempty_interiorDualBoundaryRootAdjDistancePeelData_wheelWithInnerTriangle_via_forcingInteriorEdgeWitness},
as well as
\path{not_nonempty_planarBoundaryAnnulusConstructionFaceLayerData_wheelWithInnerTriangle_via_forcingInteriorEdgeWitness}.
The bundled theorem
\path{wheelWithInnerTriangle_tait_nonempty_carrier_and_forcing_obstruction}
records the sharper status: the benchmark has an explicit Tait coloring and
nonempty purified carrier, but it carries the same obstruction that blocks the
current annulus-root and face-layer construction routes on the fixed
embedding.
The same source-side obstruction is now visible without choosing a benchmark.
In \path{Theorem49PlanarBoundaryPeeling.lean},
\path{PlanarBoundaryHeightOrderedFacePeelWitnessData.exists_terminal_face_selectedBoundary_remainders}
extracts a terminal peeled face from any nonempty height-ordered witness-cover
package: all non-witness remainders on that face already lie on selected
boundary support.  For
\path{InteriorDualBoundaryRootAdjDistancePeelData}, peeled faces are disjoint
from selected-boundary support, so
\path{InteriorDualBoundaryRootAdjDistancePeelData.exists_terminal_peelFace_witness_erase_eq_empty_of_mem_interiorEdgeSupport}
forces that terminal face to have no non-witness boundary edges.  The corollary
\path{InteriorDualBoundaryRootAdjDistancePeelData.exists_peelFace_card_le_one_of_mem_interiorEdgeSupport}
therefore says that any IDBRAD package covering even one interior edge contains
a peeled face with boundary cardinality at most one.  This is a direct
source-side necessary condition: a proposed boundary-order derivation of IDBRAD
must explain the terminal singleton-edge face rather than only supplying
coverage and root distances.
The wheel-with-inner-triangle regression now instantiates this criterion on a
concrete source.  Lean proves
\path{wheelWithInnerTriangleFaceBoundary_card_eq_three} and
\path{wheelWithInnerTriangle_two_le_faceBoundary_card}, so every ambient face in
that benchmark has at least two boundary edges.  With
\path{wit01_mem_interiorEdgeSupport}, the theorem
\path{not_nonempty_interiorDualBoundaryRootAdjDistancePeelData_wheelWithInnerTriangle_via_terminal_card_lower_bound}
refutes IDBRAD without appealing to the forcing-edge selector interface.  The
source-facing package
\path{wheelWithInnerTriangle_closedWalkSource_tait_hasUnblockedEndpoint_without_interiorDualBoundaryRootAdjDistancePeelData_via_terminal_card}
keeps the honest closed-walk source, Tait coloring, and local unblocked
endpoint on the same embedding while deriving the contradiction through the
terminal-card obstruction.
The same concrete source is now connected to the actual projected Theorem 4.9
surface.  Lean adds
\path{wheelWithInnerTriangleGraph_edgeSet_fintype} and
\path{hasUnblockedInteriorEndpoint_wheelWithInnerTriangle}, then proves
\path{wheelWithInnerTriangle_theorem49BoundaryRootNonemptyProjectedSynthesis_of_interiorDualBoundaryRootAdjDistancePeelData}:
if the wheel embedding had a same-embedding IDBRAD datum, the existing source
bridge would yield
\path{Theorem49BoundaryRootNonemptyProjectedSynthesis} for the concrete wheel
Tait coloring.  The raw Definition 4.8 endpoint is also instantiated as
\path{wheelWithInnerTriangle_theorem49BoundaryRawQuotientErrorPackage_of_interiorDualBoundaryRootAdjDistancePeelData},
and
\path{wheelWithInnerTriangle_source_to_theorem49_and_terminal_card_contradiction_of_interiorDualBoundaryRootAdjDistancePeelData}
records that this same hypothetical datum is impossible by the terminal-card
criterion.
The regression now states the same obstruction directly at the blocked
projected endpoint:
\path{not_exists_wheelWithInnerTriangle_closedWalkSource_and_interiorDualBoundaryRootAdjDistancePeelData_and_theorem49BoundaryRootNonemptyProjectedSynthesis}
rules out coexistence of the closed-walk source, a same-embedding IDBRAD datum,
and the projected Theorem 4.9 endpoint on the wheel embedding.  The bundled
source-facing theorem
\path{wheelWithInnerTriangle_closedWalkSource_tait_hasUnblockedEndpoint_without_interiorDualBoundaryRootAdjDistancePeelData_based_projectedSynthesis}
keeps the honest source, concrete Tait coloring, and unblocked endpoint while
ruling out an IDBRAD-based projected-endpoint instantiation.
The same wheel regression now adds the parent-route counterpart
\path{wheelWithInnerTriangle_closedWalkSource_tait_hasUnblockedEndpoint_without_boundaryFaceRootsCanonicalParent_based_projectedSynthesis}:
the positive source, Tait, and endpoint facts remain present, but no projected
endpoint can be obtained from source boundary-face roots plus canonical parent
child-membership on this embedding.
The cyclic-source derivation file has now been repaired into the corresponding
checked failed universal.  In
\path{Theorem49PlanarBoundaryAnnulusHonestWitnessRegression.lean}, Lean proves
\path{sharedInteriorPair_sipFace0_interiorEdgeFilter_card}: on the honest
\path{sharedInteriorPairAnnulusEmbedding}, the face \path{sipFace0} filters to
the two interior edges \path{sip01} and \path{sip12}.  Hence
\path{not_sharedInteriorPair_atMostOneInteriorEdgePerFace}.  The route-facing
module \path{Theorem49CyclicSourceAtMostOneDerivation.lean} now proves
\path{route_impossibility_at_most_one_not_derivable_from_cyclic_source} and
\path{not_forall_closedWalkSource_implies_atMostOneInteriorEdgePerFace}: an
honest closed-walk annulus boundary source can exist while the face-local
at-most-one premise fails.  Thus the at-most-one condition remains an
independent sufficient hypothesis, not a theorem of cyclic or closed-walk
boundary-order source data.
The strengthened-condition derivation file now makes the same point on the
route-facing successor-cycle surface.  In
\path{Theorem49StrengthenedConditionAtMostOneDerivation.lean}, the theorem
\path{exists_boundaryReachabilityData_and_dartSuccessorCycleEmbeddingData_and_selectedBoundaryArc_without_atMostOneInteriorEdgePerFace_sharedInteriorPair}
packages the shared-interior-pair embedding with annulus boundary reachability,
dart-successor facial geometry, and selected-boundary arcs, while
\path{not_forall_boundaryReachabilityData_and_dartSuccessorCycleEmbeddingData_and_selectedBoundaryArc_and_boundaryComponentConstant_implies_atMostOneInteriorEdgePerFace_sharedInteriorPair}
adds the boundary-component constancy normally obtained from the honest source
fields.  The same `sipFace0` two-interior-edge obstruction still refutes the
at-most-one conclusion.
The route summary file now packages these separate calibrations as the checked
theorem
\path{theorem49_route_obstruction_matrix} in
\path{Theorem49DerivationImpossibilitySummary.lean}.  Its four components
record the cyclic-source counterexample to at-most-one, the chained-diamonds
counterexample to deriving a nonempty purified carrier from source plus
at-most-one, the strengthened successor-cycle failed universal with
boundary-component constancy, and the wheel-with-inner-triangle
Tait/nonempty-carrier model carrying a forcing interior edge that blocks the
current annulus-root and face-layer construction routes.
This has now been extended to a second, strictly larger positive model rather
than another constructor theorem.  The same regression file defines
\path{chainedDiamondsWithTriangleGraph}, whose outer boundary component
contains two chained diamond cells and whose inner boundary component is a
separate triangle.  The explicit source
\path{chainedDiamondsClosedWalkAnnulusBoundarySource} supplies honest facial
closed walks and selected-boundary arcs on all five faces.  The direct finite
witness
\path{chainedDiamondsClosedWalkCanonicalWitnessChoice} chooses the two
distinct interior shared edges \path{cdt12} and \path{cdt45} on their
respective diamond face pairs and a boundary edge on the inner triangle.  Lean
now proves \path{chainedDiamondsWithTriangle_interiorEdgeSupport_eq}: the
interior edge support is exactly \(\{\mathtt{cdt12}, \mathtt{cdt45}\}\), and
\path{chainedDiamondsWithTriangle_atMostOneInteriorEdgePerFace}: each face has
at most one interior edge on its boundary, so this model satisfies the
facewise local condition enabling the canonical-choice constructor.  The
fallback and boundary-rest obligations are now explicit as
\path{chainedDiamondsCanonicalWitnessEdge_mem} and
\path{chainedDiamonds_nonInteriorOnAmbient}, and the full generic premise is
packaged as
\path{exists_closedWalkSource_with_atMostOneInteriorEdgePerFace_chainedDiamondsWithTriangleGraph}.
Lean reconstructs the canonical witness choice through the generic constructor
as
\path{chainedDiamondsClosedWalkAnnulusBoundarySource_hasCanonicalWitnessChoice_via_atMostOneInteriorEdgePerFace}
and lowers it to
\path{chainedDiamondsWithTriangleGraph_admitsPlanarBoundaryAnnulusWitnessAssignment_via_atMostOneInteriorEdgePerFace}
and
\path{chainedDiamondsWithTriangleGraph_admitsPlanarBoundaryAnnulusCollarGeometry_via_atMostOneInteriorEdgePerFace}.
The endpoint bridge also applies to this larger benchmark:
\path{chainedDiamondsWithTriangleEmbedding_theorem49BoundaryRootSynthesis_via_atMostOneInteriorEdgePerFace}
and
\path{exists_theorem49BoundaryRootSynthesis_chainedDiamondsWithTriangleGraph_via_atMostOneInteriorEdgePerFace}
show that, conditional on a Tait coloring, the local criterion reaches the
same theorem-4.9 synthesis surface as the named one-collar geometry.
The same benchmark is now calibrated against the stronger nonempty-carrier
surface:
\path{selectedBoundaryInteriorEdgeEndpointVertices_eq_empty_chainedDiamondsWithTriangle}
and
\path{not_selectedBoundaryInteriorEdgeEndpointVertices_nonempty_chainedDiamondsWithTriangle}
show that the purified carrier is empty, because every endpoint of
\path{cdt12} and \path{cdt45} is incident to a selected-boundary edge.
Together with
\path{not_exists_taitEdgeColoring_chainedDiamondsWithTriangleGraph}, the
status theorem
\path{chainedDiamonds_atMostOne_source_has_empty_carrier_and_no_tait}
records that this model supplies at-most-one source geometry but neither a
surviving carrier nor a Tait coloring.
The route-impossibility module now makes this obstruction a checked failed
universal, not just a benchmark note:
\path{route_impossibility_at_most_one_does_not_imply_nonempty_carrier} gives
the concrete \path{Fin 10} counterexample,
\path{not_forall_atMostOneInteriorEdgePerFace_implies_nonempty_carrier}
rules out deriving nonempty purified carrier from at-most-one alone, and
\path{not_forall_closedWalkSource_and_atMostOneInteriorEdgePerFace_implies_nonempty_carrier}
shows that adding the honest closed-walk annulus boundary source still does not
recover nonemptiness.
The stronger route-facing successor-cycle source shell now gives a sharper
diagnosis.  \path{Theorem49AtMostOneNonemptyCarrierImpossibility.lean}
constructs explicit facial successor maps on the diamond-with-triangle model in
\path{diamondWithTriangleDartSuccessorCycleGeometry}, proves selected-arc
contiguity by
\path{diamondWithTriangleDartSuccessorCycleGeometry_selectedBoundaryArcOnFace},
and packages the counterexample as
\path{successor_cycle_at_most_one_and_empty_carrier_independent}.  The failed
universal
\path{not_forall_boundaryReachabilityData_and_dartSuccessorCycleEmbeddingData_and_selectedBoundaryArc_and_atMostOneInteriorEdgePerFace_implies_nonempty_carrier}
is now explained by a universal obstruction:
\path{selectedBoundaryInteriorEdgeEndpointVertices_eq_empty_of_closedWalkEmbeddingData_and_facewiseAtMostOneInteriorEdge}
proves that honest facial closed-walk data plus facewise at-most-one forces
the purified interior-edge endpoint carrier to be empty on any embedding.  The
closed-walk contradiction forms
\path{not_exists_closedWalkEmbeddingData_and_facewiseAtMostOneInteriorEdge_and_nonempty_selectedBoundaryInteriorEdgeEndpointVertices}
and
\path{not_exists_closedWalkAnnulusBoundarySource_and_facewiseAtMostOneInteriorEdge_and_nonempty_selectedBoundaryInteriorEdgeEndpointVertices}
rule out the current honest-source at-most-one non-vacuous endpoint outright.
The same obstruction is now available in constructive contrapositive form:
\path{exists_two_distinct_interior_edges_on_faceBoundary_of_closedWalkEmbeddingData_and_nonempty_selectedBoundaryInteriorEdgeEndpointVertices}
and the route-facing
\path{exists_two_distinct_interior_edges_on_faceBoundary_of_closedWalkAnnulusBoundarySource_and_hasUnblockedInteriorEndpoint}
show that any honest-source endpoint witness for the current carrier must
exhibit a face boundary with two distinct interior edges.
The positive annulus-collar replacement surface now exposes this obstruction
directly:
\path{exists_two_distinct_interior_edges_on_faceBoundary_of_closedWalkAnnulusCollarPositiveProjectedGeometryOn}
consumes
\path{Theorem49ClosedWalkAnnulusCollarPositiveProjectedGeometryOn}, so that
route package is explicitly outside the facewise at-most-one branch before any
theorem-4.9 synthesis endpoint is invoked.
The same route file now also proves
\path{exists_two_distinct_interior_edges_on_faceBoundary_of_successorCycleAnnulusCollarPositiveProjectedGeometryOn},
so the successor-cycle annulus-collar shell carries that same two-interior-edge
obstruction directly rather than only through silent lowering to closed-walk
data.  Correspondingly,
\path{not_theorem49ClosedWalkAnnulusCollarPositiveProjectedGeometryOn_of_facewiseAtMostOneInteriorEdge}
and
\path{not_theorem49SuccessorCycleAnnulusCollarPositiveProjectedGeometryOn_of_facewiseAtMostOneInteriorEdge}
show that neither route-facing annulus-collar positive package can inhabit the
canonical facewise-at-most-one branch.  The stronger selected-boundary-contact
construction shell is now blocked there as well:
\path{not_theorem49ClosedWalkAnnulusConstructionBoundarySupportPositiveProjectedGeometryOn_of_facewiseAtMostOneInteriorEdge}
and
\path{not_theorem49SuccessorCycleAnnulusConstructionBoundarySupportPositiveProjectedGeometryOn_of_facewiseAtMostOneInteriorEdge}
show that even that stronger source-facing package cannot inhabit the same
branch, while
\path{exists_two_distinct_interior_edges_on_faceBoundary_of_successorCycleAnnulusConstructionBoundarySupportPositiveProjectedGeometryOn}
exposes the same local two-interior-edge obstruction directly on the
successor-cycle selected-boundary-contact shell.  The older successor-cycle theorem
\path{selectedBoundaryInteriorEdgeEndpointVertices_eq_empty_of_dartSuccessorCycleEmbeddingData_and_facewiseAtMostOneInteriorEdge}
and its contradiction form
\path{not_exists_dartSuccessorCycleEmbeddingData_and_facewiseAtMostOneInteriorEdge_and_nonempty_selectedBoundaryInteriorEdgeEndpointVertices}
shows that the current successor-cycle at-most-one branch cannot be made
non-vacuous by merely adding the existing carrier nonemptiness hypothesis.
The older direct package remains available as
\path{exists_embedding_closedWalkAnnulusBoundarySource_and_canonicalWitnessChoice_chainedDiamondsWithTriangleGraph}
and lowers it to
\path{chainedDiamondsWithTriangleGraph_admitsPlanarBoundaryAnnulusWitnessAssignment_via_closedWalkCanonicalWitnessChoice}.
Thus the positive side is no longer represented by a single one-interior-edge
toy instance: canonical witness choice is inhabited on an honest source with
two separate interior witness edges and a connected outer boundary component,
and the generic at-most-one theorem itself is non-vacuous on that model.
The same model now reaches the one-collar geometry surface as a named concrete
object:
\path{chainedDiamondsOneCollarGeometry}.  Lean proves
\path{chainedDiamondsOneCollarGeometry_numCollars} and
\path{chainedDiamondsOneCollarGeometry_boundaryData}, so this geometry has
exactly one collar and preserves the boundary split extracted from
\path{chainedDiamondsClosedWalkAnnulusBoundarySource}.  The source-level
existence theorem
\path{exists_embedding_closedWalkAnnulusBoundarySource_and_oneCollarAnnulusCollarGeometry_chainedDiamondsWithTriangleGraph}
and graph-level admissibility theorem
\path{chainedDiamondsWithTriangleGraph_admitsPlanarBoundaryAnnulusCollarGeometry_via_closedWalkCanonicalWitnessChoice}
show that the larger concrete model inhabits the same one-collar route surface
that the theorem-4.9 synthesis layer consumes.
The same model now also reaches the downstream certificate surfaces.  The
vacuous positive-collar placement theorem
\path{chainedDiamondsOneCollarWitnessOnCurrentBoundary} upgrades the one-collar
geometry to
\path{chainedDiamondsOneCollarPreviousBoundaryWitnessGeometry}; Lean then
constructs
\path{chainedDiamondsOneCollarCollarLayerFacePeelWitnessData},
\path{chainedDiamondsOneCollarHeightOrderedFacePeelWitnessData},
\path{chainedDiamondsOneCollarWellFoundedFacePeelWitnessData}, and the attached
certificate
\path{chainedDiamondsOneCollarAttachedRootedFacePeelCertificate}.  The graph
admissibility endpoints through these surfaces, together with
\path{closedWalkAnnulusBoundarySource_and_annulusCollarGeometry_and_attachedRootedFacePeelCertificate_chainedDiamondsWithTriangle},
show that the larger honest source is not merely an annulus-geometry example:
it already inhabits the certificate branch consumed by the current synthesis
layer.  Finally,
\path{chainedDiamondsWithTriangleEmbedding_theorem49BoundaryRootSynthesis_via_oneCollarGeometry}
and
\path{exists_theorem49BoundaryRootSynthesis_chainedDiamondsWithTriangleGraph_via_oneCollarGeometry}
package the concrete theorem-4.9 boundary-root synthesis endpoint for any Tait
coloring of this finite graph.  The coloring quantifier is now marked honest:
\path{not_exists_taitEdgeColoring_chainedDiamondsWithTriangleGraph} proves
that this larger graph has no Tait edge coloring, since four distinct edges
incident with vertex \(3\) would need four pairwise distinct nonzero colors
while only three are available.  Thus the chained model remains valuable as a
larger positive honest-source, one-collar, and certificate construction, but
its theorem-4.9 synthesis endpoint is formally vacuous as a Tait instance.

This criterion is now calibrated negatively as well.  The generic theorem
\path{not_nonempty_planarBoundaryCanonicalWitnessChoice_of_exists_two_distinct_interior_edges_on_faceBoundary}
in \path{Theorem49PlanarBoundaryAnnulusWitness.lean} shows that if one ambient
face carries two distinct interior edges, the canonical facewise witness choice
itself cannot exist, since such a choice would upgrade to the already
forbidden canonical witness assignment.  The abstract boundary-data
countermodel is recorded by
\path{not_forall_planarBoundaryAnnulusBoundaryData_implies_nonempty_canonicalWitnessChoice}
in \path{Theorem49PlanarBoundaryAnnulusWitnessRegression.lean}.  The honest
source countermodel is recorded by
\path{not_forall_exists_planarBoundaryCanonicalWitnessChoice_of_closedWalkAnnulusBoundarySource_sharedInteriorPair}
in \path{Theorem49PlanarBoundaryAnnulusHonestWitnessRegression.lean}.  Thus
the explicit canonical-witness-choice repair is a genuine extra geometric
hypothesis, not a consequence of annulus boundary data or of the honest
closed-walk annulus source.

\section{Threshold assessment}

The current obstruction suite should not be overstated.  It gives high
confidence, checked in Lean, against three specific shortcuts:
\begin{itemize}[leftmargin=2em]
\item weak collar geometry alone does not imply parent-witness orientation;
\item the honest closed-walk two-component boundary source is incompatible
with the current all-roots-on-the-outer-boundary adjacency-distance peel data
when both use the same boundary split;
\item edge-disjoint peeled annulus data does not imply endpoint no-touch for
the raw interior-edge endpoint carrier;
\item IDBRAD data covering an interior edge forces a terminal peeled face of
boundary cardinality at most one, so source geometries whose peeled faces are
all multi-edge faces cannot instantiate the current package directly;
\item selected-boundary purification can be vacuous in endpoint-touch models,
so non-vacuous use of the purified Theorem~4.9 endpoint needs an explicit
nonemptiness or no-touch hypothesis.
\end{itemize}

Those are serious route blockers for the manuscript as currently written.
However, they do not yet meet the \(>90\%\) threshold for declaring the whole
Goertzel proof route robustly broken.  The reason is also formal: the Lean tree
still contains positive replacement surfaces, but their scope is now narrower.
The root-distance annulus construction proves parent-witness orientation from
\path{InteriorDualBoundaryRootAdjDistancePeelData}, and the current synthesis
module now factors the purified Theorem~4.9 endpoint through generic
\path{PlanarBoundaryInteriorDualBoundaryRootAdjDistancePeelData}, re-exposing
it on both \path{PlanarBoundaryOuterBoundaryRootAdjDistancePeelData} and the
two-sided annulus-root surface.  The new closed-walk obstruction shows that
the old outer-root package cannot be instantiated from the honest closed-walk
two-component boundary split while retaining the all-outer-root condition.  The new two-sided
\path{PlanarBoundaryAnnulusRootAdjDistancePeelData} surface is the current
ground-up replacement for that exact failure mode: it retains the existing
interior-dual proof pipeline but removes the all-outer-root localization that
the closed-walk source refutes.  These surfaces therefore remain useful
theorem interfaces, but a continuation must now source the two-sided root data
from closed-walk or cyclic boundary-order geometry rather than treating the
old outer-root package as the direct annulus endpoint.

So the present verdict is deliberately conservative.  The accumulated crux
work is enough to say that the weak v23 annulus interfaces cannot be accepted
as proof; it is not enough to say that every strengthened annulus/no-touch or
selected-boundary purification strategy is blocked.  Further crux polishing at
the same endpoint is now lower value than ground-up work on the missing
geometric source: closed-walk or cyclic boundary-order data, parent-in-peeled
annulus geometry, and non-vacuous selected-boundary interior carriers.

\section{Status summary}

\begin{tabular}{@{}>{\raggedright\arraybackslash}p{0.36\textwidth}>{\raggedright\arraybackslash}p{0.16\textwidth}>{\raggedright\arraybackslash}p{0.38\textwidth}@{}}
\toprule
Claim surface & Status & Evidence in hand \\
\midrule
Older per-run purification claim \(\gamma \cdot 1_R\) & Refuted & \path{GoertzelLemma43Obstruction.lean} and the archived old branch record \\
Exact v23 Lemma~4.3 run-invariance statement & Not refuted here & manuscript audit only; no direct Lean counterexample presently attached \\
Exact v23 Step~2 claim \(c \cdot 1_{D \cap \partial f}\) & Proved & \path{GoertzelLemma44.lean} \\
Manuscript spanning-forest leaf reduction & Refuted as stated; repaired in live Lean route & \path{Theorem49ForestRegression.lean}, \path{Theorem49InteriorDualForest.lean}, \path{PlanarBoundaryAnnulusConstruction.lean} \\
Manuscript adjacency-only well-founded peeling reduction & Refuted as stated; repaired in live Lean route & \path{Theorem49ForestRegression.lean}, \path{Theorem49InteriorDualForest.lean} \\
Live boundary-root adjacency-distance / annulus-construction route & Formalized theorem surface & \path{Theorem49InteriorDualForest.lean}, \path{Theorem49PlanarBoundaryOuterBoundaryRootInteriorDual.lean}, \path{PlanarBoundaryAnnulusConstruction.lean} \\
Weak collar geometry implies parent-witness orientation & Refuted by lowered countermodel & \path{Theorem49PlanarBoundaryAnnulusWitnessRegression.lean}, \path{PlanarBoundaryAnnulusConstruction.lean} \\
Honest source plus arbitrary weak collar implies repaired previous-boundary witness placement & Refuted on a connected source embedding & \path{Theorem49HonestSourceForcesPreviousBoundaryWitness.lean} \\
Closed-walk boundary source plus all-outer-root peel data on the same split & Refuted by structural contradiction & \path{Theorem49ClosedWalkOuterRootObstruction.lean} \\
Global \(>90\%\) proof-route crux block & Not reached; shift to ground-up geometry/synthesis obligations & obstruction files plus positive replacement surfaces in \path{PlanarBoundaryAnnulusConstruction.lean} and \path{Theorem49Synthesis.lean} \\
Current embedding-side boundary-order dependency & Still the main live blocker & \path{fourcolor-boundary-order-notes.tex} and regression files \\
Interior-vertex side condition for boundary erasure & Purification bridge proved; annulus-construction edge-disjoint peeled faces do not imply the exact peeled-face endpoint no-touch condition; raw-carrier endpoint no-touch is necessary or purification must be kept & \path{Theorem49InteriorVertices.lean}, \path{Theorem49InteriorVerticesEndpointObstruction.lean} \\
\bottomrule
\end{tabular}

\section{What this means for continuation}

The correct present verdict is:
\begin{itemize}[leftmargin=2em]
\item the current route is still blocked in the later Theorem~4.9 /
boundary-order part;
\item the obstruction is \emph{not} just lack of labor;
\item the older per-run contradiction remains important historical evidence;
\item the exact v23 Step~2 purification claim is now formally supported;
\item the manuscript's naive leaf-peeling reductions are false as stated, but
the live Lean route already replaces them with stronger adjacency-distance,
root-distance, height, and well-founded parent interfaces;
\item honest closed-walk source data plus arbitrary same-embedding weak collar
geometry does not by itself repair the previous-boundary witness placement; the
collar must be tied to the source split, strengthened, or accompanied by the
explicit current-boundary witness condition;
\item the all-outer-root version of the adjacency-distance route cannot be
instantiated from the honest closed-walk two-component boundary split;
\item any sound v23.5 theorem targeting the downstream construction face-layer
shell must already discharge the local boundary-free-selector / no-forcing
obligation there; that burden is not postponed to later decomposition or
projection wrappers;
\item the remaining blocker is to derive those stronger interfaces from honest
annulus / boundary-order geometry.
\item the current evidence does not justify declaring the full route \(>90\%\)
blocked; it justifies shifting effort from more endpoint counterexamples to
ground-up formalization of the stronger annulus geometry.
\end{itemize}

This also means that v23 should not be described as the first version to have
faced a machine-level challenge.  Rather, it is the version that followed an
earlier contradiction and was supposed to repair that weakness.  The present
formalization now shows that this specific local repair succeeds, while the
later route is still incomplete.

This does \emph{not} by itself prove that no repair of the overall Goertzel
route can exist.  But it does show that any viable continuation must revise the
current theorem surfaces and interfaces.  In particular, a repair would have to
do some combination of the following:
\begin{enumerate}[leftmargin=2em]
\item continue from the now-checked local purification claim into the later
Theorem~4.9 route without reintroducing stronger unsupported reductions;
\item derive the already-formalized boundary-root adjacency-distance package
from the intended annulus geometry, rather than falling back to the refuted
leaf heuristic;
\item derive the already-formalized root-distance / height / well-founded
peeling data from that geometry, rather than using adjacency-only parent
arguments;
\item derive ordered boundary-walk information from stronger embedding-side data,
or else postulate and track that stronger data explicitly;
\item identify the intended annulus interior vertex set with the checked
selected-boundary purification used to discharge the boundary-erasure side
condition.
\end{enumerate}

\section{Relevant files}

\begin{tabular}{@{}>{\raggedright\arraybackslash}p{0.52\textwidth}>{\raggedright\arraybackslash}p{0.38\textwidth}@{}}
\toprule
File & Role \\
\midrule
\path{.../GoertzelLemma43Obstruction.lean} & Older per-run purification obstruction \\
\path{.../GoertzelLemma44.lean} & Exact v23 Step 2 purification theorem \\
\path{.../Theorem49ForestRegression.lean} & Current later-route forest / peeling regressions \\
\path{.../Theorem49*.lean} & Theorem~4.9 theorem surfaces \\
\path{.../Theorem49InteriorVertices.lean} & Selected-boundary interior-vertex purification and discharge \\
\path{.../fourcolor-boundary-order-notes.tex} & Boundary-order gap note \\
\path{public/literature/math/fourcolor/README.md} & Canonical Coq/Rocq 4CT proof line \\
\path{tmp/ai-agents-merge-check/lean-projects/fourcolor/README.md} & Earlier Lean route branch and older Lemma~4.3 contradiction record \\
\path{tmp/ai-agents-merge-check/isabelle/4CP/README.md} & Archival Isabelle route branch \\
\bottomrule
\end{tabular}

\section{Update (2026-06-10): uniform-oracle refutations and the restructured frontier}

The Lean tree was restructured to expose the audit's conclusion as a small
formal surface.  Three new files sit on top of the development:
\path{Goal.lean} (the target and the oracle reductions/refutations below),
\path{Shells.lean} (bundled hypothesis shells \path{ClosedWalkExactShell} and
\path{SuccessorCycleExactShell} replacing the historical 8--10-hypothesis
telescopes), and \path{Frontier.lean} (the maximal positive and negative
results restated over the bundles).  The historical obstruction ledger (19
leaf files) moved to \path{Mettapedia/GraphTheory/FourColor/Legacy/} with
declaration names unchanged; \path{ROADMAP.md} is the new entry point.

The substantive sharpening over the April threshold assessment is this.
\path{Goal.lean} defines the manuscript-shaped target
\path{Theorem49ShellClaim} --- full \path{Theorem49BoundaryRootSynthesis} on
every embedding carrying the endpoint-bearing exact one-collar/v23 shell ---
and proves it would follow from any of four uniform geometric oracles:
repaired previous-boundary witness geometry
(\path{PreviousBoundaryGeometryOracle}), explicit well-founded peel-witness
data (\path{WellFoundedPeelOracle}), generic interior-dual root-distance peel
data (\path{InteriorDualPeelOracle}), or the v23.5 residual/current-boundary
positive wrapper (\path{ResidualBoundaryGeometryOracle}).  All four uniform
oracles are now \emph{proved false}
(\path{not_previousBoundaryGeometryOracle},
\path{not_wellFoundedPeelOracle}, \path{not_interiorDualPeelOracle},
\path{not_residualBoundaryGeometryOracle}): the wheel-with-inner-triangle and
shared-interior-pair benchmarks inhabit the full shells while refuting each
geometric layer.

The fourth refutation is the formal verdict on the v23.5 revision memo's
repair recommendation.  The memo's residual/current-boundary lane, as
currently specified, is conservative: Lean proves its positive wrapper
fixed-embedding equivalent to the collar-layer surface
(\path{theorem49ResidualBoundaryPositiveProjectedGeometryOn_iff_collarLayerPositiveProjectedGeometryOn}),
so it inherits the shared-interior-pair refutation verbatim.  The memo's
``algebraic cancellation certificate'' obligation therefore cannot be
discharged by any interface equivalent to the existing geometry ladders; a
surviving v23.5 route must introduce a genuinely new interface, and that
interface must provably fail on the benchmarks, since they inhabit the shells.

What remains open is a clean dichotomy, recorded in \path{ROADMAP.md}: either
construct one of the oracle inputs from natural hypotheses the benchmarks do
not satisfy (candidates should be benchmark-checked before any route surface
is built), or decide \path{Theorem49BoundaryRootSynthesis} directly on the
shell-bearing \path{wheelWithInnerTriangleEmbedding}.  The latter appears
undecided: the ledger proves the synthesis on a degenerate wheel embedding
only, and every wheel refutation concerns intermediate geometry layers rather
than the synthesis endpoint itself.  If the endpoint fails there,
\path{Theorem49ShellClaim} --- and with it the manuscript's Theorem~4.9 as
stated over its own hypothesis shells --- is refuted by an explicit
counterexample.

\section{Conclusion}

The Lean development is currently more valuable as a route-auditing tool than as
a completed proof artifact.  It has now done three distinct jobs.  First, it
preserves the historical negative result against the older per-run purification
claim and locates that claim precisely.  Second, it proves that the revised v23
local purification step is sound on its own theorem surface.  Third, it refutes
the manuscript's naive forest-leaf and adjacency-only peeling reductions while
also recording a sound replacement for them inside the live Lean route.  Fourth,
it isolates the interior-vertex side condition for boundary erasure and provides
a checked purification bridge instead of silently applying quotient endpoints to
arbitrary vertex sets.  The main remaining live blocker is now the missing
ordered boundary / annulus geometry needed to instantiate that stronger
replacement and identify the intended interior vertices.

That is useful progress.  It means future work can proceed with a cleaner map:
\begin{itemize}[leftmargin=2em]
\item do not conflate the older \(\gamma \cdot 1_R\) contradiction with the
exact v23 Step~2 claim;
\item do treat the exact v23 purification statement as locally checked;
\item do not build further on the current Theorem~4.9 / boundary-order route
without instantiating the already formalized stronger adjacency-distance and
root-distance interfaces from honest geometry;
\item do not apply boundary-erasure endpoints to an unpurified vertex set unless
the selected-boundary avoidance condition has been discharged.
\end{itemize}

In that sense the present note is not a declaration of final impossibility.  It
is a sharper ledger of what is already known to fail, what remains unverified,
and where honest continuation should focus next.

\end{document}
